\documentclass[11pt]{article}
\usepackage[utf8]{inputenc}
\usepackage[T1]{fontenc}
\usepackage{lmodern}
\usepackage[margin=1in]{geometry}
\usepackage{amsmath,amssymb,amsthm,mathtools}
\usepackage{booktabs,longtable,array,enumitem}
\usepackage[colorlinks=true,linkcolor=black,citecolor=black,urlcolor=black]{hyperref}

\usepackage{fancyhdr}
\pagestyle{fancy}
\fancyhf{}
\fancyhead[L]{\small APF Technical Supplement}
\fancyhead[R]{\small E.\,S.\ Brooke}
\fancyfoot[C]{\thepage}
\renewcommand{\headrulewidth}{0.4pt}
\usepackage{setspace}
\setstretch{1.05}
\setlength{\parskip}{0.6em}
\setlength{\parindent}{0pt}
\emergencystretch=3em

\newtheorem{theorem}{Theorem}[section]
\newtheorem{lemma}[theorem]{Lemma}
\newtheorem{proposition}[theorem]{Proposition}
\newtheorem{corollary}[theorem]{Corollary}
\newtheorem{definition}[theorem]{Definition}
\newtheorem{assumption}[theorem]{Assumption}
\newtheorem{example}[theorem]{Example}
\newtheorem{remark}[theorem]{Remark}

\newcommand{\A}{\mathfrak A}
\newcommand{\Cont}{\mathsf{Cont}}
\newcommand{\Adm}{\mathsf{Adm}}
\newcommand{\Cost}{\mathsf{Cost}}
\newcommand{\Sep}{\mathrm{Sep}}
\newcommand{\IJC}{\mathrm{IJC}}
\newcommand{\Def}{\mathrm{Def}}
\newcommand{\Rec}{\mathrm{Rec}}
\newcommand{\Bool}{\mathsf{Bool}}
\newcommand{\End}{\mathrm{End}}
\newcommand{\supp}{\mathrm{supp}}
\newcommand{\id}{\mathrm{id}}
\newcommand{\R}{\mathbb R}
\newcommand{\C}{\mathbb C}
\newcommand{\N}{\mathbb N}
\newcommand{\AQtest}{\mathcal A_{Q,\mathcal T}}
\newcommand{\eps}{\varepsilon}
\newcommand{\preceqcont}{\preceq_{\mathrm{cont}}}
\newcommand{\coderef}[2]{\smallskip\noindent\textit{Code reference:} \texttt{#1} in \texttt{#2}.}

\newenvironment{apfbox}{\par\medskip\hrule\vspace{0.5em}}{\vspace{0.5em}\hrule\par\medskip}
\newenvironment{warningbox}{\par\medskip\hrule\vspace{0.4em}\textbf{Load-bearing condition}\par\vspace{0.3em}}{\vspace{0.5em}\hrule\par\medskip}
\newenvironment{auditbox}{\par\medskip\hrule\vspace{0.4em}\textbf{Audit table}\par\vspace{0.3em}}{\vspace{0.5em}\hrule\par\medskip}
\newenvironment{proofbox}{\par\medskip\hrule\vspace{0.4em}\textbf{Proof spine}\par\vspace{0.3em}}{\vspace{0.5em}\hrule\par\medskip}

\title{Technical Supplement to Paper 1:\\
The Enforceability of Distinction\\[2pt]
\large Admissible Possibility Spaces, Continuation Algebra, and the Sep/IJC Representation Theorem}
\author{E.\,S.\ Brooke}
\date{May 2026 \quad\textbar\quad v8.31}

\begin{document}
\maketitle

\begin{abstract}
This supplement gives a self-contained formal foundation for Paper~1. We keep the primitive spine small: physical identity is finite admissible continuation identity, a physical distinction is a finite enforceable separator of continuation profiles, and physical distinctions carry positive enforcement cost. From this spine we construct the physical quotient, descent of distinctions and filters, an enforcement ledger, convex mixtures from finite classical control, and a continuation-word algebra for finite sequential interrogation. The main theorem is the Sep/IJC representation theorem: a finite joint-query interface is Sep exactly when its tested continuation-word behavior admits a faithful commutative record representation, or equivalently a finite Boolean-record defender. IJC names the failure of every such faithful commuting completion. We do not use it as a synonym for quantum behavior.

The scope is finite-interface foundations. The finite query family, finite tested word set, operational resolution, and positive resolved floor arise from finite interrogation. Compact robust-interface hypotheses extend this finite normal form under explicit continuity and lower-semicontinuity assumptions. Cost monotonicity requires no-added-resource coarse-graining; additivity and independent counting require factorized or ledger-independent support. We state these trigger conditions, give concrete examples and counterexamples, relate the static finite fragment to GPT, Spekkens, and sheaf-theoretic criteria, and separate classical Boolean feasibility from downstream quantum-admissibility questions such as NPA certification. Hilbert space and the Born rule are deferred to the quantum-structure supplement, where additional record-positivity and composite-tomography assumptions enter.
\end{abstract}

\begin{apfbox}
A physical state is a finite admissible continuation profile: it is defined by what it can still do. A physical distinction is a finite positive-cost separator of such profiles. Classical behavior occurs when the relevant distinctions fit inside one shared Boolean record. IJC occurs when no faithful commuting continuation completion exists.
\end{apfbox}

\tableofcontents
\newpage



\section*{Reader guide}
\addcontentsline{toc}{section}{Reader guide}

This supplement is organized to keep primitive commitments, forced constructions, and theorem-local hypotheses separate.

\begin{itemize}[leftmargin=2em]
\item \textbf{Primitive spine.}  Sections~\ref{sec:assumption-economy}--\ref{sec:distinctions-cost} state continuation identity, finite enforceable distinctions, and positive enforcement cost.
\item \textbf{Physical quotient and ledger.}  Sections~\ref{sec:raw-quotient-descent}--\ref{sec:distinctions-cost} construct quotient descent, enforcement cost, and the finite capacity ledger.
\item \textbf{Finite-interface normal form.}  Sections~\ref{sec:finite-interrogation-normal-form}--\ref{sec:robust-finite} show how finite interrogation induces finite query families, tested word sets, operational resolution, and positive floors.
\item \textbf{Sequential algebra.}  Section~\ref{sec:continuation-word-algebra} defines continuation words as a typed partial process algebra, including the domain of product, incompatible concatenations, order effects, contextual congruence, and the two-sided ideal quotient.
\item \textbf{Sep/IJC theorem.}  Section~\ref{sec:sep-ijc-representation} gives the self-contained representation theorem, the Boolean-feasibility form, the defender definition, and the proof architecture showing the nontrivial obstruction beyond definition-unfolding.
\item \textbf{Examples and bridges.}  Later sections give Sep and IJC examples, calibration and no-added-resource tests, finite robustness and limiting-scope statements, minimal defenders, and static bridges to GPT, Spekkens, and sheaf-theoretic formulations.
\end{itemize}

\section{Purpose and theorem status}\label{sec:purpose-status}

This supplement is self-contained. A reader does not need to consult any earlier paper in order to understand the formal definitions, hypotheses, and theorem statements used here.

The purpose is narrow. We establish the minimal continuation/enforcement foundation and the exact Sep/IJC algebraic classification. We do \emph{not} derive Hilbert space, complex amplitudes, the Born rule, quantum tensor products, or the Standard Model. Those require additional representation conditions and are deferred.

\medskip
\noindent\textit{Relation to the continuation-calculus language paper.}\quad
The primitive spine of this supplement --- physical identity as finite admissible continuation identity, distinctions as positive-cost separators of continuation profiles --- is the algebraic formalization of the continuation calculus introduced as primitive in Paper~10 (\emph{The Calculus of Finite Continuability}).  Paper~10's primitive judgment $\Gamma\vDash_C D\rightsquigarrow E$ (``in context $\Gamma$, distinction-expression $D$ can continue as $E$ using capacity at most $C$'') is the language-layer reading of the present supplement's enforceability statements.  Enforcement is the case where $E\in[D]_\Gamma$; cost is the infimum capacity over admissible continuations; the finite-interface normal form developed below in \S\ref{sec:finite-interrogation-normal-form} is the constructive specialization of the continuation calculus to finite tested protocols.  The two papers are stacked, not parallel: Paper~10 names the primitive language; the present supplement formalizes it as Admissible Possibility Spaces with continuation-equivalence classes, the continuation-word algebra of \S\ref{sec:continuation-word-algebra}, and the Sep/IJC representation theorem.  No claim in the present supplement requires Paper~10 to land first; the dependency is interpretive.  Readers seeking the most primitive available statement of what this supplement formalizes should read Paper~10 in parallel.

\medskip
\noindent\textit{PLEC's four features as derived consequences.}\quad
The four constitutive features of the Principle of Least Enforcement Cost named in the broader APF corpus --- A1 (capacity bound), MD (positive cost floor), A2 ($\arg\min$ selection), BW (cost-spectrum non-degeneracy) --- are derivable consequences of the primitive spine of this supplement under the continuation-calculus reading of Paper~10 v1.11, \S3.5.  In that reading: A1 is one half of the finite-physical-regime hypothesis ($C_\Gamma < \infty$); the value of MD ($\varepsilon^* > 0$) is the other half, with the tested/gauge cleavage (positive-cost separators versus zero-cost continuation-profile-preserving relabelings) derived from the continuation primitives; A2 is derived (Lemma~A2 of Paper~10) from cost-as-infimum and a no-waste condition under saturation; BW is derived from MD (Lemma~BW) as the cost-resolution reading of the marginal floor.  This dependency is interpretive: the present supplement does not assume the four PLEC features as separate axioms, and downstream papers (Paper~4, Paper~5, Paper~6, Paper~8) that invoke PLEC by name can read the invocation as shorthand for the same primitives unfolded through Paper~10's reductions.

\medskip
\noindent\textit{On the use of operational language.}\quad
The Admissibility Physics framework is, at its deepest level, a structural framework. The continuation-word algebra constructed in \S\ref{sec:continuation-word-algebra}, the per-interface ledger of \S\ref{sec:distinctions-cost}, and the Sep/IJC representation theorem of \S\ref{sec:sep-ijc-representation} are static algebraic objects; what reads in the prose throughout this supplement as enforcement, maintenance, and perturbation is the local reading of those structures under a temporal slicing. Time itself is derived in the broader corpus: Paper~3 derives the arrow of time as the loss of global admissible coordination, and Paper~6 derives the metric and dimensionality of spacetime. The supplement's primitive spine --- continuation identity, finite enforceable separator of continuation profiles, positive enforcement cost --- uses operational vocabulary because that is the working vocabulary of physics and because it grounds the abstract structure in the experimental world physicists actually inhabit. The apparent circularity --- describing what derives time using time-language --- reflects a limit of human descriptive vocabulary rather than a substantive contradiction in the framework. A reader who wants the structural reading can take every ``enforceable,'' ``maintain,'' and ``protocol'' as shorthand for the local reading of the corresponding admissibility-space structure. The technical apparatus developed below is already structural; only the pedagogical wrapping is temporal. This convention matches Paper~0 v6.0.5's ``A note on language'' in the Prelude \S1 and Paper~0's Descriptive Reading chapter on temporal language and the eternalist commitment. The framework is, in this respect, in the company of block-universe readings of general relativity, the Wheeler--DeWitt timeless quantum gravity programme, and the Page--Wootters mechanism for time emergence in entangled quantum systems --- structural at the foundation, with time as a derived feature of how the structure is read locally.

\subsection*{Terminology: IJC and standard language}
IJC stands for \emph{Intrinsic Joint Constraint}. We do not use it as a synonym for every use of the words contextuality, nonclassicality, incompatibility, or noncommutativity. We give it the precise finite-interface meaning
\[
\begin{aligned}
\IJC(\Gamma,Q)
\quad\Longleftrightarrow\quad
&\text{no faithful commuting-extension defender exists}\\
&\text{for the queried continuation words.}
\end{aligned}
\]
Thus Bell nonlocality, Kochen--Specker contextuality, measurement incompatibility, and controlled order effects are witnesses only when they rule out such a defender for the specified finite query family.  Conversely, a disturbance, coupling, or nonzero interface cost is not by itself IJC unless it obstructs every faithful commuting continuation completion.

\begin{auditbox}
\begin{center}
\scriptsize
\begin{tabular}{@{}p{0.28\linewidth}p{0.24\linewidth}p{0.40\linewidth}@{}}
\toprule
Item & Role in this supplement & Notes \\
\midrule
Continuation identity & Primitive spine & Physical identity is equality of finite admissible continuation profile. \\
Finite enforceable distinction & Primitive spine & Physical distinctions are positive-cost separators. \\
Raw substrate and quotient & Forced construction & Raw labels are quotiented by continuation identity. \\
Physical descent & Forced construction & Distinctions and filters must factor through the quotient. \\
Ledger bookkeeping & Forced construction & Finite enforceable distinctions consume capacity. \\
Finite tested interface & APF normal form & Finite protocol gives finite $Q$, finite $\mathcal T$, finite resolution, and positive resolved floor. \\
Compact robust floor & Regime extension under R1--R4 & Compactness/robustness/LSC extend the floor beyond explicitly finite protocols. \\
Cost monotonicity & Regime theorem & Requires no-added-resource coarse-graining. \\
Continuation-word algebra & Forced construction & Quotient of finite process words; no quantum operators assumed. \\
Sep/IJC branch split & Classification theorem & Existence/nonexistence of a faithful defender. \\
Sep iff commutative record representation & Main theorem & Full proof included. \\
Hilbert/Born & Deferred preview & Requires positive records and tomography. \\
\bottomrule
\end{tabular}
\end{center}
\end{auditbox}

We use the term \emph{defender} with care. A commuting-extension defender is a proposed classical completion of the queried interface. If such a defender exists, the interface is Sep for the queried family. If every such defender fails, the interface is IJC for that family.

\section{Assumption economy and reading map}\label{sec:assumption-economy}

The formal structure below is organized to separate the primitive spine of the framework from the regime hypotheses needed for particular theorems.  This is essential: the supplement should not be read as a list of independent axioms.  It is one continuation/enforcement idea unfolded under finite-interface conditions.

\begin{apfbox}
\textbf{Primitive spine.}  The supplement uses three primitive commitments:
\begin{enumerate}[label=(\roman*),leftmargin=2em]
\item physical identity is finite admissible continuation identity;
\item a physical distinction is a finite enforceable separator of continuation profiles;
\item physical distinctions carry positive enforcement cost.
\end{enumerate}
\end{apfbox}

\begin{warningbox}
\textbf{Irreducible primitive boundary.}  Positive finite enforcement cost is not derived later in this supplement.  It is the physical content of the Paper 1 spine: a distinction that can be freely maintained at zero cost is not a physical enforced distinction in APF.  All later floors, ledgers, and capacity bounds use this primitive together with finite interrogation or compact robust-interface hypotheses.  The supplement derives what follows from positive-cost distinguishability; it does not attempt to reduce that primitive to Hilbert space, probability, thermodynamics, or dynamics.
\end{warningbox}

\begin{remark}[Naturally stable structure and the positive-cost primitive]\label{rem:naturally-stable-cost}
The positive-cost primitive applies to \emph{interface-level enforced distinctions}: distinctions that must be separated, accessed, copied, stabilized, maintained, or read out as finite records in the queried context.  Stable symmetries, conserved labels, or background regularities that are not being actively separated at the interface are treated as part of the context $\Gamma$, not as zero-cost ledger entries.  If such a label becomes a queried record, then the finite operation that separates and maintains it enters the enforcement ledger.  Thus APF does not deny that some background structure is naturally robust; it denies that a \emph{physical enforced distinction} at a finite interface is maintained for free.
\end{remark}

\begin{remark}[Zero-cost alternatives and what would fail]\label{rem:zero-cost-alternative}
The positive-cost postulate is an explicit primitive boundary, not a hidden theorem.  One may study an idealized alternative theory in which some resolved interface distinctions are freely maintained at zero cost.  Such a theory is not ruled out by definition, but it is not the APF developed here.  In that alternative setting, the positive-floor theorem, finite capacity counting bound, and ledger monotonicity consequences used below would no longer follow without additional hypotheses excluding infinitely many freely maintained separators.  This supplement therefore proves conditional structure: if physical distinctions at finite interfaces require positive enforcement, then the quotient, ledger, floor, and Sep/IJC classification follow under the stated finite or compactness hypotheses.
\end{remark}


\begin{auditbox}
\textbf{Residual input inventory after compression.}  The only primitive inputs are continuation identity, finite enforceable distinction, and positive enforcement cost.  A physically executed finite interrogation supplies its own finite query family, finite tested word set, finite word-depth bound, and operational resolution.  After quotienting unresolved, unwriteable, unreadable, or unstable labels, such an interrogation induces the finite tested interface used in the Sep/IJC theorem.  The conditions that remain later in the supplement are theorem-local trigger conditions, not new primitives: no-added-resource coarse-graining is the trigger for monotonicity; factorized support is the trigger for additivity; ledger-independent support is the trigger for independent counting; and compact robust-interface hypotheses only extend the finite-protocol floor theorem to nonexplicit compact families.  When any trigger condition is absent, the corresponding conclusion is simply not invoked.
\end{auditbox}

Everything else is either a construction forced by this spine, a normal form of finite interrogation, or a theorem under explicit local hypotheses.

The quotient state space is forced because raw descriptions with the same finite admissible continuations are physically indistinguishable.  Physical distinctions and filters must descend to this quotient; otherwise they separate raw labels rather than physical possibilities.  The ledger is forced because finite enforceable distinctions consume capacity.  Convexity is not assumed as background linear structure; it appears when admissible finite classical control records are available.

The finite query family, finite tested word set, finite depth, finite operational resolution, and positive resolved floor are induced by a finite APF interrogation protocol.  The extra hypotheses that appear later are local theorem conditions.  Robust finite interfaces R1--R4 are compactness, robustness, lower-semicontinuity, and finite-capacity assumptions extending the positive-floor result beyond explicitly finite protocols.  No-added-resource coarse-graining is the condition under which cost monotonicity is true: a filter that only forgets, identifies, restricts, or lowers resolution cannot make a distinction more expensive unless fresh enforcement resources are introduced.  Factorized support is the condition under which ledger costs add.  Ledger-independent operational resolution is the condition under which the finite-resolution theorem is allowed to count word classes as independent capacity burdens.

The continuation-word algebra is also not an independent physical assumption.  It is the algebraic bookkeeping forced by finite sequential query behavior after quotienting by context-complete continuation equivalence.  Its linear span is a formal representation device; physical probabilistic mixtures enter separately through admissible classical control.

With these distinctions in place, the Sep/IJC theorem has a narrow role.  It does not derive Hilbert space or the Born rule.  It proves the exact branch criterion for a finite joint-query interface:
\[
\begin{aligned}
(\Gamma,Q,\mathcal T)\in\Sep
\quad\Longleftrightarrow\quad
&\text{the finite continuation-word behavior admits}\\
&\text{a faithful commutative record representation.}
\end{aligned}
\]
Thus IJC is precisely the failure of every faithful Boolean-record completion.

\begin{center}
\scriptsize
\begin{tabular}{@{}p{0.24\linewidth}p{0.65\linewidth}@{}}
\toprule
Layer & Contents \\
\midrule
Primitive spine & Continuation identity; finite enforceable distinctions; positive enforcement cost. \\
Forced constructions & Quotient state space; descent; ledger bookkeeping; convex mixtures from classical control; continuation-word quotient algebra. \\
Finite-protocol normal form & Finite $Q$, finite $\mathcal T$, finite depth, finite operational resolution, positive resolved floor. \\
Local theorem hypotheses & Compact robust extension; no-added-resource filters; factorized support; ledger independence. \\
Classification theorem & Sep iff faithful commutative record representation; IJC iff no such representation exists. \\
Deferred representation theory & Hilbert space, complex amplitudes, tensor products, and Born probabilities. \\
\bottomrule
\end{tabular}
\end{center}

\begin{auditbox}
\textbf{Trigger-condition failure modes.}  The local hypotheses used later are included to prevent overclaiming.  If the trigger is absent, the theorem that depends on it is disabled rather than repaired by assumption.
\begin{center}
\scriptsize
\begin{tabular}{@{}p{0.24\linewidth}p{0.32\linewidth}p{0.32\linewidth}@{}}
\toprule
Trigger condition & Conclusion licensed & What fails without it \\
\midrule
No-added-resource filtering & Cost monotonicity & A filter may secretly add ancillas, stabilization, amplification, or fresh records, so cost can increase or be displaced. \\
Factorized enforcement support & Additivity of ledger costs & Shared infrastructure or coupled implementation may make total cost subadditive or nonseparable. \\
Ledger-independent resolution & Independent counting bound & Multiple formal distinctions may be carried by one coarser record and must not be counted separately. \\
Compact robust family & Positive floor beyond explicit finite tables & Noncompact or unstable families may have costs accumulating toward zero without a minimum. \\
Finite tested word set & Finite tested algebra & Infinite word families require additional analytic completion not supplied by Paper 1. \\
\bottomrule
\end{tabular}
\end{center}
\end{auditbox}

A compact way to read the supplement is therefore:
\begin{quote}
If physical identity is what can still be done, then distinctions are finite-cost separations of continuation profiles; finite sequential queries therefore generate a quotient word algebra, and an interface is classical exactly when that algebra admits a faithful commutative record representation.
\end{quote}

\section{Admissible possibility spaces}\label{sec:admissible-possibility-spaces}

\begin{definition}[Raw possibility substrate]
A \emph{raw possibility substrate} in context $\Gamma$ is a set $X_\Gamma$ of candidate configurations, descriptions, histories, preparations, boundary states, or other raw representatives before quotienting by physical equivalence.
\end{definition}

\begin{definition}[Finite admissible continuation]
For $x\in X_\Gamma$, $\Cont_\Gamma(x)$ denotes the class of finite admissible continuations available from $x$.  A continuation may be a finite test, interaction, perturbation, composition, record-writing operation, maintenance operation, or finite process fragment that remains admissible in context $\Gamma$.
\end{definition}

A continuation is not a metaphysical future.  It is a finite operational possibility: something the system can still physically do, undergo, record, or be used for without leaving the admissible context.

\begin{definition}[Continuation identity]
For raw possibilities $x,y\in X_\Gamma$, define
\[
 x\sim_\Gamma y
 \quad\Longleftrightarrow\quad
 \Cont_\Gamma(x)=\Cont_\Gamma(y).
\]
The physical possibility space is the quotient
\[
 \Omega_\Gamma := X_\Gamma/\sim_\Gamma .
\]
The equivalence class of $x$ is written $[x]_\Gamma$.
\end{definition}

\begin{definition}[Admissible possibility space]
An \emph{admissible possibility space} is a tuple
\[
\A_\Gamma=
(X_\Gamma,\Cont_\Gamma,\sim_\Gamma,\Omega_\Gamma,
\mathcal D_\Gamma,\kappa_\Gamma,C_\Gamma,\circ),
\]
where $X_\Gamma$ is a raw substrate, $\Cont_\Gamma$ is finite admissible continuation behavior, $\sim_\Gamma$ is continuation identity, $\Omega_\Gamma$ is the physical quotient, $\mathcal D_\Gamma$ is the family of enforceable distinctions, $\kappa_\Gamma$ is an enforcement-cost valuation, $C_\Gamma<\infty$ is the finite available capacity, and $\circ$ denotes admissible composition where defined.
\end{definition}

\begin{remark}[Not a hidden Hilbert space]
No linear structure, probability measure, inner product, tensor product, or Hilbert space is assumed in the definition of $\A_\Gamma$.
\end{remark}

\begin{definition}[Continuation preorder]
For physical possibilities $[x],[y]\in\Omega_\Gamma$, define
\[
 [x]\preceqcont [y]
 \quad\Longleftrightarrow\quad
 \Cont_\Gamma(y)\subseteq \Cont_\Gamma(x).
\]
Thus $[y]$ is downstream of $[x]$ when it has no more admissible continuations than $[x]$.
\end{definition}

\begin{lemma}[Continuation preorder]
The relation $\preceqcont$ is a preorder on $\Omega_\Gamma$.
\end{lemma}
\begin{proof}
Reflexivity follows from $\Cont(x)\subseteq\Cont(x)$.  If $[x]\preceqcont[y]$ and $[y]\preceqcont[z]$, then $\Cont(z)\subseteq\Cont(y)\subseteq\Cont(x)$, hence $[x]\preceqcont[z]$.
\end{proof}

\begin{definition}[Irreversible continuation]
An admissible transition $[x]\to[y]$ is \emph{continuation-contracting} when
\[
\Cont_\Gamma(y)\subsetneq\Cont_\Gamma(x).
\]
It is \emph{irreversible in context $\Gamma$} when the lost continuation classes cannot be restored by any finite admissible continuation without additional enforcement cost.
\end{definition}

\section{Raw substrate, quotient, and physical descent}\label{sec:raw-quotient-descent}

A common source of confusion is whether definitions are made on raw descriptions or on physical possibilities after quotienting.  This section makes the distinction explicit.

\begin{definition}[Quotient map]
Let
\[
 q_\Gamma:X_\Gamma\to\Omega_\Gamma,
 \qquad q_\Gamma(x)=[x]_\Gamma
\]
be the canonical quotient map.
\end{definition}

\begin{definition}[Physical descent of a distinction]
A raw binary predicate $d:X_\Gamma\to\{0,1\}$ \emph{descends to the quotient} if there exists a map
\[
 \bar d:\Omega_\Gamma\to\{0,1\}
\]
such that
\[
 d=\bar d\circ q_\Gamma .
\]
Equivalently, $x\sim_\Gamma y$ implies $d(x)=d(y)$.
\end{definition}

\begin{lemma}[Only quotient-descending predicates can be physical distinctions]
If $d$ separates two raw representatives $x,y$ with $x\sim_\Gamma y$, then $d$ is not a physical distinction in context $\Gamma$.
\end{lemma}
\begin{proof}
If $x\sim_\Gamma y$, then $x$ and $y$ have identical finite admissible continuation behavior.  A predicate that differs on them separates raw notation rather than physical possibility.  Hence it cannot represent a finite admissible continuation-separating distinction.  Equivalently, it fails to descend through $q_\Gamma$.
\end{proof}

\begin{definition}[Physical descent of filters]
A raw map $F:X_\Gamma\to X_{\Gamma'}$ descends to a physical map when
\[
 x\sim_\Gamma y\quad\Longrightarrow\quad F(x)\sim_{\Gamma'}F(y).
\]
In that case there is an induced map
\[
 \bar F:\Omega_\Gamma\to\Omega_{\Gamma'},
 \qquad \bar F([x])=[F(x)].
\]
\end{definition}

\begin{proposition}[Well-defined induced filter]
If $F$ respects continuation identity, then $\bar F([x])=[F(x)]$ is well-defined.
\end{proposition}
\begin{proof}
If $[x]=[y]$, then $x\sim_\Gamma y$.  By hypothesis $F(x)\sim_{\Gamma'}F(y)$, so $[F(x)]=[F(y)]$.  Thus the induced value does not depend on the chosen raw representative.
\end{proof}

\begin{remark}[Categorical reading]
Admissible possibility spaces may be treated as objects and quotient-descending admissibility-preserving filters as morphisms.  This supplement uses only the elementary set-level formulation.
\end{remark}

\section{Distinctions and enforcement cost}\label{sec:distinctions-cost}

This section contains the second and third elements of the primitive spine: physical distinctions are finite enforceable separators, and such distinctions carry positive cost.  The ledger language introduced here is bookkeeping forced by those commitments, not an additional ontology.

\begin{definition}[Physical distinction]
A \emph{physical distinction} in context $\Gamma$ is a quotient-descending separator
\[
 d:\Omega_\Gamma\to\{0,1\}
\]
for which there exists a finite admissible enforcement protocol maintaining the separated alternatives against admissible perturbations.
\end{definition}

\begin{definition}[Enforcement cost]
The enforcement cost of a distinction $d\in\mathcal D_\Gamma$ is
\[
 \kappa_\Gamma(d)\in(0,\infty),
\]
the finite capacity burden required to stabilize, record, or maintain the distinction in context $\Gamma$.
\end{definition}

\begin{definition}[Ledger admissibility]
For a finite family $S\subset\mathcal D_\Gamma$, define its ledger cost by
\[
 \kappa_\Gamma(S):=\sum_{d\in S}\kappa_\Gamma(d)
\]
when costs are independently supported.  More generally, write
\[
 \kappa_\Gamma(S)=\sum_{d\in S}\kappa_\Gamma(d)+\kappa_{\Gamma,\mathrm{int}}(S)
\]
where $\kappa_{\Gamma,\mathrm{int}}$ is an interface or coupling term.  The family $S$ is ledger-admissible when
\[
 \kappa_\Gamma(S)\le C_\Gamma .
\]
\end{definition}

\begin{remark}[Ledger valuation convention]\label{rem:ledger-valuation-convention}
For a finite family $S$, $\kappa_\Gamma(S)$ is the cost of the least admissible joint enforcement protocol for the family at the stated resolution.  When supports are independent this valuation reduces to a sum.  When supports share implementation, the interface term records the difference between the independent-sum estimate and the least joint protocol.  No theorem below assumes additivity except where factorized support or ledger independence is explicitly stated.
\end{remark}

\begin{remark}[On the structural status of $\kappa_{\Gamma,\mathrm{int}}$]\label{rem:kappa-int-placeholder}
The interface-cost term $\kappa_{\Gamma,\mathrm{int}}(S)$ is defined by residue: it is the difference between the independent-sum estimate and the least joint protocol cost.  Two structural questions about its shape come apart.

The \emph{lower-bound side} --- how negative can $\kappa_{\Gamma,\mathrm{int}}$ become through shared substrate? --- is settled in \S\ref{sec:composition-factorization} by the marginal-floor structural inequalities (Lemma~\ref{lem:marginal-floor-on-joint-cost}, Corollary~\ref{cor:sum-of-floors-lower-bound}, Theorem~\ref{thm:kappa-int-singleton-shape}, Corollary~\ref{cor:kappa-int-binary-shape}): the joint cost cannot drop below the sum of marginal floors $n\varepsilon^*_\Gamma$, hence $\kappa_{\Gamma,\mathrm{int}}(S) \ge n\varepsilon^*_\Gamma - \sum_{d\in S}\kappa_\Gamma(d)$.  Coarse-graining monotonicity (Lemma~\ref{lem:coarse-graining-monotonicity}) and monotonicity in the family (Lemma~\ref{lem:joint-cost-monotonicity-in-S}) further constrain the joint cost's structural behavior.  Each of these inequalities follows directly from the spine: the marginal-floor definition (Remark~\ref{rem:marginal-floor}) plus its positivity from MD/BW (Paper~10 v1.12 \S3.5, Lemma~BW), plus the infimum-over-admissible-protocols valuation convention (Remark~\ref{rem:ledger-valuation-convention}).  No new axiom is added; the structural shape is exposed rather than imposed.

The \emph{upper-bound side} --- how positive can $\kappa_{\Gamma,\mathrm{int}}$ become through coupling overhead, given the shared-substrate dynamics of the context? --- is settled in \S\ref{subsec:kappa-int-upper-bound} for the continuum-bridge regime characterized by C1--C5.  Theorem~\ref{thm:kappa-int-binary-upper-bound} bounds the binary-form interface term by the cooperative-cost kernel sup against the profile $L^1$ norms; Corollary~\ref{cor:kappa-int-far-separation} gives exponential suppression for far-separated supports; Corollary~\ref{cor:kappa-int-singleton-upper-bound} extends the bound to the singleton form; and Theorem~\ref{thm:kappa-int-two-sided-bound} packages the lower and upper bounds together as the explicit two-sided structural-rigidity statement.  This bound is internal to the present supplement; the apparatus is the continuum bridge of \S\ref{sec:recruitment-radius} extended to the multi-distinction setting via \eqref{eq:Erec-multi}--\eqref{eq:Ecross}.

Outside the continuum-bridge regime --- saturation regimes where C4 fails (second-order expansion breaks down), far-from-equilibrium regimes where C5 fails (linear-response approximation fails; this is the cosmogenesis regime of Paper~25 \S8), substrate-collapse regimes where C3 fails (locality fails) --- the upper bound does not transfer automatically.  The relevant substrate-mode work for these non-bridge regimes lives in Paper~25 \S\S\,2--3 (recruitment-radius foundations) and the Paper~6 supplement \S1.5 (geometric H1/H2 closure); downstream papers reasoning about coupling overhead in those regimes should supply the relevant upper bound from substrate-level commitments.

Where ledger additivity is needed (e.g., the independent-counting theorem, the resolved-floor theorem), the relevant trigger condition is factorized support or ledger-independent resolution, both of which set $\kappa_{\Gamma,\mathrm{int}} = 0$ in the regimes where they apply (Theorem~\ref{thm:ledger-additivity-protocol-factorization}).  The strengthened counting bound (Corollary~\ref{cor:counting-bound-from-marginal-floor}) shows that the counting conclusion in fact survives without the additivity hypothesis: it follows directly from the marginal-floor structural inequality applied to any admissible family, with or without support-independence.  Where the marginal-floor formulation (Remark~\ref{rem:marginal-floor}) is the operative invariant --- shared-substrate / broken-symmetry / confined-phase regimes --- the marginal floor handles the shared-commitment piece and $\kappa_{\Gamma,\mathrm{int}}$ records only the residual cooperative interface bookkeeping.
\end{remark}

\begin{remark}[Operational meanings of cost]
The symbol $\kappa_\Gamma$ is abstract by design.  In concrete models it may be bounded by minimum work to write a stable record, energy-time control cost, statistical sample complexity at fixed error, channel resources required to maintain distinguishability, or thermodynamic cost of record support.  We require only positivity, finiteness in the context, and monotonicity under admissible coarse-graining where stated.
\end{remark}

\begin{remark}[Marginal floor and shared substrate]\label{rem:marginal-floor}
The cost $\kappa_\Gamma(d)$ as defined here is an in-isolation cost: the finite capacity burden required to maintain $d$ in $\Gamma$, computed without reference to other distinctions that may share substrate support.  In regimes where many distinctions share a structural commitment --- a broken vacuum, a confined-phase substrate, a topological infrastructure --- the in-isolation cost overcounts: the shared piece is paid once by the carrying context, not once per distinction.  The right invariant for capacity counting in such regimes is the \emph{marginal floor}
\[
        \varepsilon^*_\Gamma \;:=\; \inf\Big\{\kappa_\Gamma(S \otimes d) - \kappa_\Gamma(S) : S \text{ admissible in } \Gamma,\ d \text{ tested in } \Gamma,\ S \otimes d \text{ admissible}\Big\},
\]
which measures only the genuinely-new burden of adding a tested distinction to any admissible context (Paper~10, \emph{The Calculus of Finite Continuability}, v1.11, \S3, ``Why the marginal floor'').  When the substrate has no shared structural commitments, the marginal floor and the in-isolation cost floor coincide, and the present supplement's positive-cost theorems (the resolved-floor theorem, the independent-counting bound, the finite capacity ledger) carry over verbatim.  In broken-symmetry or confined-phase regimes --- including the Standard Model interface, where Paper~4's $C_{\mathrm{total}}=61$ count reads as a count of marginal channel-distinctions on top of the broken-vacuum and confined-phase substrate --- the marginal floor is the operative invariant.  The interface-cost term $\kappa_{\Gamma,\mathrm{int}}(S)$ in the ledger admissibility definition above already records the difference between the independent-sum estimate and the least joint protocol cost; the marginal-floor convention is the systematic specialization of that interface bookkeeping to substrates with shared commitments.
\end{remark}

\subsection*{Concrete calibration models for $\kappa_\Gamma$}\label{subsec:kappa-calibration}

The abstract ledger can be instantiated in more than one physically legitimate unit system.  We do not require a universal microscopic formula for $\kappa_\Gamma$. We require only that, once a context, resolution, and readout protocol are fixed, the cost assigned to a resolved distinction is positive and finite.  Two common calibrations are useful for experimental readers.

\begin{example}[Sample-complexity calibration]\label{ex:kappa-sample-complexity}
Suppose a binary distinction $d$ is operationally resolved by estimating a Bernoulli probability to separate two continuation profiles whose predicted probabilities differ by at least $\Delta>0$.  If one demands error probability at most $\alpha$, Hoeffding's inequality gives the sufficient condition
\[
  N \geq \frac{2}{\Delta^2}\log\frac{2}{\alpha}
\]
independent trials, up to constant conventions depending on whether one estimates one mean or a difference of two means.  If each trial consumes capacity cost $c_{\rm trial}>0$ and writing the final stable bit consumes $c_{\rm rec}>0$, one may use the operational lower-bound model
\[
  \kappa_\Gamma(d;\Delta,\alpha)
  \geq
  c_{\rm rec}+c_{\rm trial}\frac{2}{\Delta^2}\log\frac{2}{\alpha}.
\]
The important APF point is not the particular concentration inequality.  It is that finite resolution and finite confidence turn ``distinguishability'' into a positive operational burden.
\end{example}

\begin{example}[Stable-record work calibration]\label{ex:kappa-landauer-calibration}
If a resolved distinction must be written into an irreversible stable bit in a bath at temperature $T$, then a conservative thermodynamic calibration assigns a work floor on the order of
\[
  W_{\rm bit}\geq k_{\rm B}T\log 2
\]
per erased or irreversibly reset bit.  If lifetime stability $\tau$ is enforced by an activation barrier with attempt time $\tau_0$, a common engineering estimate adds a barrier scale
\[
  E_{\rm barrier}\gtrsim k_{\rm B}T\log(\tau/\tau_0).
\]
Thus a stable finite record can be charged in work/energy units rather than sample units.  These are calibration choices for concrete models, not extra APF axioms.
\end{example}


\begin{example}[Combined sampling-plus-record ledger]\label{ex:combined-kappa-ledger}
In a realistic detector, resolving a binary distinction usually requires both repeated sampling and stable record support.  A conservative composite ledger can use a chosen calibration map
\[
  \chi_\Gamma:\{\text{trials, work, memory cells, control actions}\}\to\R_{\ge0}
\]
into one context-normalized capacity unit, and then charge
\[
  \kappa_\Gamma(d)
  \ge
  \chi_\Gamma(N_{\Delta,\alpha}\,\text{trials})
  +\chi_\Gamma(W_{\rm bit}\,\text{work})
  +\chi_\Gamma(M_{\rm stab}\,\text{maintained memory}) .
\]
The calibration map is part of the experimental context; it is not universal physics.  Once chosen, however, all comparisons inside that context must use the same map.  This is enough for the finite-interface theorems, which depend only on positive costs, finite capacity, and the explicitly stated monotonicity/additivity triggers.
\end{example}


\begin{example}[Numerical ledger/floor case study]\label{ex:numerical-kappa-case-study}
\textbf{Two-outcome noisy detector.}
Suppose a detector must distinguish two binary continuation profiles whose outcome probabilities differ by at least \(\Delta=0.10\), with error probability at most \(\alpha=0.01\).  The Hoeffding calibration gives
\[
  N_{\Delta,\alpha}
  =\left\lceil \frac{2}{0.10^2}\log \frac{2}{0.01}\right\rceil
  =\left\lceil 200\log 200\right\rceil
  =1060
\]
trials.  If each trial is charged one normalized unit and the stable final record is charged ten normalized units, then
\[
  \kappa_\Gamma(d)\ge 1060+10=1070.
\]
If the experiment has total capacity \(C_\Gamma=5000\) normalized units and all resolved distinctions use the same calibration, the positive floor may be taken as \(\mu_\Gamma=1070\), giving the independent-counting bound
\[
  N_{\rm indep}\le \left\lfloor \frac{5000}{1070}\right\rfloor=4 .
\]
The same detector can also be reported in work units.  At room temperature, \(k_{\rm B}T\log2\) is about \(2.9\times10^{-21}\,\mathrm J\) per irreversible bit erasure.  Requiring one-day stability with attempt time \(\tau_0=10^{-9}\,\mathrm s\) gives an activation-barrier scale
\[
  k_{\rm B}T\log(86400/10^{-9})\approx 1.3\times10^{-19}\,\mathrm J.
\]
Those two reports are not competing laws; they are two calibrations of the same APF requirement: once the context fixes the ledger unit, resolved distinctions have positive finite cost and the floor/counting theorem becomes numerically checkable.
\end{example}

\begin{example}[Calibration sensitivity of the independent-counting bound]\label{ex:kappa-sensitivity}
Continuing Example~\ref{ex:numerical-kappa-case-study}, keep \(\Delta=0.10\), \(\alpha=0.01\), so that \(N_{\Delta,\alpha}=1060\), and keep total normalized capacity \(C_\Gamma=5000\).  Vary only the calibration map \(\chi_\Gamma\).  The theorem is unchanged across the rows; what changes is the context-indexed ledger cost assigned to the same resolved distinction.
\begin{center}
\small
\setlength{\tabcolsep}{3pt}
\renewcommand{\arraystretch}{1.15}
\begin{tabular}{@{}p{0.25\linewidth}p{0.12\linewidth}p{0.20\linewidth}p{0.16\linewidth}p{0.14\linewidth}@{}}
\toprule
case & trial cost & record/stability cost & $\kappa_\Gamma(d)$ & $\lfloor C_\Gamma/\kappa_\Gamma(d)\rfloor$\\
\midrule
baseline & 1 & 10 & 1070 & 4\\
expensive trials & 2 & 10 & 2130 & 2\\
memory constrained & 1 & 500 & 1560 & 3\\
long lifetime target & 1 & 1000 & 2060 & 2\\
cheap sampling, costly record & 0.5 & 1500 & 2030 & 2\\
\bottomrule
\end{tabular}
\end{center}
Thus unequal per-trial costs, memory constraints, and lifetime targets do not alter the representation theorem; they alter the calibrated interface capacity available for independent resolved distinctions.  This is why APF treats \(\kappa_\Gamma\), \(\mu_\Gamma\), and \(C_\Gamma\) as context-indexed quantities rather than universal constants.
\end{example}

\begin{remark}[Platform comparisons and scope of empirical calibration]\label{rem:platform-comparison-scope}
The calibrations above are not claimed as empirical validation across physical platforms.  APF supplies a ledger template: once an experimental context specifies the cost unit, resolution, confidence level, record lifetime, and available capacity, the same finite-interface accounting computes \(\kappa_\Gamma\), \(\mu_\Gamma\), and the corresponding capacity bound.

A cross-platform comparison therefore requires platform-specific engineering data.  For example, a photonic implementation might charge trial cost, detector dark counts, timing-bin stability, optical loss correction, and coincidence-window synchronization.  An ion-trap implementation might charge state-preparation time, fluorescence readout error, memory lifetime, heating, sympathetic cooling, and control-pulse stability.  APF does not assert that these ledgers have the same absolute units. We assert only that, after a context-specific calibration map \(\chi_\Gamma\) is fixed, resolved distinctions carry positive finite cost, and no additivity, monotonicity, or independent-counting claim is made without the relevant trigger condition.

Thus platform comparison is a downstream empirical application of the ledger formalism, not an additional assumption of the Sep/IJC theorem.  The present supplement provides the accounting rules and finite toy cases; full photonic-versus-ion-trap ledgers require measured instrument specifications and are outside the finite representation theorem proved here.
\end{remark}

\begin{remark}[Resolution scaling of the floor]
In the sample-complexity calibration, decreasing the resolution gap \(\Delta\) generally raises the required trial count like \(\Delta^{-2}\) at fixed confidence.  Thus finer resolution does not threaten the finite-interface theorem; it changes the context and normally increases the cost of maintaining a resolved distinction.  In singular ideal limits where the required resolution tends to zero without a compensating finite record/stability floor, \(\mu_\Gamma\) may collapse to zero.  That is outside the finite-interface theorem and must be handled as a separate limiting construction.
\end{remark}
\begin{remark}[Units of the ledger]
The ledger may be expressed in trials, work, memory cells, control actions, or normalized capacity units.  Comparing costs across contexts requires a chosen calibration.  The representation theorem itself is invariant under monotone rescaling of $\kappa_\Gamma$ within a fixed context because it uses only positivity, finiteness, and explicitly stated monotonicity/additivity triggers.
\end{remark}

\paragraph{Calibration equivalence and empirical content.}
Two ledger calibrations $\kappa_\Gamma$ and $\kappa'_\Gamma$ are
representation-equivalent on a fixed finite tested interface when they
preserve positivity, finiteness, the no-added-resource monotonicity
order, and the admissible/inadmissible status of the finite distinction
families actually used in the theorem.  Under such a calibration change,
the Sep/IJC representation verdict is unchanged: the same tested
continuation quotient either does or does not admit a faithful Boolean
defender.

The empirical content of a calibration is therefore not the arbitrary
choice of unit.  It lies in the experimentally checkable claims that
resolved distinctions require positive finite resources, that the
declared capacity bound is finite, and that the
monotonicity/additivity/independence trigger conditions invoked in a
given application are actually satisfied.  Different platforms may
realize different numerical ledgers while remaining equivalent for the
finite representation theorem.

\paragraph{Empirical status of the ledger.}
The ledger formalism is an empirical template, not a claim that all
platforms share a universal microscopic cost unit.  Its testable
content is local and context-indexed.  Once an experiment declares a
resolution scale, confidence level, record lifetime, available control
resources, and calibration map
\[
    \chi_\Gamma:\mathcal R_\Gamma\to\R_{\ge 0},
\]
the APF ledger makes checkable claims about which distinctions can be
jointly maintained, which filters are resource-nonincreasing, and
which families may be counted independently.

Changing units or platform-specific engineering details does not
change the Sep/IJC representation theorem, provided the calibration
preserves positivity, finiteness, the no-added-resource preorder, and
the admissible/inadmissible status of the finite distinction families
used in the theorem.  What can be experimentally challenged is not
the arbitrary choice of ledger unit, but the asserted satisfaction of
the trigger conditions: positive resolved cost, finite available
capacity, resource-nonincreasing filtering, factorized support, and
support-independent resolution.

\paragraph{Boundary of the positive-cost primitive.}
The positive-cost primitive applies only to \emph{resolved enforced
distinctions} at a finite interface.  It does not say that every
formal label, background symmetry, reversible microscopic coordinate,
or unqueried degree of freedom carries an independent ledger charge.
A distinction enters the APF ledger only when it is required to be
separated, stabilized, accessed, re-read, copied, maintained, or used
as a finite record in the queried context.  The primitive thus has
the following boundary conditions.
\begin{enumerate}[label=(\roman*),leftmargin=2em]
\item \textbf{Resolvedness.}  The alternatives must remain distinct
at the declared operational resolution $\delta$.  Differences below
resolution are quotiented away and are not charged as physical
distinctions.
\item \textbf{Recordability.}  The distinction must be writeable,
readable, or otherwise operationally usable in the finite protocol.
A formal label that cannot be accessed by any tested continuation is
part of raw syntax, not a physical record.
\item \textbf{Stability.}  The distinction must persist against the
admissible perturbations of the context for the declared lifetime or
protocol depth.  Transient microscopic differences that cannot
support the tested continuation table are not counted as enforced
distinctions.
\item \textbf{Context-indexed cost.}  The cost is not a universal
scalar attached to an abstract bit.  It is the least
context-calibrated resource burden required to enforce the
distinction at the stated resolution, confidence, and stability
target.
\item \textbf{No free hidden support.}  Error correction,
amplification, ancillas, stabilizers, reference frames, clocks,
recovery memories, or control registers may lower the apparent error
rate of a record, but they are part of the enforcement protocol and
must be included in the ledger unless they are already declared as
background context.
\end{enumerate}

\paragraph{Ideal reversible and zero-cost limits.}
Ideal reversible evolution is not excluded by APF.  A reversible
process may transport or transform an already supported distinction
without adding a new ledger entry.  What APF denies is that a
\emph{new resolved enforced record} at a finite interface can be
created, stabilized, and made available for tested continuations at
zero cost.

If one takes an ideal limit in which the resolution threshold tends
to zero, the record lifetime tends to infinity, or the number of
resolved alternatives tends to infinity, then the finite-interface
theorem need not apply.  Such a limit requires its own limiting
construction.  In particular, if the resolved cost floor collapses,
$\mu_{\Gamma,\delta}\to 0$, then the independent-counting and
finite-capacity conclusions are no longer licensed.

A theory in which arbitrarily many resolved, stable, re-readable
distinctions can be maintained at zero marginal cost is a legitimate
mathematical alternative, but it is not the APF regime studied here.
In that alternative, the positive-floor theorem, capacity-counting
bounds, and ledger monotonicity consequences would fail unless
replaced by additional hypotheses.

\paragraph{Error-corrected records.}
Error correction does not evade the positive-cost primitive.  It
changes the enforcement protocol.  A logical record protected by
redundancy, syndrome extraction, active feedback, dissipation, or
recovery control is charged by the least admissible joint protocol
that maintains the logical distinction at the declared resolution and
lifetime:
\[
    \kappa_\Gamma(d_{\mathrm{logical}})
    \;=\;
    \inf_{P\in\mathsf{Prot}_\Gamma(d_{\mathrm{logical}})}
    \chi_\Gamma\bigl(r_\Gamma(P)\bigr).
\]
The physical bits, stabilizers, syndrome records, recovery
operations, clocks, and control references used by $P$ are included
in the resource load $r_\Gamma(P)$, unless they are explicitly part
of the background context $\Gamma$.

Thus an error-corrected record may reduce logical failure
probability, but it does not constitute a zero-cost distinction.  It
is a different implementation of the same resolved logical
separator, with a different context-indexed ledger cost.

\paragraph{Empirical content of positive cost.}
The positive-cost primitive is not tested by the choice of ledger
units.  It is tested by whether a declared finite interface can
support resolved, stable, re-readable distinctions without consuming
any finite resource under the chosen calibration.  A counterexample
would be a finite experimental context in which arbitrarily many
independent resolved records can be created and maintained with
zero marginal resource burden, while still satisfying the stated
resolution, confidence, and stability requirements.  Such a context
would violate the APF positive-floor premise.

Conversely, ordinary changes of calibration do not alter the finite
representation theorem when they preserve positivity, finiteness,
no-added-resource monotonicity, and the admissible/inadmissible
status of the finite distinction families used in the theorem.

\begin{definition}[Distinction refinement]
For $d,e\in\mathcal D_\Gamma$, write $d\preceq e$ when every pair of physical possibilities separated by $d$ is also separated by $e$.  Thus $e$ is at least as fine a distinction as $d$.
\end{definition}

\begin{lemma}[Physical distinctions are continuation separators]
If $d$ is a physical distinction separating $[x]$ and $[y]$, then $\Cont_\Gamma(x)\ne\Cont_\Gamma(y)$.
\end{lemma}
\begin{proof}
If $\Cont_\Gamma(x)=\Cont_\Gamma(y)$, then $[x]=[y]$ in $\Omega_\Gamma$.  A quotient-descending distinction cannot separate a class from itself.
\end{proof}

\section{Finite interrogation normal form}\label{sec:finite-interrogation-normal-form}

The finite tested objects used later are not independent assumptions.  They are the normal form of asking finitely many questions of a finite interface at finite resolution.

\begin{definition}[Finite APF interrogation protocol]\label{def:finite-apf-interrogation}
A finite APF interrogation protocol in context $\Gamma$ consists of:
\begin{enumerate}[label=(\roman*),leftmargin=2em]
\item a finite family $Q$ of admissible query operations or discriminators;
\item a finite word-depth bound $L$;
\item a finite readout alphabet for each tested discriminator;
\item an operational resolution scale identifying continuation profiles that no tested discriminator separates at that scale;
\item a finite interface ledger with available capacity $C_\Gamma<\infty$.
\end{enumerate}
The associated tested word set $\mathcal T\subseteq Q^{\le L}$ contains exactly the finite words whose continuation behavior the protocol interrogates.
\end{definition}

\begin{definition}[Induced resolved interface]\label{def:induced-resolved-interface}
Given a finite APF interrogation protocol, its induced resolved interface is obtained by quotienting raw labels, preparations, effects, records, and finite words by the equivalence relation of agreement on all tested admissible continuations at the stated resolution.  Labels that are not writeable, readable, re-readable, or stable under the admissible perturbations of $\Gamma$ are not stable finite records at this interface and are removed from the resolved record family.
\end{definition}

\begin{theorem}[Finite tested regime as APF normal form]\label{thm:finite-tested-normal-form}
Every finite APF interrogation protocol induces a finite tested continuation quotient $\mathcal A_{Q,\mathcal T}$ and a finite resolved distinction family.  If the resolved family is nonempty, then it has a positive distinction floor
\[
\mu_{\Gamma,\mathcal T}:=\min\{\kappa_\Gamma(d): d \text{ is a resolved tested distinction}\}>0.
\]
Thus finite query family, finite tested word set, finite operational resolution, and positive finite-tested distinction floor are not additional postulates; they are consequences of a finite act of APF interrogation plus the positive-cost spine.
\end{theorem}
\begin{proof}
The protocol supplies only finitely many primitive query operations and a finite word-depth bound, so $\mathcal T\subseteq Q^{\le L}$ is finite.  Finite readout resolution quotients all raw labels and word behaviors that cannot be separated by the tested discriminators at that resolution.  After this quotient, only finitely many resolved tested distinction classes remain.  By the positive-cost spine, each physical resolved distinction has strictly positive enforcement cost.  A finite nonempty set of strictly positive numbers has a positive minimum, giving $\mu_{\Gamma,\mathcal T}>0$.  Unwriteable, unreadable, non-re-readable, or perturbatively unstable labels do not define stable finite records at the stated resolution and are removed by the resolved-interface normal form, rather than counted as additional physical distinctions.
\end{proof}

\begin{corollary}[Finite tested ingredients are not independent inputs]\label{cor:finite-tested-not-inputs}
For any physically executed finite APF interrogation protocol, the data used in the Sep/IJC theorem---finite $Q$, finite tested words $\mathcal T$, finite depth, finite operational resolution, and a positive floor for resolved tested distinctions---are induced by the protocol and quotient normal form.  They need not be added as separate physical postulates.
\end{corollary}

\begin{remark}[What remains theorem-local]\label{rem:what-remains-theorem-local}
Theorem~\ref{thm:finite-tested-normal-form} derives the finite-tested regime for physically executed finite protocols.  The compact robust-interface hypotheses R1--R4 below are a controlled extension of the same logic to compact nonexplicit families.  No-added-resource, factorized-support, and ledger-independence assumptions are not primitives; they are local hypotheses used only for monotonicity, additivity, and independent-counting theorems.
\end{remark}

\section{Robust finite interfaces}\label{sec:robust-finite}

For explicitly finite APF interrogation protocols, Theorem~\ref{thm:finite-tested-normal-form} already gives a positive floor for the resolved tested distinctions.  The present section gives the compact-family extension: when one studies a nonexplicit but compact queried family, R1--R4 are the regularity hypotheses under which the same positive-floor conclusion remains valid.

\begin{assumption}[R1--R4 robust finite interface]
A context $\Gamma$ is a \emph{robust finite interface} for a queried family $Q$ when:
\begin{enumerate}[label=R\arabic*.,leftmargin=2em]
\item The queried family $Q$ is finite, or is a nonempty compact subset of an admissible generator family equipped with the topology induced by the operational resolution of $\Gamma$.
\item The family contains only robust distinctions, i.e. distinctions stable under the admissible perturbation class of $\Gamma$.
\item The enforcement-cost map $\kappa_\Gamma$ is lower semicontinuous on the relevant generator family.
\item The available capacity is finite: $C_\Gamma<\infty$.
\end{enumerate}
\end{assumption}

\paragraph{Why R1--R4 are stated (and why they are not load-bearing).}
Each robust-interface condition addresses a specific failure mode that arises if one tries to derive the positive floor purely from the supplement's primitive spine (positive cost on each physical distinction) without invoking MD's uniform floor.  Without R1, one might worry about a sequence of resolved distinctions $d_n$ with costs $\kappa_\Gamma(d_n)=1/n$, accumulating to zero.  Without R3, costs might jump downward along a convergent sequence and the infimum might not be attained.  R2 and R4 are about physical realizability and counting respectively.

\emph{However, MD precludes the central failure mode directly.}  The Principle of Least Enforcement Cost names MD --- a uniform positive cost floor $\varepsilon^*_\Gamma > 0$ on every physical distinction --- as one of its four constitutive features (\S\ref{sec:purpose-status}).  MD is a derivable consequence of the supplement's primitive spine under the continuation-calculus reading of Paper~10 v1.12 \S3.5: the marginal-floor definition (Remark~\ref{rem:marginal-floor}) plus the finite-physical-regime hypothesis ($\varepsilon^*_\Gamma > 0$ as the second half of the regime hypothesis).  Once MD is in scope, the failure mode ``$\kappa(d_n) = 1/n$'' cannot occur: every $d_n$ in the queried family is a physical distinction, hence $\kappa_\Gamma(d_n) \ge \varepsilon^*_\Gamma > 0$ uniformly.  No compactness, no lower semicontinuity, no Weierstrass-style infimum-attainment theorem is needed; the positive floor follows directly from MD's uniform lower bound.

R1--R4 therefore \emph{are not load-bearing} for Theorem~\ref{thm:minimum-distinction-floor}.  They constitute a sufficient alternative regularity package for an audit-style proof route that does not invoke MD directly, but the structural source of the positive floor is MD itself.  We state R1--R4 to make the alternative route auditable and to clarify which conditions would be needed if MD were not in scope; the main proof of Theorem~\ref{thm:minimum-distinction-floor} below uses MD's uniform floor and does not invoke R1--R4.

\subsection{What R1--R4 mathematically import (and what they derive from)}\label{subsec:R1-R4-mathematical-import}

R1--R4 are stated above in operational language as if they were four added regularity hypotheses on top of the spine.  This subsection makes two points: (i) each of R1--R4 is a derivable consequence of the spine plus the operational structure of physical interrogation, not an additional commitment; and (ii) once MD is in scope (which it is, via Paper~10 v1.12 \S3.5 reductions of the supplement's primitive spine), \emph{no} mathematical content beyond the spine is imported.  Theorem~\ref{thm:minimum-distinction-floor} above uses MD's uniform floor directly.  R1--R4 give a sufficient alternative regularity package for a Weierstrass-style proof route, and the Weierstrass theorem is described below as the content one would need to import \emph{if MD were declined} --- but the main proof does not invoke it.  The supplement's stated stance --- ``physics-first, math is descriptive'' --- is honored.  R1--R4 do not add new physics, and once MD is read into the spine they do not require new mathematics either.

\paragraph{R1 (compactness or finiteness): automatic for finite interrogation; operationally automatic for compact instruments.}
For an explicitly finite APF interrogation protocol, R1 is automatic from Theorem~\ref{thm:finite-tested-normal-form}: the protocol induces a finite query family by construction, and finiteness gives compactness vacuously.  The compact-family case is invoked when one studies a continuous parameter family (e.g., a calibration parameter ranging over a closed bounded interval, an instrument family parametrized by a closed bounded subset of $\mathbb R^n$, or a finite-dimensional simplex of stochastic matrices as in Example~\ref{ex:finite-noisy-instrument-r1r4}).  In such cases compactness is operationally automatic: physical instruments are always parametrized by closed bounded subsets of finite-dimensional Euclidean spaces, and the operational-resolution-induced topology on those subsets makes compactness a consequence of the parameter space's physical specification, not an additional mathematical hypothesis.  The mathematical content invoked is the Bolzano--Weierstrass property (every sequence has a convergent subsequence) on the operational topology --- which is the topology induced by the operational resolution $\delta$: two generators are close when their induced readout distributions on the queried family differ by less than $\delta$ in some operationally-meaningful metric (total variation, $L^p$, or a calibration-specific norm).  No measure-theoretic or functional-analytic machinery is invoked beyond the topology itself.

\paragraph{R2 (robustness): built into the spine's definition of physical distinction.}
R2 is not an additional hypothesis at all.  The Definition of physical distinction in \S\ref{sec:distinctions-cost} already requires that every physical distinction admit a finite admissible enforcement protocol \emph{maintaining the separated alternatives against admissible perturbations}.  Stability under admissible perturbation is therefore the spine's primitive content: any object that fails this stability is not a physical distinction in APF and is not in $\mathcal D_\Gamma$ to begin with.  R2 is the operational restatement of this primitive at the queried-family level: the family $Q$ contains only physical distinctions, all of which are robust by definition.  No machinery beyond the spine is added.  Mathematically, R2 reads as an open-neighborhood operational stability in the topology of R1 (every $d \in Q$ has a neighborhood $U_d$ where the same physical distinction is recognized), which is a consequence of the operational-resolution structure.  R2 is what prevents the supplement from counting boundary or limit-only distinctions that have no stable physical realization, but the spine already excludes those.

\paragraph{R3 (lower semicontinuity): structurally automatic from the cost-positive-only-on-physical structure.}
R3 is structurally automatic for APF cost functionals.  The cost map $\kappa_\Gamma : \mathcal D_\Gamma \to (0,\infty)$ takes positive values exactly on physical distinctions --- those for which a finite admissible enforcement protocol exists.  At boundary or limit points where no such protocol exists, the distinction is unresolved and $\kappa_\Gamma$ is either zero or undefined (the boundary is not in $\mathcal D_\Gamma$).  As Example~\ref{ex:lsc-not-continuous-cost} shows explicitly, this gives LSC by construction: cost can jump \emph{upward} at the boundary where a distinction becomes physically enforceable (acquiring the stable-record floor $c_0 > 0$), but it cannot jump \emph{downward}, because every physical distinction has positive cost bounded below.  Standard real-analysis $\liminf$:
\[
    \kappa_\Gamma(d) \;\le\; \liminf_{n\to\infty} \kappa_\Gamma(d_n)
\]
holds vacuously at boundary points (where $\kappa_\Gamma(d) = 0$) and by continuity in the cost calibration on the physical interior.  Continuity itself is not required and would be too strong: the upward-jump pattern is the physical content of ``becoming enforceable,'' which LSC captures and continuity would forbid.  R3 is therefore not an added hypothesis but the standard real-analysis articulation of a structural fact about APF cost functionals.

\paragraph{R4 (finite capacity): A1 itself.}
R4 is just A1 --- the finite-physical-regime hypothesis already named in the spine (and again in Paper~10 v1.12 \S3.5: A1 is one half of the finite-physical-regime hypothesis $C_\Gamma < \infty$).  No new commitment is added; R4 is the spine's primitive content under a different name.  Mathematically it invokes only the standard ordering on $\mathbb R$.

\paragraph{The Weierstrass theorem as alternative-route import (not invoked in the main proof).}
With R2 built into the spine, R3 structurally automatic, R4 identical to A1, and MD's uniform floor in scope (Paper~10 v1.12 \S3.5), Theorem~\ref{thm:minimum-distinction-floor} follows directly: $\mu_\Gamma(Q) \ge \varepsilon^*_\Gamma > 0$ pointwise on physical distinctions.  No mathematical content beyond the spine is invoked.

R1--R4 nevertheless support a Weierstrass-style \emph{alternative} proof route (Remark~\ref{rem:weierstrass-alternative-route}): a lower semicontinuous function on a (sequentially or topologically) compact set attains its infimum, and the attained infimum is positive because every value of $\kappa_\Gamma$ on a physical distinction is strictly positive.  The Weierstrass extremum theorem in its continuous form, with the LSC generalization (see, e.g., any standard real-analysis text covering Tonelli or the direct method of the calculus of variations), is the mathematical content the alternative route would import.  We state R1--R4 to make this alternative route auditable: a reader who declines to import MD from the broader corpus can verify the floor theorem via Weierstrass + R1--R3.  The main proof above does not require it.

\paragraph{The physics-first reading made explicit.}
The supplement's stated stance is that physics specifies the operational structure (finite resolution, robust distinctions, finite capacity) and mathematics describes the consequences of that structure.  Once MD is in scope --- which it is, since MD is a derivable consequence of the supplement's primitive spine via Paper~10 v1.12 \S3.5 reductions, and is invoked elsewhere in this supplement (\S\ref{sec:composition-factorization}, marginal-floor lemma) --- the floor theorem requires no real-analysis machinery whatsoever:

\begin{itemize}[leftmargin=2em,topsep=4pt,itemsep=4pt]
\item For any $Q\subset\mathcal D_\Gamma$ (finite, compact, noncompact, with or without LSC, with or without robustness conditions), $\mu_\Gamma(Q) \ge \varepsilon^*_\Gamma > 0$ by MD's uniform floor applied pointwise.  The positive floor is the floor on every physical distinction, not a property of the family's topology.
\item R1--R4 give a sufficient alternative regularity package for a Weierstrass-style proof route that does not invoke MD directly.  The Weierstrass infimum-attainment theorem is the mathematical content this alternative route imports.  We retain R1--R4 to make the alternative route auditable but do not invoke them in the main proof.
\item When $Q$ is explicitly finite (the typical APF interrogation protocol), the floor is the trivial minimum of finitely many positive costs.  Both proof routes collapse to inspection.
\item When $Q$ is compact-but-not-finite (e.g., a continuous calibration parameter, an instrument-family parametrized by a closed bounded subset of $\mathbb R^n$, or the finite-dimensional simplex of Example~\ref{ex:finite-noisy-instrument-r1r4}), the MD route still gives the floor pointwise; R1--R4 + Weierstrass is the alternative.
\item Outside the spine itself --- e.g., regimes where MD fails because the finite-physical-regime hypothesis fails ($\varepsilon^*_\Gamma = 0$ or $C_\Gamma = \infty$) --- the supplement makes no claim and the positive floor argument is unavailable.  Such regimes lie outside the scope of the present spine.
\end{itemize}

The honest summary is short: with MD's uniform floor in scope (and it is, as a derivable consequence of the spine), Theorem~\ref{thm:minimum-distinction-floor} requires no real-analysis machinery at all.  R1--R4 do not add four regularity hypotheses to the spine; they are derivable consequences of the spine plus the operational structure of physical interrogation, and they support a Weierstrass-style alternative proof route for readers who decline the MD import.  A reader who accepts MD has nothing further to import.  A reader who declines MD imports a single standard real-analysis theorem (Weierstrass).  Either way, the supplement stands at the ``physics-first, math is descriptive'' level: the heavy lifting is done by the spine's primitive content, not by topology.

\begin{theorem}[Minimum distinction floor]\label{thm:minimum-distinction-floor}
For any nonempty queried family $Q\subset\mathcal D_\Gamma$ of physical distinctions in a finite physical regime, there exists a positive floor
\[
 \mu_\Gamma(Q):=\inf_{d\in Q}\kappa_\Gamma(d) \;\ge\; \varepsilon^*_\Gamma \;>\; 0.
\]
\end{theorem}
\begin{proof}
By MD (the second half of the finite-physical-regime hypothesis; Paper~10 v1.12 \S3.5 Lemma~BW; equivalently, the marginal floor of Remark~\ref{rem:marginal-floor} plus its positivity), every physical distinction $d \in \mathcal D_\Gamma$ satisfies $\kappa_\Gamma(d) \ge \varepsilon^*_\Gamma > 0$.  Since $Q\subset\mathcal D_\Gamma$ is nonempty and contains only physical distinctions, this uniform lower bound applies pointwise: $\kappa_\Gamma(d) \ge \varepsilon^*_\Gamma$ for every $d\in Q$.  Taking the infimum, $\mu_\Gamma(Q) \ge \varepsilon^*_\Gamma > 0$.

No compactness, lower semicontinuity, or infimum-attainment theorem is invoked.  The positive floor follows from MD's uniform lower bound on every physical distinction; it does not depend on any topological structure on $Q$.
\end{proof}

\begin{remark}[Alternative proof route via R1+R3]\label{rem:weierstrass-alternative-route}
A reader who declines to import MD from the broader corpus may obtain Theorem~\ref{thm:minimum-distinction-floor} via R1--R3 and the standard Weierstrass-style infimum-attainment theorem: a lower semicontinuous function on a compact set attains its infimum, and the attained infimum is positive because every value of $\kappa_\Gamma$ on a physical distinction is strictly positive.  This alternative route is internally consistent, but it requires importing one nontrivial real-analysis theorem (the LSC form of Weierstrass) and importing a topology on the generator family.  The MD route above imports neither.  Both routes give the same conclusion; we use the MD route as the primary proof and retain R1--R4 to support the alternative.
\end{remark}

\begin{corollary}[Finite jointly enforceable independent distinctions]\label{cor:finite-independent-distinctions}
If every independent queried distinction in $Q$ costs at least $\mu_\Gamma(Q)>0$ and total capacity is $C_\Gamma$, then any jointly enforceable independent family has cardinality at most
\[
 N_\Gamma(Q)\le \left\lfloor \frac{C_\Gamma}{\mu_\Gamma(Q)}\right\rfloor .
\]
\end{corollary}
\begin{proof}
Let $S$ be a jointly enforceable independent family.  By the distinction floor, each $d\in S$ satisfies $\kappa_\Gamma(d)\ge\mu_\Gamma(Q)$.  Independence means the joint ledger cost is at least the sum of the individual supports, so $\kappa_\Gamma(S)\ge |S|\mu_\Gamma(Q)$.  Joint enforceability gives $\kappa_\Gamma(S)\le C_\Gamma$.  Therefore $|S|\le C_\Gamma/\mu_\Gamma(Q)$, and integrality gives the displayed floor bound.
\end{proof}

\begin{remark}[Verifying R1--R4]
The robust-finite hypotheses are satisfied in standard finite operational models with explicit nonzero resolution thresholds: finite classical registers, finite detector arrays, finite-state control systems, and finite-dimensional noisy instrument models with closed bounded effect sets and continuous cost functionals.  They are not claimed for arbitrary idealized infinite systems without additional limiting arguments.
\end{remark}


\begin{example}[Finite noisy instruments verifying R1--R4]\label{ex:finite-noisy-instrument-r1r4}
Let the admissible generator family consist of binary finite instruments on a finite input alphabet $X$ and output alphabet $Y$.  Each instrument is represented by a stochastic matrix $P(y\mid x)$, so the generator family is a closed subset of a finite product of simplices and is compact in the Euclidean topology.  Fix an operational resolution $\delta>0$ and declare two continuation profiles resolved only when their induced readout distributions differ by at least $\delta$ in total variation distance.

R1 holds because the admissible family is compact.  R2 holds after quotienting all pairs separated by less than $\delta$: remaining distinctions are stable under perturbations smaller than $\delta/3$.  R3 holds for any calibration such as
\[
  \kappa_\Gamma(d)=c_0+c_1\,N_{\delta,\alpha}(d),
  \qquad c_0,c_1>0,
\]
where $N_{\delta,\alpha}$ is a sufficient sample size bound chosen as an upper semicontinuous integer envelope; the positive constant $c_0$ supplies the stable-record floor and prevents accumulation to zero.  R4 is the experimental capacity budget: finite trials, finite stable memory, and finite permitted control actions.  Thus this standard finite noisy-instrument context satisfies the robust finite-interface hypotheses without assuming Hilbert space.
\end{example}


\begin{example}[Lower semicontinuous but not continuous cost]\label{ex:lsc-not-continuous-cost}
Let a one-parameter family of binary discriminators be indexed by a stability margin $m\in[0,1]$, where $m=0$ is unresolved and $m>0$ is a resolved distinction.  Define
\[
  \kappa_\Gamma(m)=
  \begin{cases}
  0, & m=0,\\
  c_0+c_1 m, & m>0,
  \end{cases}
  \qquad c_0>0.
\]
On the physical resolved family $m\ge m_*>0$ this cost is continuous and has a positive floor.  On the closure including the unresolved boundary, it is lower semicontinuous at $m=0$ because
\[
  \kappa_\Gamma(0)\le \liminf_{m\downarrow0}\kappa_\Gamma(m)=c_0,
\]
but it is not continuous there.  This illustrates why R3 asks for lower semicontinuity rather than unnecessary continuity: costs may jump upward when a distinction becomes physically enforceable, while unresolved boundary points are not counted as positive-cost physical distinctions.
\end{example}

\section{Filters, cost monotonicity, and convex control}\label{sec:filters-convexity}

\begin{definition}[Admissibility-preserving filter]
A filter $F:\A_\Gamma\to\A_{\Gamma'}$ is an induced quotient map
\[
 F:\Omega_\Gamma\to\Omega_{\Gamma'}
\]
that maps admissible continuations of $\Gamma$ to admissible or coarser continuations of $\Gamma'$.  It may erase, coarse-grain, restrict, or identify possibilities, but it may not invent distinctions unsupported by finite admissible continuation.
\end{definition}

\begin{definition}[Coarse-graining filter]
A filter $F$ is a coarse-graining when every distinction $d'$ in the image context pulls back to a distinction $F^*d'$ in the source context by
\[
 (F^*d')([x])=d'(F[x]).
\]
\end{definition}

\begin{assumption}[No-added-resource coarse-graining]\label{ass:no-added-resource}
A coarse-graining filter $F:\A_\Gamma\to\A_{\Gamma'}$ is \emph{no-added-resource} when it only erases, identifies, forgets, restricts, or lowers the resolution of continuation distinctions.  It does not introduce fresh ancillary records, active amplification, error correction, adaptive control registers, external stabilization, or catalytic enforcement resources in the target context.
\end{assumption}

The no-added-resource condition is a theorem condition.  If a map adds an amplifier, ancilla, error-correcting code, control register, or stabilizer, that resource must be charged to the target ledger before comparing costs.


\begin{definition}[Register-channel filter model]\label{def:register-channel-no-added}
A concrete no-added-resource filter is a finite protocol
\[
 (R,\sigma_R)\xrightarrow{\Phi} (R',\sigma_{R'})
\]
where $R$ is a finite record register supporting distinctions in context $\Gamma$, $R'$ is obtained from $R$ by a stochastic channel, deterministic relabeling, deletion of coordinates, coarse binning, or restriction to a subregister, and no new ancilla record, amplifier, stabilization bath, error-correction register, or fresh synchronization resource is introduced.  The induced distinction in the filtered context is admissible only when its pullback to $R$ was already supported by the original record protocol.
\end{definition}

\begin{definition}[Ledger resource preorder]\label{def:ledger-resource-preorder}
For each context $\Gamma$, let $\mathsf{Prot}_\Gamma(d)$ denote the
finite enforcement protocols that realize a distinction
$d\in\mathcal D_\Gamma$ at the stated operational resolution.  Each
protocol $P\in\mathsf{Prot}_\Gamma(d)$ carries a resource load
\[
    r_\Gamma(P)\in\mathcal R_\Gamma,
\]
where $(\mathcal R_\Gamma,\oplus,0,\le)$ is a commutative preordered
resource monoid.  The preorder $a\le b$ means that resource load $a$
is no greater than resource load $b$ in the chosen context calibration.

A \emph{context calibration} is a monotone valuation
\[
    \chi_\Gamma:\mathcal R_\Gamma\to\R_{\ge 0},
\]
i.e.\ a map satisfying $a\le b\Rightarrow \chi_\Gamma(a)\le\chi_\Gamma(b)$
and $\chi_\Gamma(0)=0$.  The cost of a distinction is the minimum, or
infimum when a minimum is not attained, of the calibrated load over
its realizing protocols:
\[
    \kappa_\Gamma(d)
    \;=\;
    \inf_{P\in\mathsf{Prot}_\Gamma(d)} \chi_\Gamma\bigl(r_\Gamma(P)\bigr).
\]
In the finite resolved setting used in the main theorem, this infimum
is a minimum after quotienting by operational resolution.  Comparisons
between costs in the no-added-resource simulation
(Definition~\ref{def:no-added-resource-simulation}) and the additivity
theorem (Theorem~\ref{thm:ledger-additivity-protocol-factorization})
are stated either directly in the preorder $\le$ on $\mathcal R_\Gamma$
or in the calibrated values $\chi_\Gamma(r_\Gamma(P))$, which the
context calibration relates monotonically.
\end{definition}

\begin{definition}[No-added-resource simulation]\label{def:no-added-resource-simulation}
Let $F:\A_\Gamma\to\A_{\Gamma'}$ be an admissibility-preserving
coarse-graining.  We say that $F$ is \emph{no-added-resource} for a
target distinction $d'\in\mathcal D_{\Gamma'}$ when there exists a
protocol-level simulation map
\[
    \Sigma_F:\mathsf{Prot}_\Gamma(F^*d')
    \longrightarrow
    \mathsf{Prot}_{\Gamma'}(d')
\]
such that, for every source protocol $P$,
\[
    r_{\Gamma'}(\Sigma_F(P)) \;\le\; r_\Gamma(P).
\]
Equivalently, every enforcement of the pulled-back distinction $F^*d'$
can be pushed forward to an enforcement of $d'$ without increasing
resource load.

A filter is no-added-resource on a finite family
$S'\subset\mathcal D_{\Gamma'}$ when the same condition holds for every
$d'\in S'$, and for every joint protocol realizing finite subfamilies
of $S'$.  This is the formal trigger that the operational
Assumption~\ref{ass:no-added-resource} captures in prose.
\end{definition}

\begin{proposition}[Passive register channels are no-added-resource]\label{prop:register-channels-no-added-resource-formal}
A register-channel filter $R\xrightarrow{K}R'$ implemented by stochastic
post-processing, deterministic relabeling, coordinate deletion, binning,
or restriction, with no fresh persistent register, amplifier,
stabilizer, error-correcting redundancy, clock, reference frame, or
adaptive controller, induces a no-added-resource simulation in the
sense of Definition~\ref{def:no-added-resource-simulation}.
\end{proposition}
\begin{proof}
A passive channel $K$ extends to a protocol-level map $\Sigma_F$ that
applies $K$ to the readout coordinates of any source protocol while
leaving its persistent records, amplifiers, stabilizers, and control
registers untouched.  Because $K$ adds no maintained register or
catalytic resource, the image protocol's resource load satisfies
$r_{\Gamma'}(\Sigma_F(P)) \le r_\Gamma(P)$ in any calibration that
charges only persistent maintained resources.  Hence $\Sigma_F$ is a
no-added-resource simulation in the formal sense.
\end{proof}

\begin{proposition}[Operational sufficient tests for no-added-resource maps]\label{prop:no-added-resource-sufficient-tests}
A filter is safely no-added-resource for the distinctions under comparison if it admits an implementation as passive post-processing of already maintained records:
\[
  R \xrightarrow{K} R',
\]
where $K$ is a stochastic channel, deterministic relabeling, coordinate deletion, binning map, or restriction map acting only on the source register $R$, and where the implementation uses no fresh persistent memory, no active error correction, no amplification stage, no new synchronization reference, and no external stabilization bath whose state must be maintained to realize the target distinction.  Under these sufficient conditions, every target discriminator pulls back to a discriminator already implementable on $R$.
\end{proposition}
\begin{proof}
A passive channel $K$ cannot increase the sigma-algebra of finite records supported by $R$; it can only merge, relabel, randomize, or discard already recorded alternatives.  Therefore any target event in $R'$ has a preimage event in $R$.  Since no additional maintained register or stabilizer is introduced, the preimage event is supported by the original ledger resources.  This is exactly the no-added-resource condition needed in the monotonicity theorem.
\end{proof}

\begin{remark}[Checklist: certifying no-added-resource]\label{check:no-added-resource}
For experimental use, treat a filter as no-added-resource only after checking the following sufficient conditions.
\begin{enumerate}[label=(\roman*),leftmargin=2em]
\item The map is a declared post-processing channel on existing records, not a new measurement.
\item It adds no persistent register, clock, reference frame, memory cell, ancilla, error-correcting redundancy, amplifier, or feedback controller.
\item Its output alphabet is a quotient, subalphabet, or randomized relabeling of the input record alphabet.
\item The target distinction's pullback is measurable using the source record alone.
\item Any laboratory hardware needed to implement the map is either already included in the source context or is explicitly charged to the target ledger before costs are compared.
\end{enumerate}
Failure of any item does not make the map inadmissible; it means the map is not free for the monotonicity theorem.
\end{remark}

\begin{remark}[Why monotonicity is not a free axiom]\label{rem:cost-monotonicity-model}
In the register-channel model, a filter can forget, merge, relabel, or restrict records, but it cannot create a separator that was not already maintained by the source register.  Thus any implementation of a filtered distinction can be pulled back to an implementation before filtering.  Ambiguous cases are exactly those where a passive-looking filter hides stabilization, amplification, error correction, or a new reference register; those are not no-added-resource filters and must be charged explicitly in the ledger.
\end{remark}
\begin{theorem}[Cost monotonicity under no-added-resource simulation]\label{thm:cost-monotonicity-no-added}
Let $F:\A_\Gamma\to\A_{\Gamma'}$ be an admissibility-preserving
coarse-graining, and let $d'\in\mathcal D_{\Gamma'}$ with
$F^*d'\in\mathcal D_\Gamma$.  If $F$ is no-added-resource for $d'$ in
the sense of Definition~\ref{def:no-added-resource-simulation}, then
\[
 \kappa_{\Gamma'}(d')\;\le\;\kappa_\Gamma(F^*d').
\]
\end{theorem}
\begin{proof}
For every protocol $P\in\mathsf{Prot}_\Gamma(F^*d')$, the
no-added-resource simulation gives a protocol
\[
    \Sigma_F(P)\in\mathsf{Prot}_{\Gamma'}(d')
\]
with
\[
    r_{\Gamma'}(\Sigma_F(P))\;\le\;r_\Gamma(P).
\]
Taking the infimum over all protocols $P$ realizing $F^*d'$ gives
\[
    \kappa_{\Gamma'}(d')
    \;\le\;
    \inf_{P\in\mathsf{Prot}_\Gamma(F^*d')} r_{\Gamma'}(\Sigma_F(P))
    \;\le\;
    \inf_{P\in\mathsf{Prot}_\Gamma(F^*d')} r_\Gamma(P)
    \;=\;
    \kappa_\Gamma(F^*d').
\]
Thus cost cannot increase under a resource-nonincreasing simulation.
If a map uses active amplification, new ancillas, adaptive controls,
or error correction, then it does not admit a no-added-resource
simulation in the sense above; those resources must be included
explicitly in the target ledger before any monotonicity comparison is
made.
\end{proof}


\begin{example}[Resource-adding filter: why monotonicity can fail]\label{ex:resource-adding-filter-counterexample}
Start with a noisy binary record $R\in\{0,1\}$ whose readout error is $p=0.15$.  A passive coarse-graining that forgets $R$ or merges outcomes is no-added-resource and cannot make a harder distinction cheaper.  Now compare a different map that first copies the signal into $m$ fresh ancilla bits, majority-votes them, and outputs a cleaned bit $R'$.  The target distinction may have lower error and lower apparent sample cost than the source distinction, but the map has introduced fresh memory, amplification, and error correction.  It is therefore not a no-added-resource coarse-graining.  APF does not forbid the construction; it requires the ancilla, control, and stabilization costs to be charged in the target ledger before any monotonicity comparison is made.
\end{example}

\begin{definition}[Admissible classical control]
A finite classical control register $R=\{1,\ldots,m\}$ is admissible in context $\Gamma$ if its alternatives are finitely recordable and its total record cost lies within the capacity budget.
\end{definition}

\begin{theorem}[Convexity from admissible classical control]
If preparations $\omega_1,\ldots,\omega_m$ are admissible and a classical control register selects preparation $\omega_i$ with frequency weights $p_i\ge0$, $\sum_i p_i=1$, then the operational preparation
\[
 \omega=\sum_{i=1}^m p_i\omega_i
\]
is admissible whenever the control record and selected preparation fit within the ledger budget.
\end{theorem}
\begin{proof}
The control register is itself a finite record.  Conditional on record value $i$, the preparation is $\omega_i$.  Coarse-graining over the control record gives observed statistics equal to the weighted average of the conditional statistics.  Thus the operational preparation is the convex mixture shown.
\end{proof}

\paragraph{Status of the trigger conditions.}
The trigger conditions above are not additional physical laws.  They are
sufficient structural hypotheses under which particular ledger
inequalities are licensed.  No theorem in this supplement assumes that
arbitrary filters are no-added-resource, that arbitrary composites are
additive, or that arbitrary resolved word classes are independently
counted.  When the protocol-level simulation, support factorization, or
support-independence hypothesis is absent, the corresponding
monotonicity, additivity, or counting theorem is simply unavailable.

\section{Composition, factorization, and interface terms}\label{sec:composition-factorization}

\begin{definition}[Admissible composition]
Given admissible possibility spaces $\A_{\Gamma_1}$ and $\A_{\Gamma_2}$, their composite $\A_{\Gamma_1\circ\Gamma_2}$ is defined where there exists a context in which representatives of both spaces can be jointly continued, queried, and recorded without exceeding the available capacity.
\end{definition}

\begin{definition}[Factorized continuation]
A composite is factorized at $([x_1],[x_2])$ when
\[
 \Cont_{12}(x_1,x_2)
 \cong
 \Cont_1(x_1)\times\Cont_2(x_2)
\]
up to the quotient identity of the composite context.
\end{definition}

\begin{definition}[Coupled composition and interface cost]
A composite is coupled when its continuation set does not factorize.  For a family of distinctions $S_1\cup S_2$ supported on the two components, write
\begin{equation}\label{eq:kappa-12-interface}
 \kappa_{12}(S_1,S_2)
 =\kappa_1(S_1)+\kappa_2(S_2)+\kappa_{\mathrm{int}}(S_1,S_2).
\end{equation}
The term $\kappa_{\mathrm{int}}$ is the interface cost.
\end{definition}

\begin{definition}[Enforcement-factorized support]\label{def:enforcement-factorized-support}
Let $S_1\subset\mathcal D_{\Gamma_1}$ and $S_2\subset\mathcal D_{\Gamma_2}$.
Their enforcement support is \emph{factorized} in the composite context
$\Gamma_{12}$ when:
\begin{enumerate}[label=(\roman*),leftmargin=2em]
\item every joint protocol realizing $S_1\cup S_2$ decomposes, up to
contextual equivalence, as a tensor/disjoint union
\[
    P_{12}\;\simeq\;P_1\boxtimes P_2;
\]
\item the resource monoid factorizes:
\[
    \mathcal R_{\Gamma_{12}}
    \;\cong\;
    \mathcal R_{\Gamma_1}\oplus\mathcal R_{\Gamma_2};
\]
\item resource load is additive on factorized protocols:
\[
    r_{\Gamma_{12}}(P_1\boxtimes P_2)
    \;=\;
    r_{\Gamma_1}(P_1)+r_{\Gamma_2}(P_2);
\]
\item no shared stabilizer, record, memory, reference frame, control
bus, or interface protocol contributes a nonzero joint term.
\end{enumerate}
\end{definition}

\begin{theorem}[Ledger additivity under protocol factorization]\label{thm:ledger-additivity-protocol-factorization}
If $S_1$ and $S_2$ have enforcement-factorized support in $\Gamma_{12}$
in the sense of Definition~\ref{def:enforcement-factorized-support},
then
\[
    \kappa_{\Gamma_{12}}(S_1\cup S_2)
    \;=\;
    \kappa_{\Gamma_1}(S_1)+\kappa_{\Gamma_2}(S_2),
\]
equivalently $\kappa_{\mathrm{int}}(S_1,S_2)=0$ in
\eqref{eq:kappa-12-interface}.
\end{theorem}
\begin{proof}
By support factorization (i)--(iii), any admissible joint protocol is
equivalent to $P_1\boxtimes P_2$, and its resource load is
\[
    r_{\Gamma_{12}}(P_1\boxtimes P_2)
    \;=\;
    r_{\Gamma_1}(P_1)+r_{\Gamma_2}(P_2).
\]
Taking the infimum over all factorized realizing protocols gives
\[
    \kappa_{\Gamma_{12}}(S_1\cup S_2)
    \;=\;
    \inf_{P_1,P_2}\bigl(r_{\Gamma_1}(P_1)+r_{\Gamma_2}(P_2)\bigr)
    \;=\;
    \kappa_{\Gamma_1}(S_1)+\kappa_{\Gamma_2}(S_2).
\]
By condition (iv), no shared resource contributes a positive joint
term; equivalently $\kappa_{\mathrm{int}}(S_1,S_2)=0$.  The equality
fails exactly when the factorization hypotheses fail, in which case
the difference is recorded as the interface term $\kappa_{\mathrm{int}}$.
\end{proof}

\begin{definition}[Support-independent resolution]\label{def:support-independent-resolution}
Let $W_1,\ldots,W_N$ be operationally resolved continuation-word
classes, with witnessing discriminators $r_1,\ldots,r_N$.  The
witnesses are \emph{support-independent} when there exist resource
supports $\sigma_i\subseteq\Sigma_\Gamma$ such that:
\begin{enumerate}[label=(\roman*),leftmargin=2em]
\item $r_i$ cannot be realized at the stated resolution without support
on $\sigma_i$;
\item the supports are pairwise disjoint or resource-orthogonal:
$\sigma_i\perp\sigma_j$ for $i\ne j$;
\item resource load is additive across the supports:
\[
    r_\Gamma(P_1\boxtimes\cdots\boxtimes P_N)
    \;=\;
    \sum_{i=1}^N r_\Gamma(P_i);
\]
\item no single coarser record, relabeling, or uncharged control
variable realizes all $r_i$ at the same resolution.
\end{enumerate}
\end{definition}

\paragraph{Signs of the interface term.}
The interface term $\kappa_{\mathrm{int}}(S_1,S_2)$ in
\eqref{eq:kappa-12-interface} records the difference between the
least joint enforcement protocol and the formal sum of separately
maintained protocols.  It admits the following typology.  It vanishes,
$\kappa_{\mathrm{int}}=0$, when the supports factorize per
Definition~\ref{def:enforcement-factorized-support}.  It is positive,
$\kappa_{\mathrm{int}}>0$, when joint enforcement requires
synchronization, stabilization, a shared reference frame, an
interface memory, or active recovery not needed by either subsystem
alone.  It may be negative, $\kappa_{\mathrm{int}}<0$, when two
distinctions share an already maintained record or stabilizer, so
that the least joint protocol is cheaper than maintaining two
independent supports.  A negative interface term is not a violation
of positive cost; it means that the joint implementation shares
support, and independent-counting bounds are licensed only after
support independence has been verified.

\subsection*{Structural shape of the joint cost and the interface term}

The interface term $\kappa_{\mathrm{int}}$ in \eqref{eq:kappa-12-interface} (and the singleton-decomposition residue $\kappa_{\Gamma,\mathrm{int}}(S)$ in the ledger admissibility definition of \S\ref{sec:distinctions-cost}) was introduced by residue: it absorbs whatever deviation from additivity the joint protocol exhibits.  The residue framing is honest, but it leaves the structural shape of $\kappa_{\Gamma,\mathrm{int}}$ apparently free.  We now show that the joint cost $\kappa_\Gamma(S)$ inherits four structural inequalities directly from the spine --- the marginal-floor definition of \S\ref{sec:distinctions-cost}, the positivity of the marginal floor in any finite physical regime, and the infimum-over-admissible-protocols valuation convention (Remark~\ref{rem:ledger-valuation-convention}).  Together these inequalities fix the structural shape of $\kappa_{\Gamma,\mathrm{int}}$ on its lower-bound side.  No new commitment beyond the spine is added.

\begin{lemma}[Marginal-floor lower bound on every admissible increment]
\label{lem:marginal-floor-on-joint-cost}
Let $S\subset\mathcal D_\Gamma$ be admissible and let $d\in\mathcal D_\Gamma$ be a tested distinction with $S\cup\{d\}$ also admissible.  Then
\[
    \kappa_\Gamma(S\cup\{d\}) - \kappa_\Gamma(S)
    \;\ge\;
    \varepsilon^*_\Gamma.
\]
\end{lemma}
\begin{proof}
By Remark~\ref{rem:marginal-floor}, the marginal floor is
\[
    \varepsilon^*_\Gamma
    \;=\;
    \inf\bigl\{\kappa_\Gamma(S\otimes d) - \kappa_\Gamma(S) :
            S \text{ admissible},\, d \text{ tested},\,
            S\otimes d \text{ admissible}\bigr\}.
\]
The displayed inequality follows from the definition of infimum.  The substantive content is that $\varepsilon^*_\Gamma > 0$ in any finite physical regime, which is the cost-resolution reading of MD given by Paper~10 v1.12 \S3.5, Lemma~BW: the cost spectrum on tested distinctions is graded at scale $\varepsilon^*$, so there is no admissible move that resolves a tested distinction for an increment $< \varepsilon^*$.  Without this positivity the bound would be trivial; with it, every admissible increment costs at least the floor.
\end{proof}

\begin{corollary}[Sum-of-floors lower bound on the joint cost]
\label{cor:sum-of-floors-lower-bound}
For an admissible family $S\subset\mathcal D_\Gamma$ with $|S|=n$,
\[
    \kappa_\Gamma(S) \;\ge\; n\,\varepsilon^*_\Gamma.
\]
\end{corollary}
\begin{proof}
Order $S=\{d_1,\ldots,d_n\}$ and build it by single-element extension:
\[
    \emptyset \;\subset\; \{d_1\} \;\subset\; \{d_1,d_2\} \;\subset\; \cdots \;\subset\; S,
\]
each intermediate family admissible (any sub-family of an admissible family is admissible by the infimum-over-admissible-protocols definition: a least-cost admissible protocol for $S$ realizes every sub-family as a sub-resolution).  Each increment satisfies the marginal-floor bound (Lemma~\ref{lem:marginal-floor-on-joint-cost}).  Telescoping, $\kappa_\Gamma(S) - \kappa_\Gamma(\emptyset) \ge n\varepsilon^*_\Gamma$, and $\kappa_\Gamma(\emptyset)=0$ by convention.
\end{proof}

\begin{lemma}[Monotonicity of the joint cost in the family]
\label{lem:joint-cost-monotonicity-in-S}
If $S\subseteq S'\subset\mathcal D_\Gamma$ are both admissible, then $\kappa_\Gamma(S)\le\kappa_\Gamma(S')$.
\end{lemma}
\begin{proof}
Any admissible joint protocol realizing $S'$ also realizes $S$ as a sub-resolution: the protocol resolves every distinction in $S'$, hence in particular every distinction in $S\subseteq S'$.  Therefore the infimum of admissible-protocol costs realizing $S$ is at most the infimum of admissible-protocol costs realizing $S'$.
\end{proof}

\begin{lemma}[Coarse-graining monotonicity]
\label{lem:coarse-graining-monotonicity}
Let $S\subset\mathcal D_\Gamma$ be admissible and let $\pi:S\to S'$ be a contextual coarse-graining requiring no new resources beyond those already enforced for $S$ (i.e., $S'$ is a contextually-equivalent quotient of $S$ realized by reading off the coarser partition from any admissible protocol for $S$).  Then
\[
    \kappa_\Gamma(S')\;\le\;\kappa_\Gamma(S).
\]
\end{lemma}
\begin{proof}
Let $P^\star$ be a least-cost admissible protocol realizing $S$ at cost $\kappa_\Gamma(S)$.  Since $\pi$ requires no new resources, $P^\star$ also realizes $S'$ at the same cost: one reads off the coarser equivalence classes from the same record content.  Therefore the infimum of admissible-protocol costs realizing $S'$ is at most $\kappa_\Gamma(S)$.
\end{proof}

\begin{theorem}[Structural lower bound on the singleton-form interface term]
\label{thm:kappa-int-singleton-shape}
Let $S=\{d_1,\ldots,d_n\}$ be admissible and let $\kappa_{\Gamma,\mathrm{int}}(S)$ be the residue defined by
\(
    \kappa_\Gamma(S) = \sum_{d\in S}\kappa_\Gamma(d) + \kappa_{\Gamma,\mathrm{int}}(S).
\)
Then
\[
    \kappa_{\Gamma,\mathrm{int}}(S)
    \;\ge\;
    n\,\varepsilon^*_\Gamma \;-\; \sum_{d\in S}\kappa_\Gamma(d).
\]
The bound is attained when every per-distinction surplus $\kappa_\Gamma(d_i) - \varepsilon^*_\Gamma$ is fully amortized by shared substrate, i.e., when the joint cost reduces to $n\varepsilon^*_\Gamma$.
\end{theorem}
\begin{proof}
Substitute the sum-of-floors lower bound (Corollary~\ref{cor:sum-of-floors-lower-bound}) for $\kappa_\Gamma(S)$ into the residue equation and rearrange:
\[
    \kappa_{\Gamma,\mathrm{int}}(S)
    \;=\;
    \kappa_\Gamma(S) - \sum_{d\in S}\kappa_\Gamma(d)
    \;\ge\;
    n\varepsilon^*_\Gamma - \sum_{d\in S}\kappa_\Gamma(d).
\]
The bound is saturated exactly when $\kappa_\Gamma(S) = n\varepsilon^*_\Gamma$, the sum-of-floors floor itself.
\end{proof}
\coderef{check\_T\_kappa\_int\_lower\_bound}{kappa\_int\_bounds.py}

\begin{corollary}[Structural lower bound on the binary-form interface term]
\label{cor:kappa-int-binary-shape}
For the binary-decomposition form
\(
    \kappa_{12}(S_1,S_2) = \kappa_1(S_1) + \kappa_2(S_2) + \kappa_{\mathrm{int}}(S_1,S_2)
\)
of \eqref{eq:kappa-12-interface},
\[
    \kappa_{\mathrm{int}}(S_1,S_2)
    \;\ge\;
    (|S_1|+|S_2|)\,\varepsilon^*_\Gamma - \kappa_1(S_1) - \kappa_2(S_2).
\]
\end{corollary}
\begin{proof}
Apply Corollary~\ref{cor:sum-of-floors-lower-bound} to the joint family $S_1\cup S_2$ in the composite context $\Gamma_{12}$: $\kappa_{12}(S_1,S_2) \ge (|S_1|+|S_2|)\varepsilon^*_\Gamma$.  Subtract $\kappa_1(S_1)+\kappa_2(S_2)$ from both sides.
\end{proof}

\paragraph{Reading.}
The lower bound on $\kappa_{\Gamma,\mathrm{int}}$ encodes a sharp structural commitment: shared substrate may reduce the joint cost below the in-isolation independent-sum, but cannot reduce it below the sum of marginal floors $n\varepsilon^*_\Gamma$.  The marginal floor is the irreducible per-distinction cost regardless of sharing; the in-isolation surplus $\kappa_\Gamma(d) - \varepsilon^*_\Gamma$ is the amortizable piece.  This is the structural reading of MD lifted to the joint-family setting.  Although the singleton-form residue $\kappa_{\Gamma,\mathrm{int}}(S)$ depends only on the family $S$ while the binary-form residue $\kappa_{\mathrm{int}}(S_1,S_2)$ depends on the chosen decomposition, both are bounded below by the same sum-of-floors structural inequality applied to the underlying admissible family.

A complementary \emph{upper} bound on $\kappa_{\Gamma,\mathrm{int}}$ as a function of the substrate's coupling structure --- the cost of synchronization, shared reference frames, interface memory, and active recovery in regimes where joint enforcement requires more than parallel separate protocols --- is derived in \S\ref{subsec:kappa-int-upper-bound} via the continuum-bridge apparatus of \S\ref{sec:recruitment-radius}.  Together with the lower bound above, the two sides combine in Theorem~\ref{thm:kappa-int-two-sided-bound} (two-sided structural bound): in the continuum-bridge regime C1--C5, $\kappa_{\Gamma,\mathrm{int}}$ is bounded between the marginal-floor structural inequality below and the kernel-norm against profile-$L^1$-norms inequality above, with both endpoints structurally explicit.  Far-separated anchors get exponential suppression (Corollary~\ref{cor:kappa-int-far-separation}).  Outside this regime (saturation, far-from-equilibrium cosmogenesis, substrate collapse), the upper bound does not transfer automatically; the relevant substrate-mode work is supplied by Paper~25 \S\S\,2--3 and the Paper~6 supplement \S1.5.

\begin{theorem}[Independent-counting bound]\label{thm:independent-counting-bound-formal}
Assume a resolved distinction floor $\mu_\Gamma>0$, finite capacity
$C_\Gamma<\infty$, and witnesses $r_1,\ldots,r_N$ that are
support-independent in the sense of
Definition~\ref{def:support-independent-resolution}.  Then
\[
    N\;\le\;\left\lfloor \frac{C_\Gamma}{\mu_\Gamma}\right\rfloor.
\]
\end{theorem}
\begin{proof}
Each witness satisfies $\kappa_\Gamma(r_i)\ge\mu_\Gamma$ by the
distinction-floor hypothesis.  Support independence (iii) gives the
additive lower bound
\[
    \kappa_\Gamma(\{r_1,\ldots,r_N\})
    \;\ge\;
    \sum_{i=1}^N\kappa_\Gamma(r_i)
    \;\ge\;
    N\,\mu_\Gamma.
\]
Joint enforceability requires
\[
    \kappa_\Gamma(\{r_1,\ldots,r_N\})\;\le\;C_\Gamma.
\]
Therefore $N\,\mu_\Gamma\le C_\Gamma$, hence
$N\le\lfloor C_\Gamma/\mu_\Gamma\rfloor$.  Conditions (i), (ii), and
(iv) of Definition~\ref{def:support-independent-resolution} ensure
that no single coarser record subsumes the witnesses; if any of these
conditions fails, the additive lower bound fails and the theorem is
unavailable.
\end{proof}

\begin{corollary}[Counting bound from the marginal-floor structural inequality]
\label{cor:counting-bound-from-marginal-floor}
The same counting bound $N\le\lfloor C_\Gamma/\varepsilon^*_\Gamma\rfloor$ holds for any admissible family of $N$ tested distinctions in any finite physical regime, without requiring support-independence (Definition~\ref{def:support-independent-resolution}).
\end{corollary}
\begin{proof}
By Corollary~\ref{cor:sum-of-floors-lower-bound}, $\kappa_\Gamma(S)\ge N\varepsilon^*_\Gamma$ for any admissible $S$ with $|S|=N$.  Combined with the ledger-admissibility constraint $\kappa_\Gamma(S)\le C_\Gamma$ and integrality, $N\le\lfloor C_\Gamma/\varepsilon^*_\Gamma\rfloor$.  Theorem~\ref{thm:independent-counting-bound-formal} above is the special case where support-independence licenses the additive lower bound by an alternative route; the present corollary shows that the additive lower bound is not the only route to the counting conclusion --- the marginal-floor structural inequality (Lemma~\ref{lem:marginal-floor-on-joint-cost}) suffices.
\end{proof}

\begin{remark}[Positive interface cost is not IJC]
A positive interface term $\kappa_{\mathrm{int}}>0$ does not imply IJC.  Classical systems with shared baths, shared control buses, or shared memory may have positive interface overhead while still admitting a common Boolean record defender.  IJC requires failure of all faithful commuting continuation completions, not merely nonzero coupling cost.
\end{remark}

\section{Continuation words and contextual quotient algebra}\label{sec:continuation-word-algebra}

This section supplies the algebraic construction used by the Sep/IJC theorem.  The construction is a forced bookkeeping device for finite sequential query behavior; it does not assume quantum operators, Hilbert space, or physical superposition.  It starts from finite process fragments.

\begin{definition}[Query continuations]
Let $Q=\{d_1,\ldots,d_n\}\subset\mathcal D_\Gamma$ be a finite family of queried distinctions.  Let $E_i$ denote an admissible finite process that queries, enforces, or records $d_i$ in the context where it is applied.
\end{definition}

\begin{definition}[Continuation word]
A continuation word is a finite formal product
\[
 W=E_{i_k}\cdots E_{i_1}
\]
interpreted operationally as: apply $E_{i_1}$, then $E_{i_2}$, and so on through $E_{i_k}$, wherever the sequential process remains admissible.
\end{definition}

\begin{definition}[Typed process fragments]\label{def:typed-process-fragments}
Each finite query fragment is typed by its input and output contexts,
\[
  E_i:\Gamma_{\mathrm{in}}(E_i)\leadsto \Gamma_{\mathrm{out}}(E_i),
\]
and by the finite ledger support required to execute it.  A fragment is admissible only relative to these data.  Thus the symbol $E_i$ is not a context-free operator; it is a finite continuation move with a domain, codomain, and cost support.
\end{definition}


\begin{definition}[Finite APF process category]\label{def:finite-apf-process-category}
For a fixed family of admissible finite-resolution contexts, the finite APF process category \(\mathsf{Cont}_\Gamma\) is the following typed partial process category.
\begin{enumerate}[label=(\roman*),leftmargin=2em]
\item Objects are admissible finite-resolution contexts \(\Gamma_a,\Gamma_b,\ldots\).
\item Arrows \(K:\Gamma_a\leadsto\Gamma_b\) are finite admissible continuation fragments carrying specified ledger support.
\item Each object has an identity continuation \(\mathrm{id}_{\Gamma_a}\).
\item A composite \(L\circ K\) is defined exactly when the output context of \(K\) matches the input context of \(L\), all required recovery or coarse-graining maps are specified, and the combined ledger load remains admissible.
\item Whenever both composites are defined, associativity holds:
\[
  (M\circ L)\circ K=M\circ(L\circ K).
\]
\item The tested subcategory \(\mathsf{Cont}_{\Gamma,Q,\mathcal T}\) contains only the finite words actually tested, together with the tested products, identity words, recovery maps, abort/failure words, and coarse-grainings required by \(\mathcal T\).
\end{enumerate}
Thus the continuation-word algebra used below is the real path algebra of a finite typed partial process category, quotiented by contextual continuation equivalence.  Incompatible concatenations are not zero morphisms; they are absent from the composition domain unless explicitly tested as abort or failure arrows.
\end{definition}

\begin{definition}[Partial word product]\label{def:partial-word-product}
For two typed words $U$ and $V$, the product $UV$ means ``perform $V$ first and then $U$.''  It is defined exactly when
\begin{enumerate}[label=(\roman*),leftmargin=2em]
\item the output context of $V$ matches the input context of $U$ after the stated recovery/coarse-graining maps;
\item the composed process remains admissible in the context under discussion; and
\item the composed word is part of the tested family, or has been explicitly added to the tested family at the chosen depth and resolution.
\end{enumerate}
An incompatible concatenation is not silently assigned the value zero.  It is outside the product domain unless the protocol explicitly tests an abort/failure outcome, in which case that abort is a charged tested word.  Order effects are different: if both $UV$ and $VU$ are defined but have different continuation profiles, they are distinct word classes rather than undefined products.
\end{definition}

\begin{lemma}[Finite continuations form a partial process category]
Finite admissible process fragments form a partial category: identity continuations exist, and sequential composition $L\circ K$ is defined whenever the output context of $K$ matches the input context of $L$ and the combined process remains admissible.
\end{lemma}
\begin{proof}
The identity continuation does nothing and preserves the continuation profile.  If $K:x\leadsto y$ and $L:y\leadsto z$ are finite admissible fragments and the combined enforcement cost remains within the context budget, then performing $K$ followed by $L$ is again a finite admissible process from $x$ to $z$.  Associativity holds wherever all composites are defined, because concatenating finite process fragments is associative.
\end{proof}

\begin{definition}[Contextual continuation equivalence of words]
For continuation words $W,W'$, define
\[
 W\equiv_{\mathrm{cont}}W'
\]
iff for every admissible prefix $P$, suffix $S$, and compatible input possibility $[x]$, the processed continuations $SWP[x]$ and $SW'P[x]$ have the same physical continuation profile whenever both sides are defined.
\end{definition}

\begin{lemma}[Contextual equivalence is a congruence]
If $W\equiv_{\mathrm{cont}}W'$, then for all admissible words $U,V$,
\[
 U W V\equiv_{\mathrm{cont}} U W' V
\]
where the composites are defined.
\end{lemma}
\begin{proof}
The definition of $\equiv_{\mathrm{cont}}$ quantifies over all admissible prefixes and suffixes.  Taking the prefix to include $V$ and the suffix to include $U$ gives the result.
\end{proof}

\begin{definition}[Contextual null space]
Let $\mathbb R[\mathcal W_Q]$ be the real path algebra of the finite typed partial process category generated by the query symbols.  Its basis elements are admissible finite words; multiplication is typed concatenation when defined and is used only on product-closed tested subfamilies or explicitly tested products.  Define
\[
 \mathcal N_{\mathrm{cont}}
 :=\mathrm{span}\{W-W': W\equiv_{\mathrm{cont}}W'\}.
\]
\end{definition}


\begin{remark}[Free-word algebra reading]\label{rem:free-word-algebra-reading}
Equivalently, if one temporarily forgets typing and admissibility, this is the free word algebra on the query symbols modulo the relations generated by contextual continuation equivalence.  The typed partial-category formulation is used because APF must distinguish undefined concatenations from zero-valued processes, and must preserve order effects when both \(UV\) and \(VU\) are defined but operationally different.
\end{remark}

\begin{proposition}[The contextual null space is a two-sided ideal]\label{prop:null-space-two-sided-ideal}
$\mathcal N_{\mathrm{cont}}$ is a two-sided ideal of the defined product algebra $\mathbb R[\mathcal W_Q]$; equivalently, for every defined left and right multiplication by tested words, contextual null differences remain contextual null differences.
\end{proposition}
\begin{proof}
It suffices to check generators of $\mathcal N_{\mathrm{cont}}$.  If $W-W'$ is a generator and $U,V$ are words for which the products are defined in the typed partial word algebra, then by contextual congruence $UWV\equiv_{\mathrm{cont}}UW'V$, so
\[
 U(W-W')V=UWV-UW'V\in\mathcal N_{\mathrm{cont}}.
\]
Linearity gives closure under all defined algebra multiplications on the left and right.  Incompatible concatenations are not algebra multiplications and therefore impose no ideal condition unless an explicit failure word has been added to the tested algebra.
\end{proof}

\begin{definition}[Continuation-word algebra]
The continuation-word algebra is
\[
 \mathcal A_Q
 :=\mathbb R[\mathcal W_Q] / \mathcal N_{\mathrm{cont}}.
\]
By Proposition~\ref{prop:null-space-two-sided-ideal}, this quotient is a well-defined associative algebra whose multiplication is induced by sequential concatenation of finite continuation words.
\end{definition}

\begin{definition}[Finite tested word set]\label{def:finite-tested-word-set}
For a finite joint-query interface, a \emph{tested word set} $\mathcal T$ is a finite set of continuation-word classes containing the identity word, the query generators, and every recovery, coarse-graining, or sequential word whose behavior the defender is required to reproduce.  We assume $\mathcal T$ is closed under taking tested subwords.  Products of tested words are used only when the resulting word is itself included in $\mathcal T$ or when the product is explicitly added to the tested set.
\end{definition}

\begin{definition}[Finite tested continuation quotient]\label{def:finite-tested-continuation-algebra}
Given a tested word set $\mathcal T$, define
\[
 \AQtest := \mathrm{span}_{\R}(\mathcal T)/\bigl(\mathcal N_{\mathrm{cont}}\cap \mathrm{span}_{\R}(\mathcal T)\bigr).
\]
This is the finite-dimensional tested quotient used in the Sep/IJC representation theorem.  It is a partial algebraic object: products are required only for tested concatenations whose product word is included in $\mathcal T$.  If a full finite algebra is desired, one first enlarges $\mathcal T$ to a finite product-closed test suite at the chosen depth and resolution.  The full formal word algebra $\mathcal A_Q$ remains useful for bookkeeping, but branch classification is always applied to a specified finite tested/resolved quotient.  When the tested set is clear from context, we write $\mathcal A_Q$ informally for this finite tested quotient.
\end{definition}

\begin{theorem}[Existence and universal property of the tested continuation quotient]\label{thm:universal-property-tested-quotient}
Let $\Cont_{\Gamma,Q,\mathcal T}$ be the finite typed partial process
category generated by the tested query fragments, tested composites,
identities, abort/failure arrows, recovery maps, and coarse-grainings
included in $\mathcal T$.  Let
$\R[\Cont_{\Gamma,Q,\mathcal T}]$ be the real tested path algebra whose
basis elements are tested words and whose multiplication is defined
exactly on tested compatible composites.

Let $\equiv_{\Gamma,\delta}$ be context-complete continuation
equivalence at operational resolution $\delta$, and let
\[
    \mathcal N_{\Gamma,\delta}
    \;=\;
    \mathrm{span}\{\,U-V \;:\; U\equiv_{\Gamma,\delta}V\,\}.
\]
Then $\mathcal N_{\Gamma,\delta}$ is a two-sided ideal for all defined
tested products, and the quotient
\[
    A^{\Gamma,\delta}_{Q,\mathcal T}
    \;=\;
    \R[\Cont_{\Gamma,Q,\mathcal T}]
    \big/\mathcal N_{\Gamma,\delta}
\]
exists.  It is unique up to unique isomorphism among real tested-process
algebras in which contextually equivalent tested words are identified.

More precisely, if $B$ is any real tested-process algebra and
\[
    F:\R[\Cont_{\Gamma,Q,\mathcal T}]\to B
\]
is a tested-product-preserving linear map satisfying
\[
    U\equiv_{\Gamma,\delta}V \;\Longrightarrow\; F(U)=F(V),
\]
then there exists a unique homomorphism
\[
    \bar F:A^{\Gamma,\delta}_{Q,\mathcal T}\to B
\]
such that $F=\bar F\circ\pi$, where $\pi$ is the quotient map.

The object $A^{\Gamma,\delta}_{Q,\mathcal T}$ is the resolution-explicit
form of the finite tested continuation quotient $\AQtest$ of
Definition~\ref{def:finite-tested-continuation-algebra}; the supplement's
working notation suppresses $\Gamma$ and $\delta$ when they are clear
from context.
\end{theorem}

\begin{proof}
Context-complete continuation equivalence is a congruence: if
$U\equiv_{\Gamma,\delta}V$, then for every tested prefix $P$ and suffix
$S$ for which the composites are defined,
\[
    SUP \;\equiv_{\Gamma,\delta}\; SVP.
\]
Therefore every generator $U-V$ of $\mathcal N_{\Gamma,\delta}$ remains
inside $\mathcal N_{\Gamma,\delta}$ after multiplication on either side
by any tested compatible word.  Hence $\mathcal N_{\Gamma,\delta}$ is a
two-sided ideal for the tested product structure.

The quotient algebra therefore exists.  If $F$ identifies all
contextually equivalent words, then
$\mathcal N_{\Gamma,\delta}\subseteq\ker F$.  The standard quotient
property gives a unique homomorphism $\bar F$ with
$F=\bar F\circ\pi$.  Uniqueness follows from surjectivity of $\pi$.
Any two objects satisfying this property factor uniquely through one
another and are therefore uniquely isomorphic.
\end{proof}


\begin{definition}[Tested state and effect functionals]\label{def:tested-state-effect-functionals}
For a finite tested quotient $\AQtest$, a tested preparation or input condition $s$ determines a linear functional
\[
  \omega_s:\AQtest\to\mathbb R,
\]
where $\omega_s([W])$ is the operational value, probability, expectation, or resolved readout statistic assigned to the tested word $W$ from $s$.  A tested effect, discriminator, or readout event $f$ determines a linear functional on the represented output records, written equivalently as an evaluation map on tested words.  The unit word $1$ is normalized by $\omega_s(1)=1$ for normalized preparations.
\end{definition}

\begin{definition}[Order and positivity on the tested quotient]\label{def:tested-order-positivity}
The tested positive cone $\AQtest^+$ is the cone generated by tested events and finite positive combinations of tested events whose operational values are nonnegative on all tested preparations.  A representation of $\AQtest$ is order-preserving on the tested interface when
\[
  a\in\AQtest^+ \quad\Longrightarrow\quad \rho(a)(\lambda)\geq 0
  \quad\text{for every Boolean atom }\lambda.
\]
It is unit-preserving when $\rho(1)=1_\Lambda$.  These requirements apply only to the finite order/unit structure actually resolved by $\mathcal T$; no untested no-restriction hypothesis or global GPT cone is assumed.
\end{definition}

\begin{remark}[How states and effects are represented]\label{rem:states-effects-represented}
Under a faithful commutative representation $\rho:\AQtest\to\mathbb R^\Lambda$, normalized tested states become probability weights $p_s\in\Delta(\Lambda)$ whenever the tested statistics are probabilistic:
\[
  \omega_s(a)=\sum_{\lambda\in\Lambda}p_s(\lambda)\rho(a)(\lambda),
  \qquad a\in\AQtest .
\]
Tested effects become functions $e:\Lambda\to[0,1]$ when the readouts are probabilistic events.  Thus the defender is not merely an algebra homomorphism: it is a finite Boolean-record model preserving the tested unit, positivity/order, normalization, and all state--effect pairings included in the finite interface.
\end{remark}

\begin{remark}[Formal linearization versus operational mixtures]\label{rem:linearization-not-mixture}
The generators and words of $\mathcal A_Q$ represent finite admissible process fragments.  The real linear span is not an assertion that arbitrary real linear combinations of processes are themselves physical operations.  Rather, it is the standard algebraic linearization of a process monoid modulo contextual continuation equivalence, used to form quotients and representation maps.  Physical probabilistic mixtures enter only through admissible classical control registers, as in the convexity theorem.  Thus the quotient algebra is a bookkeeping and representation object for finite word behavior, not an additional postulate of operational superposition.
\end{remark}

\begin{example}[Formal linearization is not operational mixing]\label{ex:linearization-vs-mixture}
\textbf{Compatible binary readout.}
Let \(E_0\) and \(E_1\) denote two tested process words.  The formal element \(E_0-E_1\in\AQtest\) is meaningful as an algebraic probe: a tested state functional may evaluate \(\omega(E_0-E_1)=\omega(E_0)-\omega(E_1)\), so it can witness that the two words have different continuation behavior.  But \(E_0-E_1\) is not itself a physical operation, because negative execution frequency is not an admissible control protocol.

By contrast, \(pE_0+(1-p)E_1\) is operational only when there exists an admissible classical control bit that selects \(E_0\) with frequency \(p\) and \(E_1\) with frequency \(1-p\), and the control record fits in the ledger.  Thus linearization supplies a representation calculus; operational mixtures require a charged finite control mechanism.
\end{example}
Only distinctions that require independent records consume independent capacity.  The next definitions make this restriction explicit so the finite-resolution bound does not count merely formal algebraic differences.

\begin{definition}[Resolved continuation-word classes]\label{def:operationally-resolved-words}
A finite family of continuation-word classes
\[
 \mathcal W=\{[W_1],\ldots,[W_N]\}\subset\mathcal A_Q
\]
is \emph{operationally resolved in context $\Gamma$} when, for each $i$, there exists a finite admissible discriminator $r_i\in\mathcal D_\Gamma$ whose outcome separates the continuation profile generated by $W_i$ from the competing classes at the stated operational resolution.
\end{definition}

\begin{definition}[Ledger-independent resolution]\label{def:ledger-independent-resolution}
An operationally resolved family
\[
 \mathcal W=\{[W_1],\ldots,[W_N]\}
\]
is \emph{ledger-independent} when its witnessing discriminators cannot be implemented by a single shared coarser record.

Explicitly, for witnesses $r_1,\ldots,r_N$, ledger independence requires
\[
 \kappa_\Gamma(\{r_1,\ldots,r_N\})
 \ge \sum_{i=1}^N \kappa_\Gamma(r_i).
\]
Equivalently, any shared implementation must still maintain $N$ independent separation degrees of freedom at the specified resolution.
\end{definition}

\begin{remark}[Algebraic difference is not automatically capacity cost]\label{rem:algebraic-vs-ledger-independence}
Two algebraically distinct word classes need not consume additive ledger capacity.  They may be separated by the same record, differ only below the chosen resolution, or be distinguished only by a coupled discriminator with nonadditive cost.  The finite-resolution bound below counts only ledger-independent operationally resolved classes, not every formal word in the free algebra.
\end{remark}


\begin{remark}[Practical criteria for ledger-independent resolution]\label{rem:ledger-independence-practice}
Ledger independence is not meant to double-count algebraic labels.  In a concrete finite model, a family of discriminators is ledger-independent only when all three checks hold at the stated resolution:
\begin{enumerate}[label=(\roman*),leftmargin=2em]
\item \emph{No shared coarser record:} there is no single record variable whose finite bins implement all of the discriminators without maintaining their individual separators.
\item \emph{Removal sensitivity:} removing the record support for any one discriminator changes at least one resolved continuation profile that remains fixed when the other discriminator supports are retained.
\item \emph{No free recovery by control:} classical control, filtering, or relabeling cannot recover the missing discriminator without introducing a new charged record or stabilizer.
\end{enumerate}
Borderline cases should be treated conservatively: if two discriminators share an implementation through a common physical record, the shared implementation is charged once and the claimed independence must be justified by an explicit support decomposition.
\end{remark}


\begin{remark}[Bookkeeping prescription for shared records]\label{rem:shared-record-bookkeeping}
To avoid double-counting, use the following ledger prescription for any finite tested family.
\begin{enumerate}[label=(\roman*),leftmargin=2em]
\item List the actual maintained physical records, control registers, stabilizers, and recovery memories used by the protocol.
\item Map each formal distinction or word-class discriminator to the maintained records that implement it at the stated resolution.
\item Merge all formal distinctions supported by the same maintained record before applying an independent-counting bound.
\item Charge the least admissible joint protocol, including any interface term \(\kappa_{\rm int}\), rather than the formal sum of names.
\item Claim ledger independence only for the remaining supports whose removal changes distinct resolved continuation profiles and cannot be recovered by uncharged relabeling or control.
\end{enumerate}
We keep the prescription conservative.  It lets one use the algebra to find candidate distinctions without letting the algebra inflate the capacity count.
\end{remark}

\begin{example}[Shared coarse record: not ledger-independent]\label{ex:shared-coarse-record}
Suppose a finite record register $R\in\{0,1,2,3\}$ is already maintained at the interface, and two discriminators read the first and second binary digit of $R$.  Algebraically these are two binary distinctions, but their enforcement support is one four-valued maintained record.  They are not counted as two ledger-independent burdens unless the implementation actually maintains two independent separators.  The correct charge is the least joint cost $\kappa_\Gamma(\{r_1,r_2\})$, not the formal sum of two labels.
\end{example}

\begin{example}[Independent records: ledger-independent]\label{ex:independent-records}
Suppose two discriminators require two separately maintained physical records $R_A$ and $R_B$, and removing the support for $R_A$ changes an $A$-resolved continuation profile while leaving the $B$ profile intact, and conversely for $R_B$.  If no admissible relabeling or shared coarse register recovers both distinctions at the stated resolution, then the family satisfies the ledger-independence checks and the counting bound may charge both burdens.
\end{example}

\begin{example}[Coupled discriminator: count conservatively]\label{ex:coupled-discriminator}
If two word classes are separated only by a coupled discriminator that probes their joint behavior, then the pair is operationally resolved but not automatically ledger-independent.  The coupled discriminator may be a single interface record with a nonadditive cost.  In that case the finite-resolution theorem applies to the independent resolved rank after decomposing, or failing to decompose, the support into independent maintained separators.
\end{example}

\begin{theorem}[Finite-resolution bound]\label{thm:finite-resolution-bound}
In a robust finite interface with finite query family $Q$, distinction floor $\mu_\Gamma(Q)>0$, and total capacity $C_\Gamma<\infty$, the number $N$ of jointly enforceable ledger-independent operationally resolved continuation-word classes satisfies
\[
 N\le \left\lfloor \frac{C_\Gamma}{\mu_\Gamma(Q)}\right\rfloor .
\]
Consequently, the ledger-independent resolved rank of the queried word behavior at fixed context resolution is finite.  If, in addition, the tested word set is finite, the corresponding tested continuation algebra has finite dimension.
\end{theorem}
\begin{proof}
Let $\mathcal W=\{[W_1],\ldots,[W_N]\}$ be jointly enforceable and ledger-independent, with witnessing discriminators $r_1,\ldots,r_N$.  Operational resolution gives $r_i\in\mathcal D_\Gamma$ for each class.  By the distinction-floor theorem, $\kappa_\Gamma(r_i)\ge\mu_\Gamma(Q)$ at the stated resolution.  Ledger independence gives
\[
 \kappa_\Gamma(\{r_1,\ldots,r_N\})
 \ge \sum_{i=1}^N\kappa_\Gamma(r_i)
 \ge N\mu_\Gamma(Q).
\]
Joint enforceability requires $\kappa_\Gamma(\{r_1,\ldots,r_N\})\le C_\Gamma$.  Hence $N\mu_\Gamma(Q)\le C_\Gamma$, proving the bound.  At fixed operational resolution, unresolved formal word differences are quotiented away and no ledger-independent resolved subfamily can have cardinality exceeding this bound.  For a finite tested word set, the tested quotient is a finite span of tested classes modulo continuation equivalence, hence finite-dimensional.
\end{proof}

\section{Sep/IJC representation theorem}\label{sec:sep-ijc-representation}

\subsection{Main theorem objects collected}\label{subsec:main-theorem-objects}
For reference, the finite representation theorem uses the following objects and only these objects.
\begin{enumerate}[label=(\roman*),leftmargin=2em]
\item A finite query family $Q={d_1,\ldots,d_n}$ in context $\Gamma$.
\item A finite tested word set $\mathcal T$, closed under the products, identities, and recovery/coarse-graining operations that the interface actually tests at the stated resolution.
\item The finite tested continuation quotient
\[
  \mathcal A_{Q,\mathcal T},
\]
whose elements are tested continuation-word classes modulo context-complete continuation equivalence.
\item A commuting-extension defender: a finite reduced Boolean record space reproducing the tested word behavior.
\item A faithful commutative continuation representation
\[
 \rho:\mathcal A_{Q,\mathcal T}\hookrightarrow\mathbb R^\Lambda,
\]
which is injective on resolved tested classes and preserves the tested products, identities, recovery maps, and continuation values included in $\mathcal T$.
\end{enumerate}
The theorem is therefore a finite tested-interface theorem.  It does not classify untested words, infinite-depth completions, or ideal continuum protocols unless they are supplied by an additional limiting construction.

This is the classification layer.  The theorem below does not add a new physical axiom; it identifies the exact condition under which the finite continuation-word behavior of a queried interface can be defended as one shared classical record.

\begin{definition}[Finite joint-query interface]
A finite joint-query interface is a pair $(\Gamma,Q)$ where $\Gamma$ is an admissible context and $Q=\{d_1,\ldots,d_n\}\subset\mathcal D_\Gamma$ is a finite family of individually enforceable queried distinctions.
\end{definition}

\begin{definition}[Commuting-extension defender]\label{def:commuting-extension-defender}
Fix a finite tested word set $\mathcal T$ for $(\Gamma,Q)$.  A \emph{commuting-extension defender} for $(\Gamma,Q,\mathcal T)$ consists of:
\begin{enumerate}[label=(\roman*),leftmargin=2em]
\item a finite admissible record space $\widetilde\Omega$ and reduced Boolean algebra $\mathcal B=2^{\widetilde\Omega}$;
\item a unit-preserving, order-preserving tested representation of the queried events by Boolean events $\widetilde d_i\in\mathcal B$;
\item for every tested word $W\in\mathcal T$, a Boolean-record representative $\widetilde W$ obtained by the defender's tested product/recovery/coarse-graining rules;
\item a recovery/coarse-graining map $r$ from defender records back to the operational readout alphabet such that every tested state--effect pairing and every tested continuation value is reproduced;
\item no hidden uncharged resource: any added memory, order register, ancilla, stabilizer, or control variable used by the defender is part of $\widetilde\Omega$ and is therefore either included in the defended record algebra or explicitly charged as part of the finite interface.
\end{enumerate}
The defender is \emph{faithful} when
\[
  \widetilde U=\widetilde V
  \quad\Longleftrightarrow\quad
  U\equiv_{\mathrm{cont}}V
  \quad\text{inside }\AQtest
\]
for all tested words $U,V\in\mathcal T$, and when it preserves the tested unit, tested products, tested order/positivity, normalization, recovery maps, and all tested state--effect pairings.  Thus faithfulness is injectivity on the resolved finite tested quotient, not faithfulness to untested words or to an ideal infinite completion.
\end{definition}

\begin{definition}[Sep and IJC]\label{def:sep-ijc-tested}
A finite joint-query interface $(\Gamma,Q)$, together with a specified finite tested word set $\mathcal T$, is in the Sep branch when a commuting-extension defender exists for that tested word behavior.  It is in the IJC branch when no such defender exists.  When $\mathcal T$ is fixed by context, we suppress it from the notation.
\end{definition}

\begin{definition}[Boolean record algebra and its real linearization]\label{def:boolean-record-algebra}
Let $\Lambda$ be a finite record-atom space.  The associated Boolean
record algebra is
\[
    \mathcal B_\Lambda \;:=\; 2^\Lambda.
\]
Since $\Lambda$ is finite, $\mathcal B_\Lambda$ is also a Boolean
$\sigma$-algebra.

For algebraic representation it is useful to pass to the real
linearization
\[
    \R[\mathcal B_\Lambda] \;\cong\; \R^\Lambda,
\]
the algebra of real-valued functions on $\Lambda$ with pointwise
addition and multiplication.  A Boolean event $E\subseteq\Lambda$ is
identified with its indicator function
$1_E\in\R^\Lambda$.
\end{definition}

\begin{definition}[Commutative record representation]\label{def:finite-commutative-record-model}
Fix a finite tested interface $(\Gamma,Q,\mathcal T,\delta)$ and its
resolved tested continuation quotient
$A^{\Gamma,\delta}_{Q,\mathcal T}$.  A \emph{commutative record
representation} consists of a finite atom space $\Lambda$ and a
unit-preserving tested-product homomorphism
\[
    \rho:A^{\Gamma,\delta}_{Q,\mathcal T}
    \longrightarrow \R^\Lambda
\]
such that, whenever $U,V,UV\in\mathcal T$ are tested compatible words,
\[
    \rho(UV)\;=\;\rho(U)\,\rho(V),
\]
with multiplication on the right taken pointwise.  The representation
also preserves the tested recovery maps, coarse-grainings,
normalization, and tested positivity/order structure included in the
finite interface.

For a tested event $e$, $\rho(e)$ is required to be an indicator
function $1_{E_e}$ for some Boolean record event
$E_e\in\mathcal B_\Lambda$.  More general tested word values are
represented as real-valued functions on the same atom space.  No
untested continuum event structure is assumed.
\end{definition}

\begin{definition}[$\delta$-faithfulness]\label{def:delta-faithfulness}
A finite commutative record model
\[
    \rho:A^{\Gamma,\delta}_{Q,\mathcal T}\to\R^\Lambda
\]
is \emph{$\delta$-faithful} when its kernel is exactly the contextual
null ideal:
\[
    \ker\rho \;=\; 0
    \quad\text{on}\quad
    A^{\Gamma,\delta}_{Q,\mathcal T}.
\]
Equivalently, for tested words $U,V\in\mathcal T$,
\[
    \rho([U])=\rho([V])
    \;\Longleftrightarrow\;
    U\equiv_{\Gamma,\delta}V.
\]
Thus faithfulness is not faithfulness to raw labels or to untested
counterfactual words.  It is injectivity on the resolved tested quotient
after all distinctions below operational resolution $\delta$ have
already been identified.  In this sense the representation is an
embedding of the finite resolved quotient, not an embedding of the raw
process syntax.
\end{definition}

\paragraph{Faithfulness is resolution-relative.}
Faithfulness is always stated after quotienting by the operational
resolution $\delta$.  Thus a faithful defender need not distinguish raw
labels, syntactic word names, untested counterfactual continuations, or
differences lying below the resolved threshold.  It must distinguish
exactly those tested continuation-word classes that remain distinct in
\[
    A^{\Gamma,\delta}_{Q,\mathcal T}.
\]
Equivalently, faithfulness means that the defender introduces no extra
identifications beyond contextual continuation equivalence and no extra
separations of merely raw notation.  The representation is therefore an
embedding of the finite resolved tested quotient, not an embedding of
the unquotiented raw process syntax.

\begin{definition}[Two-valued record homomorphisms]\label{def:two-valued-homomorphisms}
For each atom $\lambda\in\Lambda$, define the two-valued Boolean
homomorphism
\[
    \chi_\lambda:\mathcal B_\Lambda\to\{0,1\}
\]
by $\chi_\lambda(E)=1_E(\lambda)$.  Equivalently,
$\chi_\lambda(E)=1$ exactly when $\lambda\in E$.  These are precisely
the Boolean characters of the finite record algebra.

The corresponding evaluation functional on the real linearization
$\R^\Lambda$ is
\[
    \mathrm{ev}_\lambda(f) \;=\; f(\lambda).
\]
Composing with a commutative record representation gives a tested-word
valuation
\[
    v_\lambda \;:=\; \mathrm{ev}_\lambda\circ\rho
    : A^{\Gamma,\delta}_{Q,\mathcal T}\to\R.
\]
On tested Boolean events, $v_\lambda$ is two-valued.
\end{definition}

\begin{definition}[Defender state simplex]\label{def:defender-state-simplex}
Given a commutative record representation
$\rho:A^{\Gamma,\delta}_{Q,\mathcal T}\to\R^\Lambda$, the defender
state space is the classical simplex
\[
    \Delta(\Lambda)
    \;=\;
    \Bigl\{\, p:\Lambda\to[0,1] \;:\;
    \sum_{\lambda\in\Lambda}p(\lambda)=1\,\Bigr\}.
\]
Each $p\in\Delta(\Lambda)$ defines a state on the tested quotient by
\[
    \omega_p(a)
    \;=\;
    \sum_{\lambda\in\Lambda}p(\lambda)\,\rho(a)(\lambda),
    \qquad
    a\in A^{\Gamma,\delta}_{Q,\mathcal T}.
\]
Equivalently,
\[
    \omega_p \;=\;
    \sum_{\lambda\in\Lambda}p(\lambda)\,v_\lambda.
\]
Thus the defender states are exactly the convex hull of the
deterministic two-valued record valuations, pulled back along $\rho$.
\end{definition}

\begin{definition}[Tested word behavior]\label{def:tested-word-behavior}
The \emph{tested word behavior} of $(\Gamma,Q,\mathcal T,\delta)$ is
the finite operational table
\[
    \mathsf B_{\Gamma,Q,\mathcal T,\delta}
    \;=\;
    \bigl\{\,\omega_s(a)\;:\; s\in\mathcal S_{\mathrm{test}},\;
    a\in A^{\Gamma,\delta}_{Q,\mathcal T}\,\bigr\},
\]
where $s$ ranges over tested preparations or input conditions, and
$\omega_s(a)$ is the observed probability, expectation value,
pass/fail frequency, or declared operational value associated with the
tested word, event, or recovery outcome $a$.

The algebra $A^{\Gamma,\delta}_{Q,\mathcal T}$ records the tested
continuation syntax modulo operational equivalence; the state
functionals $\omega_s$ record the finite probabilistic behavior on
that quotient.  $\mathsf B_{\Gamma,Q,\mathcal T,\delta}$ is therefore
the stochastic-data layer of the finite tested interface.  A faithful
commutative defender is required to reproduce
$\mathsf B_{\Gamma,Q,\mathcal T,\delta}$ as a convex mixture of
deterministic two-valued atom valuations on
$A^{\Gamma,\delta}_{Q,\mathcal T}$.
\end{definition}

\begin{lemma}[Finite atom dilation of a commutative representation]\label{lem:finite-atom-dilation}
Let
\[
    \rho:A^{\Gamma,\delta}_{Q,\mathcal T}\to\R^\Lambda
\]
be a finite commutative continuation representation.  Let
\[
    q_\rho:\Lambda\to\prod_{W\in\mathcal T}\mathrm{im}\,\rho(W)
\]
be the joint tested evaluation map
\[
    q_\rho(\lambda) \;=\; \bigl(\rho(W)(\lambda)\bigr)_{W\in\mathcal T}.
\]
Then $\rho$ factors through the finite Boolean algebra generated by the
nonempty fibers of $q_\rho$.  This reduced atom space is minimal among
record spaces reproducing the same tested continuation table.
\end{lemma}
\begin{proof}
Two atoms $\lambda,\lambda'$ with the same joint tested evaluation
profile cannot be distinguished by any tested word, recovery map, or
state--effect value in $\mathcal T$.  Identifying such atoms therefore
preserves all tested data.  Conversely, if two atoms have different
joint tested profiles, some tested coordinate separates them, so any
faithful record model reproducing the tested table must keep that
distinction available.  The quotient by equal tested profiles is
therefore the minimal finite record refinement through which $\rho$
factors.
\end{proof}

\begin{theorem}[Boolean feasibility and separating witnesses]\label{thm:boolean-feasibility-separating-witnesses}
For a fixed finite tested interface $(\Gamma,Q,\mathcal T,\delta)$, the
existence of a $\delta$-faithful commutative record model is equivalent
to feasibility of a finite Boolean-record constraint system.

The variables of this system are:
\begin{enumerate}[label=(\roman*),leftmargin=2em]
\item finitely many record atoms $\lambda\in\Lambda$;
\item values $\rho(W)(\lambda)$ for every tested word
$W\in\mathcal T$;
\item state weights $p_s(\lambda)$ for every tested preparation $s$;
\item recovery/readout coordinates for the tested recovery maps.
\end{enumerate}
The constraints are:
\begin{enumerate}[label=(\arabic*),leftmargin=2em]
\item unit preservation: $\rho(1)=1_\Lambda$;
\item tested product preservation:
$\rho(UV)=\rho(U)\,\rho(V)$ whenever $U,V,UV\in\mathcal T$;
\item tested recovery and coarse-graining preservation;
\item positivity/order constraints for tested positive events;
\item normalization and state--effect reproduction;
\item separation constraints:
\[
    U\not\equiv_{\Gamma,\delta}V
    \;\Longrightarrow\;
    \rho(U)\neq\rho(V)
\]
inside the resolved tested quotient.
\end{enumerate}
Consequently, IJC is not merely a terminological negation of Sep.  It
is the infeasibility of this finite Boolean-record constraint system.
When the constraint system is infeasible, there exists a finite
\emph{separating witness}: a linear functional $L$ on the tested table
such that
\[
    L(R)\le c
\]
for every Boolean record model $R$, while the observed tested table
satisfies
\[
    L(\Gamma,Q,\mathcal T,\delta) > c.
\]
\end{theorem}
\begin{proof}
Given a finite commutative record model, take its atom space $\Lambda$.
Evaluation of every tested word on every atom gives the variables
$\rho(W)(\lambda)$, and the model axioms give the listed constraints.
Thus every defender gives a feasible point.

Conversely, given a feasible point of the Boolean-record constraint
system, construct $\R^\Lambda$ with pointwise multiplication and define
$\rho(W)$ by its listed atom values.  The product constraints make
$\rho$ a tested-product morphism.  The recovery constraints supply the
required readout maps.  The state weights reproduce the tested
state--effect and continuation values.  The separation constraints make
the model $\delta$-faithful on the resolved quotient.  Hence a feasible
point gives a $\delta$-faithful commutative record model.

Because the tested interface is finite, the set of Boolean-record
tables is a finite union of finite-dimensional polytopes, and in the
usual fixed finite atom bound it is a polytope.  If the observed table
lies outside this feasible set, finite-dimensional convex separation
gives a linear functional separating it from all Boolean-record tables.
This functional is the APF analogue of a Bell, Boole, Fine, or
sheaf-contextuality witness.  The CHSH inequality is the standard
special case when the tested words are static pairwise contexts and
recovery maps are trivial.
\end{proof}

\begin{theorem}[Soundness and completeness of commutative defenders]\label{thm:soundness-completeness}
For a fixed finite tested interface $(\Gamma,Q,\mathcal T,\delta)$, the
following are equivalent.
\begin{enumerate}[label=(\arabic*),leftmargin=2em]
\item The interface is Sep: there exists a faithful commuting-extension
defender reproducing the tested continuation behavior.
\item There exist a finite atom space $\Lambda$, a $\delta$-faithful
commutative record representation
\[
    \rho:A^{\Gamma,\delta}_{Q,\mathcal T}
    \hookrightarrow \R^\Lambda,
\]
and, for each tested preparation $s$, a defender state
$p_s\in\Delta(\Lambda)$ such that every tested state--effect and
continuation value is reproduced:
\[
    \omega_s(a)
    \;=\;
    \sum_{\lambda\in\Lambda}p_s(\lambda)\,\rho(a)(\lambda)
\]
for every tested word, event, recovery outcome, or state--effect
value $a\in A^{\Gamma,\delta}_{Q,\mathcal T}$.
\item The finite tested operational table lies in the convex hull of
deterministic two-valued record valuations compatible with the tested
products, recovery maps, coarse-grainings, order constraints, and
contextual continuation equivalences.
\end{enumerate}
Consequently, $\IJC(\Gamma,Q,\mathcal T,\delta)$ holds exactly when
no such $\delta$-faithful commutative record representation exists.
\end{theorem}

\begin{proof}
\emph{(1) implies (2): soundness.}  Assume a faithful
commuting-extension defender exists.  By definition, it uses a finite
classical record space.  Let $\Lambda$ be the reduced set of record
atoms after identifying any atoms that agree on all tested readout,
recovery, and continuation coordinates.  Its Boolean record algebra is
$\mathcal B_\Lambda=2^\Lambda$, with real linearization $\R^\Lambda$.

Each tested continuation word $W\in\mathcal T$ has, under the defender,
a record representative.  Define
$\rho([W])\in\R^\Lambda$ to be the function assigning to each atom
$\lambda$ the value of that record representative on $\lambda$.
Because the defender preserves the tested products, identities,
recovery maps, coarse-grainings, and tested order structure, $\rho$ is
a unit-preserving tested-product homomorphism into $\R^\Lambda$.

Faithfulness of the defender says that two tested words have the same
record representative exactly when they are contextually equivalent at
resolution $\delta$.  Hence $\rho$ is injective on the resolved tested
quotient $A^{\Gamma,\delta}_{Q,\mathcal T}$.  For each tested
preparation $s$, the defender assigns probabilities to the finite
record atoms; these probabilities define $p_s\in\Delta(\Lambda)$.  The
defender's reproduction condition gives
$\omega_s(a) = \sum_\lambda p_s(\lambda)\rho(a)(\lambda)$ for every
tested $a$.  Thus a $\delta$-faithful commutative record representation
exists.

\emph{(2) implies (1): completeness.}  Assume a $\delta$-faithful
commutative record representation
$\rho:A^{\Gamma,\delta}_{Q,\mathcal T}\hookrightarrow\R^\Lambda$ and
state weights $p_s\in\Delta(\Lambda)$ reproducing all tested values.
Construct a defender whose hidden record is the finite atom
$\lambda\in\Lambda$.  Tested Boolean events are represented by the
Boolean subsets $E\subseteq\Lambda$ whose indicators are the
corresponding $\rho(e)$.  Tested non-Boolean word values are
represented by their functions $\rho(W)\in\R^\Lambda$.  The declared
recovery and coarse-graining maps are those preserved by $\rho$.

Because multiplication in $\R^\Lambda$ is pointwise, all tested record
representatives commute.  Because $\rho$ preserves every tested
product, identity, recovery map, order relation, normalization, and
state--effect pairing, the constructed record model reproduces the
entire tested continuation table.  Because $\rho$ is $\delta$-faithful,
this defender identifies exactly the contextually equivalent tested
word classes and no additional resolved distinctions.  Therefore it is
a faithful commuting-extension defender.

\emph{Equivalence with (3).}  The atoms $\lambda\in\Lambda$ define
deterministic valuations
$v_\lambda=\mathrm{ev}_\lambda\circ\rho$.  On Boolean tested events
these are two-valued homomorphisms.  Any defender state is a convex
mixture
$\omega_{p_s} = \sum_\lambda p_s(\lambda)\,v_\lambda$.  Thus the
represented tested table lies in the convex hull of deterministic
record valuations.  Conversely, if the tested table lies in such a
convex hull, the underlying deterministic valuations define a finite
atom space $\Lambda$, the coordinate functions define $\rho$, and the
convex coefficients define the defender states $p_s$.  The
compatibility constraints ensure that $\rho$ preserves the tested
products, recovery maps, coarse-grainings, order constraints, and
contextual equivalences.  This gives (2).

Therefore Sep is equivalent to existence of a $\delta$-faithful
commutative record representation, and IJC is exactly the failure of
such a representation.
\end{proof}

\paragraph{The gluing obstruction.}
The nontrivial obstruction in the Sep/IJC split is a finite gluing
problem.  The tested word set $\mathcal T$ determines a finite family
of \emph{local word-contexts}
\[
    \mathfrak C \;=\; \{C_\alpha\}_{\alpha\in A},
\]
where each $C_\alpha\subseteq\mathcal T$ is a product-compatible
family of tested words, recovery maps, and readouts jointly
interrogated in one local protocol.  For each context $C_\alpha$, let
\[
    \mathcal R(C_\alpha)
\]
be the finite set of local Boolean record assignments reproducing the
tested behavior restricted to $C_\alpha$.  If
$C_\alpha\cap C_\beta\neq\varnothing$, there are restriction maps
\[
    r_{\alpha\beta}:\mathcal R(C_\alpha)\to
    \mathcal R(C_\alpha\cap C_\beta)
\]
obtained by forgetting all record coordinates not belonging to the
overlap.

A commuting defender exists exactly when the compatible local record
assignments glue to one global record assignment: there is
$R\in\mathcal R(\mathcal T)$ such that for every context $C_\alpha$,
$R|_{C_\alpha}=R_\alpha$, and for every overlap,
$R_\alpha|_{C_\alpha\cap C_\beta} = R_\beta|_{C_\alpha\cap C_\beta}$.
The Sep branch is precisely the case in which this finite gluing
problem has a solution after quotienting by contextual continuation
equivalence at resolution $\delta$.  The IJC branch is the failure of
such a compatible global section.

In the static nondisturbing special case, this is exactly the usual
global-section obstruction of sheaf-theoretic contextuality.  The APF
extension is that the contexts are not merely sets of jointly
measured outcomes: they are finite compatible families of typed
continuation words, tested products, recovery maps, coarse-grainings,
and state--effect pairings.

\begin{theorem}[Sep/IJC as a finite word-gluing problem]\label{thm:sep-ijc-word-gluing}
For a fixed finite tested interface $(\Gamma,Q,\mathcal T,\delta)$,
let $\mathfrak C$ be the finite cover of $\mathcal T$ by tested
product-compatible word-contexts, and let $\mathcal R$ be the
presheaf assigning to each context $C\in\mathfrak C$ the set of local
Boolean record assignments reproducing the tested behavior on $C$.
The following are equivalent.
\begin{enumerate}[label=(\arabic*),leftmargin=2em]
\item $(\Gamma,Q,\mathcal T,\delta)$ is Sep.
\item The presheaf $\mathcal R$ has a compatible global section.
\item The resolved tested continuation quotient
$A^{\Gamma,\delta}_{Q,\mathcal T}$ admits a $\delta$-faithful
commutative record representation.
\item The tested word behavior
$\mathsf B_{\Gamma,Q,\mathcal T,\delta}$ lies in the convex hull of
deterministic two-valued record valuations compatible with all
tested products, restrictions, recovery maps, and contextual
equivalences.
\end{enumerate}
Consequently, IJC is the non-existence of a compatible global
section for the finite word-context presheaf.
\end{theorem}

\begin{proof}
\emph{(1) $\Rightarrow$ (2).}  A faithful commuting defender
determines a finite Boolean record atom space $\Lambda$.  Restricting
the defender to each product-compatible word-context $C_\alpha$ gives
a local Boolean record assignment $R_\alpha\in\mathcal R(C_\alpha)$.
Since all local assignments arise by restriction from the same atom
space $\Lambda$, they agree on every overlap $C_\alpha\cap C_\beta$.
Thus a defender gives a compatible global section of the
word-context presheaf.

\emph{(2) $\Rightarrow$ (3).}  Suppose the presheaf has a compatible
global section.  The global section assigns one common Boolean record
value to every tested word class in $\mathcal T$, with agreement on
all overlaps and preservation of the tested products, recovery maps,
coarse-grainings, and state--effect values.  Taking the finite set of
such global assignments as the atom space $\Lambda$, and representing
each tested word by its evaluation function on $\Lambda$, gives a
homomorphism
\[
    \rho:A^{\Gamma,\delta}_{Q,\mathcal T}\longrightarrow\R^\Lambda.
\]
Compatibility on overlaps makes $\rho$ well-defined; product
compatibility makes it a tested-product homomorphism; quotienting by
$\equiv_{\Gamma,\delta}$ ensures it identifies exactly the
operationally unresolved word differences, so the representation is
$\delta$-faithful on the resolved tested quotient.

\emph{(3) $\Rightarrow$ (4).}  By
Theorem~\ref{thm:soundness-completeness}, a $\delta$-faithful
commutative record representation gives, via state weights
$p_s\in\Delta(\Lambda)$, a convex decomposition of the tested word
behavior into deterministic two-valued atom valuations.

\emph{(4) $\Rightarrow$ (1).}  A convex decomposition of
$\mathsf B_{\Gamma,Q,\mathcal T,\delta}$ into deterministic atom
valuations supplies a finite atom space and probability weights
reproducing the tested table; the same atom space serves as the
defender record space, and the assumed compatibility constraints make
the resulting record model a faithful commuting-extension defender.

Hence the four conditions are equivalent.  Failure of the gluing
problem is therefore the IJC obstruction.
\end{proof}

The obstruction is therefore not the failure of a single local
record, nor the mere presence of disturbance or coupling.  It is the
impossibility of gluing all locally valid commuting record
assignments into one global Boolean record representation of the
resolved tested continuation algebra.

\paragraph{Finite defender feasibility as a linear program.}
For a fixed finite tested interface $(\Gamma,Q,\mathcal T,\delta)$,
defender existence can be written as a finite feasibility problem.
Let $\Lambda_{\mathrm{det}}$ be the finite set of deterministic
two-valued record valuations compatible with the tested product
rules, recovery maps, coarse-grainings, order constraints, and
contextual continuation equivalences.  Each
$\lambda\in\Lambda_{\mathrm{det}}$ determines a column vector
\[
    v_\lambda
    \;=\;
    \bigl(v_\lambda(a)\bigr)_{a\in\mathcal O_{\mathrm{test}}},
\]
where $\mathcal O_{\mathrm{test}}$ is the finite list of tested
events, readout values, and state--effect coordinates.  Let $b_s$
be the observed tested behavior vector for preparation $s$.  A
commutative defender exists for $s$ exactly when there are
nonnegative weights $p_s(\lambda)$ satisfying
\[
    \sum_{\lambda\in\Lambda_{\mathrm{det}}} p_s(\lambda)=1,
    \qquad
    b_s
    \;=\;
    \sum_{\lambda\in\Lambda_{\mathrm{det}}} p_s(\lambda)\,v_\lambda.
\]
Equivalently, if $M$ is the incidence/evaluation matrix whose
columns are the deterministic valuations $v_\lambda$, defender
existence is the linear feasibility problem
\[
    M p_s = b_s, \qquad p_s\ge 0, \qquad \mathbf 1^\top p_s = 1.
\]
For several preparations one imposes the same deterministic
valuation set $\Lambda_{\mathrm{det}}$ and separate probability
vectors $p_s$, together with any tested preparation relations
included in the interface.  Failure of this linear system is the
finite Boolean-record obstruction; in the static nondisturbing
fragment it reduces to the Abramsky--Brandenburger / Fine--Boole
global-section feasibility test.

The number of deterministic valuations may grow exponentially with
the number of tested generators, as in standard hidden-variable and
contextuality feasibility tests.  The point is not that the
decision problem is always small, but that for every finite tested
interface it is finite, explicit, and certifiable by linear
feasibility or by a separating hyperplane witness.

\paragraph{Minimal defender construction.}
When the feasibility problem is solvable, a minimal defender is
obtained by a two-step reduction of the deterministic valuation set.
First, delete every atom $\lambda\in\Lambda_{\mathrm{det}}$ that is
assigned weight zero for every tested preparation $s$.  Second,
quotient any remaining atoms that agree on all tested coordinates
in $\mathcal O_{\mathrm{test}}$:
\[
    \lambda\sim\lambda'
    \;\Longleftrightarrow\;
    v_\lambda(a)=v_{\lambda'}(a)
    \quad\text{for every}\quad a\in\mathcal O_{\mathrm{test}}.
\]
The resulting reduced atom space $\Lambda_{\mathrm{min}}$ is minimal
for the represented finite table: no tested word, recovery map, or
state--effect value can distinguish atoms inside one equivalence
class, while atoms in different classes are separated by at least
one tested coordinate.  This is the constructive form of the
finite atom dilation lemma
(Lemma~\ref{lem:finite-atom-dilation}): every Sep interface admits a
canonical minimal defender, unique up to relabeling of atoms.

\paragraph{Degrees of IJC obstruction.}
The exact IJC verdict can be stratified by the strength of the
failed completion.
\begin{description}
\item[\emph{Probabilistic IJC.}]
The observed tested behavior does not lie in the convex hull of
deterministic Boolean record valuations.  Equivalently, no
nonnegative probability weights $p(\lambda)$ solve the defender
feasibility problem.
\item[\emph{Possibilistic IJC.}]
Even after forgetting probabilities and retaining only which tested
events are possible or impossible, there is no global Boolean
record assignment whose restrictions match the local supports.
\item[\emph{Strong IJC.}]
No deterministic global record assignment is compatible with the
support of the tested behavior at all.  This is the APF analogue of
strong contextuality in sheaf-theoretic treatments.
\item[\emph{Signed-defender gap.}]
If a signed solution exists but no nonnegative solution exists, the
obstruction is probabilistic rather than possibilistic: the local
supports can glue algebraically, but only by using negative
weights.  APF Sep requires honest nonnegative defender states, so
signed completions do not count as defenders.
\end{description}
The hierarchy is not needed for the basic Sep/IJC theorem; it is a
diagnostic for witness strength, mirroring the standard
possibilistic/probabilistic/strong contextuality distinctions while
applying them to finite tested continuation-word behavior.

\paragraph{Approximate defenders.}
The Sep/IJC theorem is stated exactly at $\varepsilon = 0$.  For
noisy finite data, fix an operational norm
$\|\cdot\|_{\mathrm{op}}$ on the tested table and define an
\emph{$\varepsilon$-defender} as a commutative record representation
whose tested values lie within $\varepsilon$ of the observed
behavior.  Because the tested table is finite, the distance from an
observed behavior to the defender polytope is well-defined:
\[
    d_{\mathrm{Sep}}(b)
    \;=\;
    \inf_{p\in\Delta(\Lambda_{\mathrm{det}})}
    \|Mp - b\|_{\mathrm{op}}.
\]
If $d_{\mathrm{Sep}}(b) > 0$, then perturbations smaller than
$d_{\mathrm{Sep}}(b)/2$ preserve the IJC verdict at the chosen
tested depth and resolution.  Approximate classification is
therefore a finite-distance question relative to the chosen norm,
noise model, and tested word set; it does not require an additional
limiting construction beyond the finite-interface theorem.

\subsection{The defender object in one place}\label{subsec:defender-object-consolidated}

For ease of reference, the formal data of a faithful commuting
defender for $(\Gamma,Q,\mathcal T,\delta)$ are collected here as a
single block.  The constituent definitions and theorems are stated
above; this subsection cross-references them in one location.

\begin{enumerate}[label=(\arabic*),leftmargin=2em]
\item \textbf{Domain.}\quad The resolved tested continuation
quotient
\[
    A^{\Gamma,\delta}_{Q,\mathcal T}
    \;=\;
    \R[\Cont_{\Gamma,Q,\mathcal T}]/\mathcal N_{\Gamma,\delta},
\]
constructed in
Theorem~\ref{thm:universal-property-tested-quotient}.

\item \textbf{Codomain.}\quad The real linearization
$\R^\Lambda$ of the Boolean record algebra
$\mathcal B_\Lambda = 2^\Lambda$ on a finite atom space
$\Lambda$ (Definition~\ref{def:boolean-record-algebra}).

\item \textbf{Morphism.}\quad A unit-preserving tested-product
homomorphism
\[
    \rho:A^{\Gamma,\delta}_{Q,\mathcal T}\to\R^\Lambda,
\]
preserving the tested products $\rho(UV)=\rho(U)\rho(V)$,
representing tested events as indicator functions, and preserving
declared recovery maps and coarse-grainings
(Definition~\ref{def:finite-commutative-record-model}).

\item \textbf{States.}\quad For each tested preparation $s$, a
probability weight
\[
    p_s\in\Delta(\Lambda)
\]
on the defender state simplex
(Definition~\ref{def:defender-state-simplex}).

\item \textbf{Faithfulness.}\quad The morphism $\rho$ is injective
on the resolved tested quotient: for tested words
$U,V\in\mathcal T$,
\[
    \rho([U])=\rho([V])
    \;\Longleftrightarrow\;
    U\equiv_{\Gamma,\delta}V
\]
(Definition~\ref{def:delta-faithfulness}).

\item \textbf{Preservation.}\quad $\rho$ preserves tested products,
identities, recovery maps, coarse-grainings, tested
positivity/order, normalization, and all tested state--effect
pairings included in $\mathcal O_{\mathrm{test}}$.

\item \textbf{Reproduction of tested behavior.}\quad For every
tested preparation $s$ and every tested
$a\in A^{\Gamma,\delta}_{Q,\mathcal T}$,
\[
    \omega_s(a)
    \;=\;
    \sum_{\lambda\in\Lambda} p_s(\lambda)\,\rho(a)(\lambda),
\]
so that the defender reproduces the tested word behavior
$\mathsf B_{\Gamma,Q,\mathcal T,\delta}$
(Definition~\ref{def:tested-word-behavior}).
\end{enumerate}

The Sep branch is the existence of an object satisfying (1)--(7);
the IJC branch is its non-existence.  By
Theorem~\ref{thm:soundness-completeness} and
Theorem~\ref{thm:sep-ijc-word-gluing}, this is equivalent to
feasibility of the finite linear program
$Mp_s = b_s,\;p_s\ge 0,\;\mathbf 1^\top p_s=1$, to existence of a
compatible global section of the word-context presheaf, and to
membership of $\mathsf B_{\Gamma,Q,\mathcal T,\delta}$ in the
convex hull of deterministic two-valued atom valuations.  When the
feasibility problem is solvable, a canonical minimal defender is
obtained by the two-step reduction described above.

\paragraph{Relation to classical global-section theorems.}
In the static nondisturbing fragment, where $\mathcal T$ contains only
context-wise joint readouts and the recovery maps are ordinary
marginal restrictions, Theorem~\ref{thm:soundness-completeness}
reduces to the standard classical-completability criterion.  The atom
space $\Lambda$ is then a space of global outcome assignments, and the
defender states $\Delta(\Lambda)$ are exactly classical probability
distributions over those assignments.  In that case, APF Sep is the
existence of a global section or classical hidden-variable model, and
APF IJC is the standard contextuality obstruction.

The APF extension is that $\mathcal T$ may contain typed sequential
continuation words, order-sensitive products, recovery maps,
coarse-grainings, abort/failure words, and resource-trigger
conditions.  The obstruction is therefore not merely failure of a
static global section.  It is failure of a single finite Boolean
record algebra to represent the entire tested continuation calculus
at the stated operational resolution.

\paragraph{Does cost affect the Sep/IJC verdict?}
For a fixed resolved tested interface
$(\Gamma,Q,\mathcal T,\delta)$, the Sep/IJC verdict is algebraic: it
asks whether the tested continuation quotient
$A^{\Gamma,\delta}_{Q,\mathcal T}$ admits a $\delta$-faithful
commutative record representation reproducing the tested word
behavior $\mathsf B_{\Gamma,Q,\mathcal T,\delta}$.  The ledger does
not change that verdict after the finite interface has been fixed.

Cost enters one level earlier.  It determines which distinctions,
tested words, recovery maps, and composites are admissible in the
finite interface in the first place.  Enlarging the capacity budget
may make additional queries or sequential words admissible, and a
refined tested family may reveal IJC that was invisible at a coarser
or cheaper interface.  Conversely, a restricted budget may leave
only a Sep fragment visible.  Resource constraints therefore do not
turn an IJC table into a Sep table; they determine which table is
physically available to classify.

\begin{lemma}[Restriction monotonicity for tested word sets]\label{lem:tested-restriction-monotonicity}
Let $\mathcal T\subseteq\mathcal T'$ be finite tested word sets for the same queried interface.  If a commuting-extension defender exists for $(\Gamma,Q,\mathcal T')$, then its restriction to the words in $\mathcal T$ is a commuting-extension defender for $(\Gamma,Q,\mathcal T)$.
\end{lemma}
\begin{proof}
A defender for $\mathcal T'$ supplies a common Boolean record space and recovery maps reproducing every tested word in $\mathcal T'$.  Restrict the represented events, products, and recovery data to the subset $\mathcal T$.  Faithfulness on the larger tested quotient implies faithfulness on the smaller quotient, because any equality or separation visible in $\mathcal T$ is already included among the $\mathcal T'$ tests.
\end{proof}

\begin{remark}[Refinement can reveal IJC]\label{rem:refinement-can-reveal-ijc}
The converse of Lemma~\ref{lem:tested-restriction-monotonicity} need not hold.  A coarse tested word set may admit a defender while a refined set $\mathcal T'$ exposes an order effect, recovery obstruction, or contextual incompatibility.  Operationally, Sep is monotone under forgetting tests, while IJC can be revealed by adding tests.  This is a feature of an interface-relative theorem, not a defect: the branch classification always refers to the finite behavior actually tested.
\end{remark}



\begin{proposition}[Tested-set and resolution refinement stability]\label{prop:resolution-refinement-stability}
Fix a queried interface $(\Gamma,Q)$.
\begin{enumerate}[label=(\roman*),leftmargin=2em]
\item If $\mathcal T\subseteq\mathcal T'$ and the finer tested interface $(\Gamma,Q,\mathcal T')$ is Sep, then the coarser tested interface $(\Gamma,Q,\mathcal T)$ is Sep.
\item If $\mathcal T\subseteq\mathcal T'$ and the coarser tested interface $(\Gamma,Q,\mathcal T)$ is IJC, then the finer tested interface $(\Gamma,Q,\mathcal T')$ is IJC.
\item Coarsening operational resolution can turn an IJC witness invisible by quotienting away the separating word classes; refinement can reveal IJC, but cannot rescue a coarser tested family that already has no defender.
\end{enumerate}
\end{proposition}
\begin{proof}
Parts (i) and (ii) are immediate from Lemma~\ref{lem:tested-restriction-monotonicity}: any defender for the finer tested family restricts to a defender for the coarser family.  Therefore a missing defender at the coarser level also forbids a defender at the finer level.  Part (iii) records the separate role of operational resolution: changing resolution changes the quotient itself by identifying or resolving word classes, so the branch label is always stated relative to both $\mathcal T$ and the chosen resolution.
\end{proof}

\begin{proposition}[Directed finite-test limits]\label{prop:directed-finite-test-limits}
Let $(\mathcal T_i)_{i\in I}$ be a directed family of finite tested word sets with $\mathcal T_i\subseteq\mathcal T_j$ for $i\le j$, fixed context $\Gamma$, fixed query family $Q$, and compatible quotient/restriction maps.
\begin{enumerate}[label=(\roman*),leftmargin=2em]
\item If some finite stage $(\Gamma,Q,\mathcal T_i)$ is IJC, then every later stage $\mathcal T_j\supseteq\mathcal T_i$ is IJC.
\item A global Sep limit requires a coherent family of defenders whose Boolean records and recovery maps restrict compatibly along the directed system.  Sep at each finite stage separately is not by itself enough unless these compatibility maps are supplied.
\item If the finite stages have uniformly bounded atom number or compact probability-simplex data and compatible restrictions, then a projective-limit defender exists for the limit fragment.  Without such compactness/coherence assumptions, the finite theorem makes no global claim.
\end{enumerate}
\end{proposition}
\begin{proof}
Part (i) is restriction monotonicity contrapositive: a defender at a later stage would restrict to a defender at the earlier IJC stage.  Part (ii) records that defenders are structured objects, not just truth values; they must agree on overlaps.  Part (iii) is the standard finite-projective compactness argument: compatible finite Boolean/probability data have a nonempty inverse limit under the stated boundedness/compactness assumptions.  The proposition is modest; APF's main theorem remains finite-interface unless this extra limiting structure is provided.
\end{proof}


\begin{remark}[Finite-to-continuous limiting scope]\label{rem:finite-to-continuous-scope}
APF's Sep/IJC theorem is a finite-interface theorem.  Continuous outcomes or infinite word families enter only through a directed family of finite resolutions.  A typical limiting construction consists of finite outcome partitions
\(\mathcal P_1\preceq\mathcal P_2\preceq\cdots\), finite tested word sets
\(\mathcal T_1\subseteq\mathcal T_2\subseteq\cdots\), and compatible quotient/restriction maps.

Three facts are immediate.
\begin{enumerate}[label=(\roman*),leftmargin=2em]
\item If some finite stage is IJC, every finer stage containing it is IJC.
\item A Sep verdict at every finite stage yields a global Sep limit only when the defenders form a coherent projective family.
\item If the resolved floor \(\mu_{\Gamma,\delta}\) collapses to zero, if the number or cost of required records diverges, or if tested products do not remain closed under refinement, the finite-interface classification does not automatically pass to the ideal limit.
\end{enumerate}
Thus APF does not assume that ideal continuum protocols are physically available.  They are additional analytic completions of finite interrogations and require separate compactness, coherence, noncollapse, and closure hypotheses.
\end{remark}

\begin{definition}[$\varepsilon$-faithful defender]\label{def:epsilon-defender}
Fix a finite tested word set $\mathcal T$, an operational norm $\|\cdot\|_{\mathrm{op}}$, and a resolution scale.  An $\varepsilon$-faithful commuting-extension defender for $(\Gamma,Q)$ is a finite Boolean record completion whose represented continuation-word statistics agree with the queried operational statistics to within $\varepsilon$ in $\|\cdot\|_{\mathrm{op}}$ for every word in $\mathcal T$, while preserving the same common-record and recovery structure as an exact defender.
\end{definition}

\begin{remark}[Exact classification versus noisy experiments]\label{rem:epsilon-ijc}
The Sep/IJC theorem below is the exact $\varepsilon=0$ classification.  Approximate IJC requires stability bounds specifying the tested word depth, operational norm, noise model, and resolution scale.  That quantitative theory is important for experiments and simulations, but it is not needed for the exact branch theorem proved here.
\end{remark}


\begin{proposition}[Finite-margin robustness for approximate defenders]\label{prop:finite-margin-robustness}
Fix a finite tested word set $\mathcal T$ and operational norm $\|\cdot\|_{\mathrm{op}}$.  Suppose the observed finite table is separated from the nearest commuting-extension defender by margin
\[
 m:=\inf_{R\in\Def_{\Bool}}\max_{W\in\mathcal T}
 \|\Pi_\Gamma(W)-\Pi_R(W)\|_{\mathrm{op}}>0 .
\]
Then every perturbation of the observed finite table by at most $\varepsilon<m/2$ remains IJC at the tested depth: no commuting-extension defender can reproduce the perturbed table within error $\varepsilon$.
Conversely, if an exact Sep defender has tolerance slack $s>0$ relative to the stated resolution threshold, then perturbations smaller than $s$ preserve the Sep verdict for the same defender.
\end{proposition}
\begin{proof}
For any Boolean defender $R$, the unperturbed table is at distance at least $m$ from $R$.  If the observed table is perturbed by at most $\varepsilon$, the triangle inequality gives distance at least $m-\varepsilon$ from every $R$.  If $\varepsilon<m/2$, then $m-\varepsilon>\varepsilon$, so no $\varepsilon$-faithful defender exists for the perturbed table.  The Sep statement is the same finite-table continuity argument in the opposite direction: if all represented tested values lie within the admissible tolerance with slack $s$, perturbations smaller than $s$ keep them inside that tolerance.  No Hilbert-space or probability structure is used.
\end{proof}
\begin{proposition}[Boolean feasibility for defenders]\label{prop:boolean-feasibility-criterion}
For a finite tested quotient $\AQtest$, defender existence is equivalent to feasibility of the following finite Boolean-record constraint system.  There must exist a finite atom set $\Lambda$, functions $f_W:\Lambda\to\mathbb R$ for tested words $W\in\mathcal T$, and tested state weights $p_s\in\Delta(\Lambda)$ such that:
\begin{enumerate}[label=(\roman*),leftmargin=2em]
\item $f_1=1_\Lambda$;
\item for every tested product $UV\in\mathcal T$, $f_{UV}=f_U f_V$ pointwise, with the prescribed recovery/coarse-graining maps also respected;
\item if $U\equiv_{\mathrm{cont}}V$ in the tested quotient, then $f_U=f_V$;
\item if $U$ and $V$ are resolved as distinct in $\AQtest$, then $f_U\ne f_V$ on at least one atom;
\item tested positive events are represented by pointwise nonnegative functions, and probabilistic effects by functions in $[0,1]^\Lambda$;
\item all tested state--effect values are reproduced:
\[
  \omega_s(W)=\sum_{\lambda\in\Lambda}p_s(\lambda)f_W(\lambda)
  \quad\text{for every tested state }s\text{ and word }W.
\]
\end{enumerate}
Consequently IJC is the infeasibility of this finite Boolean constraint system for the stated $\mathcal T$ and resolution.
\end{proposition}
\begin{proof}
A defender supplies exactly such atom functions by evaluating each tested Boolean record on each atom.  Unit, product, recovery, positivity, and state--effect reproduction are the defender axioms written as finite equations/inequalities; faithfulness supplies the separation condition in (iv).  Conversely, a feasible family of atom functions defines the reduced Boolean record algebra generated by their level sets, with state weights $p_s$ and recovery maps inherited from the specified equations.  This is a commuting-extension defender.  The criterion is nontrivial because it is a finite simultaneous feasibility problem: the tested products, recovery maps, positivity inequalities, and state--effect equations may be mutually inconsistent, as in the CHSH witness.
\end{proof}

\begin{definition}[Faithful commutative continuation representation]\label{def:faithful-commutative-representation}
A faithful commutative continuation representation of the finite tested continuation quotient $\AQtest$ is an injective partial-algebra homomorphism
\[
 \rho:\AQtest\to \R^\Lambda
\]
for some finite set $\Lambda$, where $\R^\Lambda$ is the reduced commutative algebra of real-valued functions on the Boolean atoms $\Lambda$.
It preserves exactly the finite structures present in the tested quotient:
\begin{enumerate}[label=(\roman*),leftmargin=2em]
\item the tested unit: $\rho(1)=1_\Lambda$;
\item tested products $UV$ whenever $U,V,UV\in\mathcal T$;
\item tested order/positivity: $a\in\AQtest^+$ implies $\rho(a)\geq0$ pointwise on $\Lambda$;
\item tested recovery and coarse-graining maps included in $\mathcal T$;
\item tested continuation values and state--effect pairings resolved by the context;
\item the resolved quotient relation: $\rho(U)=\rho(V)$ iff $U\equiv_{\mathrm{cont}}V$ inside $\AQtest$.
\end{enumerate}
Thus \emph{faithful} means injective on resolved tested continuation classes and faithful to the finite operational data actually tested.  It does not mean Hilbert-space isometry, norm preservation beyond the chosen operational norm, or preservation of untested infinitary structure.
\end{definition}

\begin{remark}[Why the commutative target is finite and reduced]\label{rem:finite-lambda}
The finite set $\Lambda$ is not assumed for arbitrary infinite operational theories.  It belongs to the finite-interface regime: the queried family is finite and only finitely many word classes are operationally resolved at the chosen context resolution.  A finite reduced commutative record algebra is therefore isomorphic to a finite algebra of functions on its atoms,
\[
 \mathcal B \cong \R^\Lambda .
\]
The word \emph{reduced} is operationally load-bearing.  Nilpotent elements in a commutative algebra vanish on every classical atom and therefore cannot label a stable Boolean record outcome.  In APF language they are record-null or infinitesimal bookkeeping directions: they may occur in a non-reduced formal algebraic completion, but they have no independent finite discriminator at the stated resolution.  Sep is therefore defined by faithful representation into a reduced Boolean record algebra, not by arbitrary commutative algebra representability with nilpotent artifacts.
\end{remark}

\begin{proposition}[Minimal defender quotient and non-uniqueness]
\label{prop:minimal-defender-quotient}
If a defender exists for a finite tested interface, it need not be unique.  One may freely split a Boolean atom into several indistinguishable subatoms, add unused null atoms, or attach extra charged classical records without changing the tested continuation statistics.  However, every defender has a canonical finite reduction: identify two atoms \(\lambda,\lambda'\) when all represented tested words, recovery maps, and tested effects take the same values on them.  The quotient Boolean algebra generated by the distinct tested evaluation profiles is minimal among defenders that reproduce the same finite tested behavior, and it is unique up to Boolean isomorphism.
\end{proposition}
\begin{proof}
The equivalence relation is defined by equality of the full finite evaluation profile of all tested represented words and readout/recovery data.  Quotienting by it preserves every tested state--effect value because no tested functional separates equivalent atoms.  Any defender reproducing the same finite behavior must distinguish at least the inequivalent evaluation profiles, otherwise it would identify two tested profiles that the original finite interface separates.  Hence the reduced generated Boolean algebra is minimal and any two such reductions have the same finite set of evaluation profiles, giving uniqueness up to relabeling of atoms.
\end{proof}

\begin{proposition}[Generated Boolean algebra and kernel ideal of the minimal defender]
\label{prop:generated-boolean-kernel-minimal-defender}
Let \(\rho:\AQtest\to\R^\Lambda\) be any faithful commutative continuation representation of a finite tested interface.  Let \(\mathcal G_\rho\) be the finite Boolean algebra generated by the level sets of all represented tested effects, words, and recovery readouts \(\rho(W)\) for \(W\in\mathcal T\).  Define the profile map
\[
  q_\rho:\Lambda\longrightarrow
  \prod_{W\in\mathcal T}\mathrm{im}(\rho(W)),
  \qquad
  q_\rho(\lambda)=(\rho(W)(\lambda))_{W\in\mathcal T},
\]
including the tested recovery/readout coordinates.  The kernel ideal
\[
  I_\rho:=\{B\in\mathcal G_\rho: q_\rho(B)=\varnothing\}
\]
identifies exactly those Boolean subrecords that have no distinct tested evaluation profile.  The quotient
\[
  \mathcal B_{\min}(\rho):=\mathcal G_\rho/I_\rho
\]
is the minimal defender Boolean algebra for the tested interface: every defender reproducing the same finite tested behavior factors through \(\mathcal B_{\min}(\rho)\), and two minimal defenders are isomorphic by the unique relabeling of their tested evaluation profiles.
\end{proposition}
\begin{proof}
The finite family \(\{\rho(W):W\in\mathcal T\}\), together with the tested recovery/readout coordinates, partitions \(\Lambda\) into finitely many atoms with identical tested profiles.  Passing from \(\mathcal G_\rho\) to \(\mathcal B_{\min}(\rho)\) deletes empty profile classes and identifies no two nonempty profile classes.  Therefore every tested state--effect value and every tested product/recovery equation is preserved.  Any other defender that reproduces the same finite behavior must assign the same values to all tested words on each realized profile, so its Boolean algebra maps onto the algebra of nonempty tested profiles.  This gives the claimed factorization and uniqueness up to Boolean isomorphism.  The construction is finite and internal to the tested quotient; it does not assert uniqueness of nonminimal defenders, which may add unused atoms, split indistinguishable atoms, or carry extra charged records.
\end{proof}


\begin{theorem}[Defender implies faithful commutative representation]\label{thm:defender-to-commutative}
If $(\Gamma,Q)$ admits a faithful commuting-extension defender for a finite tested word set $\mathcal T$, then $\AQtest$ admits a faithful commutative continuation representation.
\end{theorem}
\begin{proof}
Let $(\widetilde\Omega,\mathcal B,\{\widetilde d_i\})$ be a faithful commuting-extension defender.  Because $\mathcal B$ is Boolean, every event $b\in\mathcal B$ can be represented by its indicator function $1_b\in\R^{\widetilde\Omega}$, and pointwise multiplication is commutative.  Map each generator $E_i$ to the indicator/update function representing $\widetilde d_i$.  Extend multiplicatively and linearly to words and finite linear combinations.  Since all represented events live in the same Boolean algebra, the image is commutative.  Since the defender reproduces all words in $\mathcal T$ and is faithful on the tested quotient, two tested words have the same image exactly when they are identified in $\AQtest$.  Hence the induced homomorphism from $\AQtest$ into $\R^{\widetilde\Omega}$ is injective and commutative.
\end{proof}

\begin{theorem}[Faithful commutative representation implies defender]\label{thm:commutative-to-defender}
If $\AQtest$ admits a faithful commutative continuation representation, then $(\Gamma,Q)$ admits a commuting-extension defender for the finite tested word set $\mathcal T$.
\end{theorem}
\begin{proof}
Let
\[
 \rho:\AQtest\to\R^\Lambda
\]
be a faithful commutative continuation representation.  The tested interface is finite, so fix the finite tested word set
\[
 \mathcal T\subset \AQtest
\]
containing the query generators, the recovery/coarse-graining words, and all continuation words whose behavior the defender is required to reproduce.  Define an equivalence relation on $\Lambda$ by
\[
 \lambda\sim_{\mathcal T}\lambda'
 \quad\Longleftrightarrow\quad
 \rho(W)(\lambda)=\rho(W)(\lambda')
 \text{ for every }W\in\mathcal T.
\]
Let
\[
 \widetilde\Omega:=\Lambda/{\sim_{\mathcal T}}
\]
be the finite set of joint evaluation-profile atoms, and let $\mathcal B=2^{\widetilde\Omega}$ be its Boolean algebra.  For each tested word $W\in\mathcal T$, the represented function $\rho(W)$ is constant on each atom of $\widetilde\Omega$, so it descends to a Boolean-record observable on $\widetilde\Omega$.  In particular, each queried distinction $E_i$ and each tested sequential word has a jointly defined Boolean representative.

This construction uses the joint finite evaluation profile of all tested represented words, not merely the level sets of the individual generators.  Therefore the resulting Boolean record space preserves the full tested continuation-word behavior, including recovery and coarse-graining data.  Because $\rho$ is faithful, two tested words induce the same Boolean-record behavior exactly when they are equivalent under $\equiv_{\mathrm{cont}}$.  Hence $(\widetilde\Omega,\mathcal B)$ with these represented events and recovery maps is a commuting-extension defender.
\end{proof}



\begin{center}
\fbox{\begin{minipage}{0.92\linewidth}
\textbf{Nontriviality of the representation theorem.}
The theorem is not the tautology that ``Sep means a defender exists.''
It identifies three independently specified finite objects:
\begin{enumerate}[label=(\arabic*),leftmargin=2em]
\item a common Boolean-record defender for the tested interface;
\item a $\delta$-faithful commutative representation of the tested
continuation quotient $A^{\Gamma,\delta}_{Q,\mathcal T}$
(Definitions~\ref{def:finite-commutative-record-model}
and~\ref{def:delta-faithfulness});
\item feasibility of a finite Boolean-record constraint system
preserving tested products, recovery maps, order, normalization, and
state--effect pairings
(Theorem~\ref{thm:boolean-feasibility-separating-witnesses}).
\end{enumerate}
The IJC branch is the failure of this finite simultaneous feasibility
problem.  Such failure may be certified by finite separating witnesses,
generalizing Bell, Boole, Fine, and sheaf-contextuality witnesses to
tested continuation-word behavior.
\end{minipage}}
\end{center}

\begin{theorem}[Self-contained Sep/IJC representation theorem]
\label{thm:sep-ijc-representation}
Fix an admissible context \(\Gamma\), a finite query family
\(Q=\{d_1,\ldots,d_n\}\subset \mathcal D_\Gamma\), and a fixed
operational resolution \(\delta\).  Let \(\mathcal T\) be a finite
tested word set generated by \(Q\), containing the tested identity
word and closed under every product, recovery map, and
coarse-graining operation actually tested at resolution \(\delta\).

Let \(\equiv_{\Gamma,\delta}\) denote context-complete continuation
equivalence at this resolution: \(U\equiv_{\Gamma,\delta}V\) iff no
admissible continuation in \(\Gamma\), at resolution \(\delta\), separates
the tested continuation profiles of \(U\) and \(V\).  Let
\[
  \mathcal A^{\Gamma,\delta}_{Q,\mathcal T}
\]
be the finite tested continuation-word quotient: the real span of
tested word classes in \(\mathcal T\), modulo
\(\equiv_{\Gamma,\delta}\), with multiplication defined exactly for
tested concatenations \(UV\in\mathcal T\).

Then the following are equivalent:
\begin{enumerate}[label=(\roman*),leftmargin=2em]
\item The finite interface \((\Gamma,Q,\mathcal T,\delta)\) is Sep.

\item There exists a finite Boolean record space \(\Lambda\) and a
faithful commutative continuation representation
\[
  \rho:\mathcal A^{\Gamma,\delta}_{Q,\mathcal T}\to \mathbb R^\Lambda
\]
such that:
\begin{enumerate}[label=(\alph*),leftmargin=2em]
\item \(\rho(1)=1_\Lambda\);
\item whenever \(U,V,UV\in\mathcal T\), one has
\(\rho(UV)=\rho(U)\rho(V)\) pointwise;
\item tested recovery and coarse-graining maps are preserved;
\item tested positive events are represented by pointwise nonnegative
functions, and probabilistic effects by functions in \([0,1]^\Lambda\);
\item all tested state--effect and continuation values are reproduced;
\item \(\rho(U)=\rho(V)\) iff
\(U\equiv_{\Gamma,\delta}V\) inside the finite tested quotient.
\end{enumerate}

\item The finite Boolean-record feasibility system associated to
\(\mathcal T\) has a solution: there exist atom functions
\(f_W:\Lambda\to\mathbb R\) and tested state weights
\(p_s\in\Delta(\Lambda)\) reproducing all tested products,
recovery maps, positivity constraints, and state--effect values,
while separating exactly the resolved tested word classes.
\end{enumerate}

Consequently,
\[
  (\Gamma,Q,\mathcal T,\delta)\in\IJC
  \quad\Longleftrightarrow\quad
  \mathcal A^{\Gamma,\delta}_{Q,\mathcal T}
  \text{ admits no faithful commutative continuation representation}.
\]
Equivalently, IJC is the infeasibility of the corresponding finite
Boolean-record constraint system.
\end{theorem}

\begin{proof}
A Sep interface, by definition, has a single admissible classical
record completion for all tested queries.  Its record atoms form a
finite Boolean space \(\Lambda\).  Evaluating each tested word on
those atoms gives functions \(f_W:\Lambda\to\mathbb R\), hence a
commutative representation into \(\mathbb R^\Lambda\).  Preservation
of the tested unit, products, recovery maps, positivity, and
state--effect pairings follows from the fact that the defender
reproduces the tested operational behavior.  Faithfulness is exactly
the requirement that the representation identify no more and no
less than the continuation quotient identifies.

Conversely, given such a faithful representation
\(\rho:\mathcal A^{\Gamma,\delta}_{Q,\mathcal T}\to\mathbb R^\Lambda\),
take \(\Lambda\) as the defender's atom space and let the Boolean
algebra be \(2^\Lambda\).  The represented tested effects and words,
together with the specified recovery maps and state weights, define
a common commuting record model reproducing the finite interface.
Faithfulness guarantees that this defender neither collapses a
resolved distinction nor separates words already quotient-identified
by admissible continuations.

The nontrivial step is the middle equivalence: existence of a
defender is not merely a restatement of a definition, but a finite
simultaneous Boolean feasibility problem.  The same atom space
\(\Lambda\) must satisfy all tested product equations, recovery
constraints, positivity requirements, and state--effect values at
once.  These constraints can be mutually inconsistent; the CHSH
table gives an explicit finite case where every Boolean assignment
satisfies the CHSH bound while the observed tested table violates
it.  Thus IJC is witnessed by infeasibility of the finite
Boolean-record system, not by terminological stipulation.

The equivalence therefore does not follow from naming the Sep branch.
It follows from the quotient universal property
(Theorem~\ref{thm:universal-property-tested-quotient}), the atom-dilation
lemma (Lemma~\ref{lem:finite-atom-dilation}), and the Boolean feasibility
theorem (Theorem~\ref{thm:boolean-feasibility-separating-witnesses}).
The quotient universal property ensures that only contextually resolved
continuation differences are represented.  The atom-dilation lemma
reduces any commutative representation to a finite Boolean record space.
The feasibility theorem then identifies exactly when such a record space
can preserve all tested products, recovery maps, order relations, and
state--effect pairings.  Failure of this finite constraint system is
precisely the IJC obstruction.
\end{proof}

\paragraph{What is new relative to global-section criteria.}
In the static nondisturbing special case, the theorem reduces to a
familiar form.  If the tested family $\mathcal T$ contains only
context-wise joint readouts, if all recovery maps are trivial, and if no
sequential order data are tested, then a commutative record defender is
exactly a global Boolean atom space whose marginals reproduce the
queried empirical contexts.  In that restricted fragment, the APF
criterion recovers the Fine--Boole--sheaf global-section condition
rather than replacing it.

The nontrivial content of the APF theorem lies in the additional finite
continuation structure.  The represented object is not merely a family
of static marginal distributions, but the tested continuation quotient
\[
    A^{\Gamma,\delta}_{Q,\mathcal T}.
\]
A defender must preserve tested sequential products, typed domains of
composition, recovery and coarse-graining maps, contextual continuation
equivalence, tested positivity/order, normalization, and all finite
state--effect pairings included in the interface.  Thus the APF
obstruction is not only the absence of a global assignment to static
outcomes.  It is the absence of a single finite Boolean record capable
of representing the entire tested continuation calculus at the stated
operational resolution.

\begin{corollary}[IJC forces noncommutative faithful bookkeeping]\label{cor:ijc-noncommutative}
If $(\Gamma,Q,\mathcal T)\in\IJC$, then no faithful algebraic realization of the queried continuation-word behavior can have a fully commutative image algebra.  In any faithful algebraic realization whose product represents sequential word composition, the image algebra is noncommutative; consequently some finite represented word classes have noncommuting images.
\end{corollary}
\begin{proof}
If the image algebra of a faithful algebraic realization were fully commutative, then the realized finite word behavior would define a faithful commutative continuation representation, contradicting Theorem~\ref{thm:sep-ijc-representation}.  Since the realization is generated by finite continuation-word classes, noncommutativity of the faithful image is witnessed by some finite pair of represented word classes.
\end{proof}


\begin{warningbox}
This theorem does not say that every physical system is IJC.  It says that when a finite queried interface has no faithful commuting continuation completion, any faithful bookkeeping of its queried continuation behavior must be noncommutative.  Classical systems remain Sep.
\end{warningbox}

\section{Worked models: Sep, IJC, and coupled-but-classical}\label{sec:worked-models}

\paragraph{Why no platform ledger is included here.}
This supplement is a formal finite-interface representation theorem,
not a platform-calibration paper.  The examples below are chosen to
test the logical distinction between Sep, IJC, disturbance, coupling,
and shared support.  A full apparatus ledger would require
platform-specific error models, detector efficiencies, record
lifetimes, stabilization costs, control resources, and calibration
maps.  Those data belong to downstream application papers.  The
present work supplies the theorem and the ledger rules that such
applications must satisfy.

\subsection{Formal CHSH fragment as a Boolean-defender obstruction}\label{subsec:formal-chsh-fragment}

In the CHSH fragment, $Q=\{A_0,A_1,B_0,B_1\}$, the tested word set
contains only the compatible pair words $A_iB_j$ for $i,j\in\{0,1\}$,
and the recovery maps are the marginal restrictions to $A_i$ and
$B_j$.  A commutative defender is a finite atom space
\[
    \Lambda \;\subseteq\; \{\pm 1\}^{\{A_0,A_1,B_0,B_1\}}
\]
with probability weights $p(\lambda)\in[0,1]$,
$\sum_\lambda p(\lambda)=1$.  The represented correlations are
\[
    E_{ij}
    \;=\;
    \sum_{\lambda\in\Lambda}p(\lambda)\,A_i(\lambda)\,B_j(\lambda).
\]
For any deterministic atom
$\lambda=(a_0,a_1,b_0,b_1)\in\{\pm 1\}^4$,
\[
    a_0 b_0 + a_0 b_1 + a_1 b_0 - a_1 b_1
    \;=\;
    a_0(b_0+b_1) + a_1(b_0 - b_1)
    \;\in\;
    \{-2,+2\}
\]
because exactly one of $b_0+b_1, b_0-b_1$ is zero.  Taking convex
combinations gives, for every defender,
\[
    \bigl|E_{00}+E_{01}+E_{10}-E_{11}\bigr| \;\le\; 2.
\]
Any tested table violating this inequality lies outside the convex
hull of deterministic two-valued record valuations and hence admits
no faithful commutative defender on this fragment.  This is the
standard Fine--Boole obstruction recovered as the static
nondisturbing special case of the soundness/completeness theorem
(Theorem~\ref{thm:soundness-completeness}).

The CHSH fragment is thus a worked instance of the IJC branch as
infeasibility of the finite Boolean-record constraint system.  The
same logical structure --- $\mathcal T$, $A^{\Gamma,\delta}_{Q,\mathcal T}$,
defender states $\Delta(\Lambda)$, separating linear functional
$E_{00}+E_{01}+E_{10}-E_{11}$ --- generalizes to richer tested
fragments by allowing $\mathcal T$ to contain sequential words and
recovery maps; the obstruction then ranges over the full tested
continuation calculus rather than only the static marginal table.

\begin{example}[Fully worked Sep defender: two compatible binary measurements]
\label{ex:explicit-sep-defender}
Consider a finite interface with two compatible binary readouts \(A,B\in\{0,1\}\).  The operational atom space is
\[
  \Lambda=\{00,01,10,11\},
\]
and the Boolean record algebra is \(\mathcal B=2^\Lambda\cong\R^\Lambda\) through indicator functions.  Define the represented effects
\[
  \rho(E_A)(ab)=a,
  \qquad
  \rho(E_B)(ab)=b,
\]
where \(ab\in\Lambda\).  The unit is \(\rho(1)=1_\Lambda\).  Products are pointwise, so \(\rho(E_AE_B)=\rho(E_A)\rho(E_B)=\rho(E_BE_A)\).  Recovery maps are simply the coordinate projections
\[
  r_A(ab)=a,
  \qquad r_B(ab)=b,
\]
and the joint recovery is \(r_{AB}(ab)=(a,b)\).

A tested preparation is a probability vector \(p_s\in\Delta(\Lambda)\).  For example,
\[
  \omega_s(E_A)=\sum_{ab\in\Lambda}p_s(ab)a,
  \qquad
  \omega_s(E_BE_A)=\sum_{ab\in\Lambda}p_s(ab)ba .
\]
The representation is order-preserving because indicators are nonnegative pointwise, unit-preserving by construction, and faithful on the tested quotient: two tested words have the same represented function exactly when they produce the same coordinate-readout behavior for every admissible preparation.  Thus \((\Lambda,2^\Lambda,\rho,r_A,r_B)\) is an explicit commuting-extension defender.
\end{example}

\begin{example}[Sep: two classical bits]
Let
\[
 \Omega=\{00,01,10,11\}.
\]
Let $d_1$ read the first bit and $d_2$ read the second bit.  The common record algebra is
\[
 \mathcal B=2^\Omega.
\]
Both distinctions are events in $\mathcal B$.  Query operations commute because reading the first bit and then the second yields the same record as reading the second and then the first:
\[
 E_1E_2\equiv_{\mathrm{cont}}E_2E_1.
\]
Assign finite costs \(\kappa_\Gamma(d_1)=a>0\), \(\kappa_\Gamma(d_2)=b>0\), and suppose the joint register has capacity \(C_\Gamma\ge a+b\).  Then the Boolean defender is not zero-cost, but it is ledger-admissible.  Thus \((\Gamma,\{d_1,d_2\})\in\Sep\).
\end{example}

\begin{example}[Sep with a shared bus: nonzero interface cost is still classical]\label{ex:classical-shared-bus}
Let two classical bits $A,B\in\{0,1\}$ be read through a shared finite bus $M$.  The raw record space is
\[
 \Omega=\Omega_A\times\Omega_B\times\Omega_M,
\]
and the queried distinctions are $d_A$ and $d_B$.  Suppose reading either bit requires activating the shared bus with cost $c_M>0$, while maintaining the individual bit records costs $a>0$ and $b>0$.  The least joint enforcement cost is not $a+b+2c_M$ because the bus is shared; conservatively one writes
\[
 \kappa_\Gamma(\{d_A,d_B\})=a+b+c_M.
\]
If $C_\Gamma\ge a+b+c_M$, the common Boolean record algebra $2^\Omega$ defends the queried words.  Thus the interface is Sep even though $\kappa_{\Gamma,\mathrm{int}}=c_M>0$.  The example illustrates why ledger independence is required before one may add costs, and why nonzero interface cost is not IJC.
\end{example}


\begin{example}[IJC: explicit CHSH finite table]\label{ex:chsh-finite-table}
Consider a bipartite finite interface with settings $A_0,A_1$ on Alice's side and $B_0,B_1$ on Bob's side, each with outcomes $\pm1$.  Let the tested word set $\mathcal T_{\rm CHSH}$ contain the four compatible joint queries $(A_i,B_j)$ and their outcome readouts.  Suppose the resolved correlations are
\[
  E_{00}=E_{01}=E_{10}=\frac{1}{\sqrt2},
  \qquad
  E_{11}=-\frac{1}{\sqrt2},
\]
with unbiased one-site marginals.  Then the CHSH expression is
\[
  S:=E_{00}+E_{01}+E_{10}-E_{11}=2\sqrt2>2.
\]
Any commuting-extension defender for this nondisturbing Bell interface supplies one joint Boolean record
\[
  \lambda=(a_0,a_1,b_0,b_1)\in\{\pm1\}^4
\]
with a probability distribution $p(\lambda)$ reproducing the four observed pair statistics.  For every deterministic atom $\lambda$,
\[
  a_0b_0+a_0b_1+a_1b_0-a_1b_1\in\{-2,2\},
\]
so averaging over $p(\lambda)$ gives $|S|\le2$.  The displayed finite table violates that bound.  Therefore no faithful common Boolean record defender exists for $\mathcal T_{\rm CHSH}$, and the queried interface is IJC.

This is an end-to-end finite-table witness: the tested set is specified, the proposed defender's atom space is specified, the inequality obeyed by every defender is derived, and the observed table violates it.  The argument uses only the Boolean-record defender criterion; it does not assume Hilbert space.

Equivalently, enumerate the sixteen deterministic atoms $\lambda\in\{\pm1\}^4$ and ask for probabilities $p_\lambda\ge0$ satisfying
\[
  \sum_\lambda p_\lambda=1,
  \qquad
  \sum_\lambda p_\lambda a_i b_j = E_{ij}
  \quad(i,j\in\{0,1\}),
\]
together with the one-site marginal constraints.  This finite linear program is infeasible because its feasible region obeys $|S|\le2$.  Thus the APF defender obstruction is a classical Boolean feasibility obstruction.
\end{example}

\begin{remark}[CHSH, LP certificates, and NPA]\label{rem:chsh-lp-npa}
For Bell tables, ruling out an APF Sep defender is a linear-programming problem over deterministic Boolean atoms.  NPA semidefinite relaxations answer a different downstream question: whether a table has a quantum or commuting-operator realization.  A CHSH value $2<S\le2\sqrt2$ is APF-IJC because no Boolean defender exists, but it may still be Hilbert/Born-admissible.  Accordingly there is generally no NPA infeasibility certificate for the usual quantum CHSH violation; the NPA hierarchy certifies compatibility with, or exclusion from, quantum/commuting-operator models, not failure of a classical Boolean defender.  A value beyond Tsirelson's bound would fail the downstream quantum certificate as well.  This separation is exactly the intended architecture: first ask whether a finite Boolean defender exists; only afterward ask which noncommutative representation class is admissible.
\end{remark}

\begin{example}[IJC: explicit two-query sequential table]\label{ex:two-query-table}
Consider a finite interface with two binary queried distinctions $X$ and $Z$.  Write $x_+$ and $x_-$ for the two $X$-stable continuation profiles and $z_+$ and $z_-$ for the two $Z$-stable continuation profiles.  The queried operations have the following finite operational behavior at the chosen resolution:
\begin{center}
\begin{tabular}{@{}p{0.18\linewidth}p{0.18\linewidth}p{0.54\linewidth}@{}}
\toprule
Input profile & Query applied & Resolved output profile \\
\midrule
$x_\pm$ & $E_X$ & $x_\pm$ (repeatability)\\
$z_\pm$ & $E_Z$ & $z_\pm$ (repeatability)\\
$x_\pm$ & $E_Z$ & $z_+$ or $z_-$, with the prior $X$-profile not recoverable without extra enforcement\\
$z_\pm$ & $E_X$ & $x_+$ or $x_-$, with the prior $Z$-profile not recoverable without extra enforcement\\
\bottomrule
\end{tabular}
\end{center}
The continuation test that closes the table is order sensitive:
\[
 E_XE_ZE_X \not\equiv_{\mathrm{cont}} E_ZE_XE_Z .
\]
Equivalently, after the word $XZX$ the admissible next continuations include the $X$-stable repeat profile, while after the word $ZXZ$ they include the $Z$-stable repeat profile, and these profiles are separated by a finite discriminator at the stated resolution.

The table is not merely classical readout with a nuisance disturbance.  A classical invasive model may disturb later readings, but a Sep defender must embed the whole tested word family into one common Boolean record completion and reproduce the recovery/coarse-graining behavior for every tested word.  A Boolean defender may reproduce order sensitivity only by carrying an additional order memory.  If that memory is part of the defended system, it must appear as an explicit record atom in the enlarged Boolean algebra and must be tested or charged as part of the defended word family; if it is not included, the defender has not reproduced the tested continuation behavior.  For the stated finite interface, no defender using only the queried common record reproduces the displayed continuation inequivalence.  Therefore this queried interface is IJC.
\end{example}

\begin{remark}[Relation to a qubit]
The preceding model captures the operational skeleton of incompatible finite queries, such as two noncommuting spin or polarization measurements, without assuming Hilbert space.  The Hilbert representation of such an IJC interface is a downstream representation theorem, not an input to this example.
\end{remark}

\begin{example}[Coupled but still Sep: shared classical bath]
Let two classical registers $A$ and $B$ be jointly coupled to a shared thermal bath.  Reading $A$ and $B$ may require an interface cost
\[
 \kappa_{\mathrm{int}}(A,B)>0
\]
because the bath must be stabilized or monitored.  Nevertheless, if there exists a joint classical record space
\[
 \Omega_{AB}=\Omega_A\times\Omega_B\times\Omega_{\mathrm{bath}}
\]
containing all relevant queried variables, then the interface is Sep.  Positive interface cost does not by itself certify IJC.
\end{example}

\section{IJC witnesses}\label{sec:ijc-witnesses}

A witness for IJC is any finite operational certificate that rules out all faithful commuting-extension defenders for the queried family.

\begin{remark}[Exact finite-table witnesses and noisy data]\label{rem:exact-noisy-witnesses}
All witnesses in this section are exact finite-table witnesses.  In noisy, finite-sample, or finite-resolution settings, the same ideas should be read through the $\varepsilon$-faithful defender definition: the tested word set, operational norm, noise model, and resolution threshold must be stated explicitly, and the witness certifies IJC only when every approximate defender fails at that specified tolerance.
\end{remark}


\begin{proposition}[Bell/CHSH witness]
If a bipartite queried interface has correlations violating a Bell/CHSH inequality satisfied by every common classical record model, then it is IJC.
\end{proposition}
\begin{proof}
A commuting-extension defender supplies a joint classical record model for all queried settings and outcomes.  In the nondisturbing Bell/CHSH case this is exactly the kind of joint distribution characterized by Fine's theorem; such a model satisfies the relevant Bell/CHSH inequalities.  Violation rules out the defender.  By the Sep/IJC classification, the interface is IJC.
\end{proof}

\begin{proposition}[Kochen--Specker or sheaf-contextual witness]
If queried contexts admit no global value assignment compatible with their context-wise records, then no commuting-extension defender exists.  Hence the interface is IJC.
\end{proposition}
\begin{proof}
A commuting-extension defender is exactly a global Boolean record completion for the queried distinctions.  Such a completion restricts to every measured context.  If no compatible global assignment exists, the defender cannot exist.
\end{proof}

\begin{proposition}[Controlled order witness]
Suppose queries $A,B$ have order-dependent continuation statistics that no classical invasive model satisfying the defender's common-record and recovery conditions can reproduce.  Then the interface is IJC.
\end{proposition}
\begin{proof}
The hypothesis says precisely that all commuting-extension defenders fail to reproduce the tested continuation-word behavior.  Thus the interface is not Sep; by classification it is IJC.
\end{proof}

\begin{proposition}[Complementarity witness]
If enforcing one distinction destroys an incompatible continuation profile and no common classical record can jointly preserve both profiles while reproducing the observed continuation behavior, then the queried interface is IJC.
\end{proposition}

\begin{remark}[Classical disturbance is not enough]
Order dependence alone is not sufficient for IJC, because ordinary classical invasive measurements can disturb later records.  The witness must rule out a faithful commuting-extension defender, not merely observe disturbance.
\end{remark}

\section{Continuum bridge: from finite interface ledgers to anchor profiles}\label{sec:recruitment-radius}

The results proved above are finite-interface results.  The Sep/IJC
classification, the continuation-word quotient, and the Boolean-record
representation theorem do not require a continuum substrate, a field
equation, a relaxation equation, or any assumption about Hilbert space.
This section is therefore not part of the proof of the Sep/IJC theorem.
It is a controlled bridge showing how the finite ledger language may be
represented, in appropriate coarse-grained regimes, by continuum anchor
profiles.

The guiding idea is simple.  A distinction enforced at an interface need
not draw all of its support from a single localized record.  In extended
systems, part of the burden of maintaining the distinction may be
distributed through substrate degrees of freedom.  When the relevant
substrate supports are numerous, individually below the observational
granularity, and stable under coarse-graining, it is useful to replace a
finite list of ledger supports by a profile
\[
    \phi_d(x)
\]
describing the density of enforcement commitment associated with
distinction $d$ at substrate location $x$.

This continuum profile is an effective representation of finite
enforcement bookkeeping.  It is not an additional primitive and it is
not needed for the finite representation theorem.  It is licensed only
in regimes where the finite supports admit a stable coarse-grained
limit, where the cost floor remains nonzero at the resolved scale, and
where the chosen substrate variables vary smoothly enough for a
continuum description to be meaningful.

\subsection{Regime assumptions for the continuum bridge}\label{subsec:continuum-bridge-regime}

The continuum bridge uses the following additional regime assumptions.
They are not primitive inputs to the finite Sep/IJC theorem.

\begin{description}
\item[C1.\ Coarse-grained substrate support.]
The finite enforcement supports of the distinction admit a stable
coarse-grained representation over a measurable substrate $\Sigma$.

\item[C2.\ Noncollapsed resolved floor.]
At the chosen resolution, the positive cost floor does not collapse:
resolved distinctions still carry positive finite enforcement burden.

\item[C3.\ Local substrate response.]
The cost of carrying an incremental enforcement load at location $x$
depends only on the local substrate state and on finite-range
correlations with nearby substrate regions.

\item[C4.\ Smooth small-load regime.]
The profile $\phi_d(x)$ lies in a perturbative regime in which the
effective cost functional can be expanded to second order.

\item[C5.\ Linear-response relaxation.]
For nonstationary anchors, deviations from instantaneous equilibrium
are small enough that first-order relaxation toward the instantaneous
minimizer is a valid approximation.
\end{description}

When these assumptions fail, the continuum profile language may fail,
but the finite-interface results of Sections~1--13 remain unaffected.

\subsection{Effective recruitment cost functional}\label{subsec:effective-recruitment-functional}

Under C1--C4, the distributed enforcement burden for distinction $d$
may be represented by an effective recruitment functional
\begin{equation}\label{eq:Erec}
E_{\mathrm{rec}}[d, \phi_d, S_{\mathrm{subst}}]
\;=\;
\int_\Sigma \phi_d(x)\, \varepsilon_{\mathrm{local}}(x, d)\, dx
\;+\;
\int_\Sigma\!\!\int_\Sigma \phi_d(x)\,
I_{\mathrm{int}}(x, x'; d, S_{\mathrm{subst}})\,
\phi_d(x')\, dx\, dx'.
\end{equation}

Here $\varepsilon_{\mathrm{local}}(x, d)$ is the local cost density for
supporting distinction $d$ at substrate location $x$, and
$I_{\mathrm{int}}$ is an effective finite-range interaction kernel
encoding how enforcement support at one substrate location changes the
cost of support elsewhere.

This functional should be read as the second-order effective form of the
finite ledger in a smooth small-load regime.  It is not asserted as a
universal microscopic action.  Higher-order terms, nonlocal memory
effects, saturation, and far-from-equilibrium substrate reorganization
are outside the approximation.

\subsection{Interpretive bridge to field language}\label{subsec:field-interpretive-bridge}

In regimes where the continuum bridge applies, the equilibrium profile
\[
    \phi_d^{\mathrm{eq}}
    \;=\;
    \arg\min_{\phi_d} E_{\mathrm{rec}}[d, \phi_d, S_{\mathrm{subst}}]
\]
may be interpreted as the substrate recruitment region associated with
the enforced distinction $d$.  This gives a useful translation between
APF ledger language and ordinary field language: what standard field
descriptions represent as a spatial field profile may, in this
effective regime, be read as the spatial distribution of enforcement
support required to maintain a distinction.

This is an interpretive bridge, not a replacement for the finite
Sep/IJC theorem.  Coulomb, screened, Yukawa, dielectric, or
Casimir-like profiles may be modeled by different choices of substrate
kernel and boundary data.  The present supplement does not derive those
standard field equations.  It only identifies the APF object that would
play the corresponding role: the equilibrium recruitment profile of an
enforced distinction in a specified substrate.

\subsection{Optional linear-response dynamics}\label{subsec:linear-response-dynamics}

If, in addition to C1--C4, the substrate satisfies the linear-response
condition C5, one may model the actual profile $\phi(x, t)$ as relaxing
toward the instantaneous equilibrium profile associated with the
current anchor state:
\begin{equation}\label{eq:relaxation}
\frac{\partial \phi(x, t)}{\partial t}
\;=\;
-\,\frac{\phi(x, t) - \phi^{\mathrm{eq}}[s_{\mathrm{anchor}}(t)](x)}
{\tau_{\mathrm{rec}}(x)}.
\end{equation}

Define the mismatch field
\[
    \mu(x, t)
    \;=\;
    \phi(x, t) - \phi^{\mathrm{eq}}[s_{\mathrm{anchor}}(t)](x).
\]
Then, to second order around equilibrium, the excess recruitment cost
density may be written
\[
    \delta E(x, t) \;=\; \tfrac{1}{2}\, K(x)\, \mu(x, t)^2,
\]
where $K(x)$ is the local cost-curvature of the effective functional.

In downstream dynamical applications, this excess cost can be
interpreted as the budget available for substrate excitation when the
anchor changes too rapidly for adiabatic recruitment.  A simple
dimensional estimate of the local excitation-production rate is
\[
    \left.\frac{dN}{dt}\right|_x
    \;\sim\;
    \frac{K(x)\, \mu(x, t)^2}
    {2\, \tau_{\mathrm{rec}}(x)\, \hbar\, \omega_{\mathrm{typ}}(x)}.
\]
The proportionality and mode-counting details depend on the substrate
model.  This expression is therefore not used as a theorem in the
present supplement; it is included only to show how the ledger
formalism connects to later radiation and relaxation models.

\subsection{Upper bound on the interface term in the continuum-bridge regime}\label{subsec:kappa-int-upper-bound}

The structural lower bound on $\kappa_{\Gamma,\mathrm{int}}$ from \S\ref{sec:composition-factorization} (Theorem~\ref{thm:kappa-int-singleton-shape}, Corollary~\ref{cor:kappa-int-binary-shape}) is unconditional: it follows from MD via BW and holds in any finite physical regime.  An upper bound, by contrast, requires more.  In this subsection we use the continuum-bridge apparatus of \S\S\ref{subsec:effective-recruitment-functional}--\ref{subsec:linear-response-dynamics} to derive a regime-conditional upper bound, valid wherever C1--C5 apply.  The bound is internal to this supplement and does not depend on Papers~24 or~25, although those papers extend the apparatus to regimes beyond C1--C5.

\subsubsection*{Multi-distinction extension of \texorpdfstring{$E_{\mathrm{rec}}$}{Erec}}

The recruitment functional \eqref{eq:Erec} was defined for a single distinction $d$ with profile $\phi_d$.  For an admissible family $S = \{d_1, \ldots, d_n\}$ in the same substrate context $S_{\mathrm{subst}}$, the natural multi-distinction extension is
\begin{equation}\label{eq:Erec-multi}
E_{\mathrm{rec}}^{\mathrm{multi}}\bigl[S, \{\phi_{d_i}\}, S_{\mathrm{subst}}\bigr]
\;=\;
\sum_{i=1}^n E_{\mathrm{rec}}\bigl[d_i, \phi_{d_i}, S_{\mathrm{subst}}\bigr]
\;+\;
\sum_{i\neq j} E_{\mathrm{cross}}\bigl[d_i, d_j; \phi_{d_i}, \phi_{d_j}, S_{\mathrm{subst}}\bigr],
\end{equation}
where the cross-coupling functional is
\begin{equation}\label{eq:Ecross}
E_{\mathrm{cross}}\bigl[d_i, d_j; \phi_{d_i}, \phi_{d_j}, S_{\mathrm{subst}}\bigr]
\;:=\;
\int_\Sigma\!\!\int_\Sigma \phi_{d_i}(x)\,
I_{\mathrm{int}}^{d_i, d_j}(x, x'; S_{\mathrm{subst}})\,
\phi_{d_j}(x') \, dx \, dx'.
\end{equation}
The single-distinction kernel $I_{\mathrm{int}}(x, x'; d, S_{\mathrm{subst}})$ of \eqref{eq:Erec} extends to a cross-distinction kernel $I_{\mathrm{int}}^{d_i, d_j}(x, x'; S_{\mathrm{subst}})$ of the same finite-range form: the cooperative cost between enforcement support for $d_i$ at $x$ and enforcement support for $d_j$ at $x'$ is local in the substrate.  C3 (local substrate response) is what licenses this extension; the same locality argument that bounded the self-interaction kernel bounds the cross-interaction kernel.

In this representation, the residue $\kappa_{\Gamma,\mathrm{int}}$ of \S\ref{sec:composition-factorization} reads as the cross-coupling sum: $\kappa_{\Gamma,\mathrm{int}}(S) \approx \sum_{i\neq j} E_{\mathrm{cross}}[d_i, d_j]$, modulo higher-order corrections suppressed by the C4 small-load assumption.

\subsubsection*{Generic upper bound from kernel-norm finiteness}

\begin{theorem}[Upper bound on the binary interface term in the continuum-bridge regime]\label{thm:kappa-int-binary-upper-bound}
Let $S_1, S_2$ be admissible families in the composite context $\Gamma_{12}$ realized by continuum-bridge profiles $\phi_{S_1}, \phi_{S_2}$ on substrate $\Sigma$ satisfying C1--C5, and let $I_{\mathrm{int}}^{S_1, S_2}$ denote the cooperative-cost kernel between $S_1$ and $S_2$ in the substrate context $S_{\mathrm{subst}}$.  Define the positive-part kernel sup
\[
    I_{\mathrm{int}}^+ \;:=\; \sup_{x, x'\in\Sigma} \max\bigl\{0,\; I_{\mathrm{int}}^{S_1, S_2}(x, x'; S_{\mathrm{subst}})\bigr\}.
\]
Then $I_{\mathrm{int}}^+ < \infty$ by C3, and
\[
    \kappa_{\mathrm{int}}(S_1, S_2)
    \;\le\;
    I_{\mathrm{int}}^+ \cdot \|\phi_{S_1}\|_{L^1(\Sigma)} \cdot \|\phi_{S_2}\|_{L^1(\Sigma)}.
\]
For normalized profiles ($\|\phi_{S_i}\|_{L^1} = 1$), $\kappa_{\mathrm{int}}(S_1, S_2) \le I_{\mathrm{int}}^+$.
\end{theorem}
\begin{proof}
By the multi-distinction extension \eqref{eq:Erec-multi}--\eqref{eq:Ecross}, the cross-interaction contribution to the joint cost is
\[
    \int_\Sigma\!\!\int_\Sigma \phi_{S_1}(x)\,
    I_{\mathrm{int}}^{S_1, S_2}(x, x'; S_{\mathrm{subst}})\,
    \phi_{S_2}(x') \, dx \, dx'.
\]
Using $I_{\mathrm{int}}^{S_1, S_2}(x, x') \le I_{\mathrm{int}}^+$ pointwise, with $\phi_{S_1}, \phi_{S_2} \ge 0$:
\[
    \le \;
    I_{\mathrm{int}}^+ \int_\Sigma\!\!\int_\Sigma \phi_{S_1}(x)\,\phi_{S_2}(x') \, dx \, dx'
    \;=\;
    I_{\mathrm{int}}^+ \cdot \|\phi_{S_1}\|_{L^1} \cdot \|\phi_{S_2}\|_{L^1}.
\]
The integrals factor by Fubini.  Profile $L^1$ norms are finite by C1 (admissible coarse-grained profiles have finite total recruitment mass), and the kernel sup $I_{\mathrm{int}}^+$ is finite by C3 (local substrate response gives bounded cooperative cost in any finite neighborhood, with finite-range cutoff outside).
\end{proof}
\coderef{check\_T\_kappa\_int\_upper\_bound\_C1C5}{kappa\_int\_bounds.py}

\subsubsection*{Far-separation exponential suppression}

\begin{corollary}[Exponential suppression of the interface term for separated supports]\label{cor:kappa-int-far-separation}
Suppose, in addition to C1--C5, the substrate kernel exhibits a finite correlation length $\xi_{\mathrm{rec}}$ in the sense that
\[
    \bigl|I_{\mathrm{int}}^{S_1, S_2}(x, x'; S_{\mathrm{subst}})\bigr|
    \;\le\;
    I_0 \, e^{-|x - x'|/\xi_{\mathrm{rec}}}
\]
for some bounded $I_0 < \infty$ and a metric $|\cdot|$ on $\Sigma$.  If the supports of $\phi_{S_1}$ and $\phi_{S_2}$ are separated by distance at least $L$,
\[
    \mathrm{dist}\bigl(\mathrm{supp}(\phi_{S_1}),\,\mathrm{supp}(\phi_{S_2})\bigr)
    \;:=\;
    \inf\bigl\{|x - x'| : x\in\mathrm{supp}(\phi_{S_1}),\, x'\in\mathrm{supp}(\phi_{S_2})\bigr\}
    \;\ge\; L,
\]
then
\[
    \bigl|\kappa_{\mathrm{int}}(S_1, S_2)\bigr|
    \;\le\;
    I_0 \cdot e^{-L/\xi_{\mathrm{rec}}} \cdot \|\phi_{S_1}\|_{L^1} \cdot \|\phi_{S_2}\|_{L^1}.
\]
\end{corollary}
\begin{proof}
For any $x \in \mathrm{supp}(\phi_{S_1})$ and $x' \in \mathrm{supp}(\phi_{S_2})$ the separation hypothesis gives $|x - x'| \ge L$, hence $|I_{\mathrm{int}}^{S_1, S_2}(x, x')| \le I_0 \, e^{-L/\xi_{\mathrm{rec}}}$.  Apply this pointwise bound under the integral as in Theorem~\ref{thm:kappa-int-binary-upper-bound}.  The integrand vanishes outside the support of either profile, so the integral runs over points satisfying the separation hypothesis.
\end{proof}

\subsubsection*{Singleton-form upper bound}

\begin{corollary}[Upper bound on the singleton-form interface term in the continuum-bridge regime]\label{cor:kappa-int-singleton-upper-bound}
Let $S = \{d_1, \ldots, d_n\}$ be admissible with continuum-bridge profiles $\phi_{d_1}, \ldots, \phi_{d_n}$ on substrate $\Sigma$ satisfying C1--C5, and let
\[
    I_{\mathrm{int}}^+ \;:=\; \max_{i \neq j}\;\sup_{x, x'\in\Sigma}\, \max\bigl\{0,\; I_{\mathrm{int}}^{d_i, d_j}(x, x'; S_{\mathrm{subst}})\bigr\}
\]
denote the supremum of the positive part of any cross-distinction kernel.  Then
\[
    \kappa_{\Gamma,\mathrm{int}}(S)
    \;\le\;
    I_{\mathrm{int}}^+ \cdot \!\!\sum_{i \neq j} \!\!\|\phi_{d_i}\|_{L^1}\,\|\phi_{d_j}\|_{L^1}
    \;\le\;
    I_{\mathrm{int}}^+ \cdot n(n-1) \cdot \max_i \|\phi_{d_i}\|_{L^1}^2.
\]
For normalized profiles ($\|\phi_{d_i}\|_{L^1} = 1$ for all $i$), $\kappa_{\Gamma,\mathrm{int}}(S) \le I_{\mathrm{int}}^+ \cdot n(n-1)$.
\end{corollary}
\begin{proof}
By the multi-distinction extension \eqref{eq:Erec-multi}--\eqref{eq:Ecross}, the singleton-form residue reads $\kappa_{\Gamma,\mathrm{int}}(S) = \sum_{i\neq j} E_{\mathrm{cross}}[d_i, d_j; \phi_{d_i}, \phi_{d_j}]$ modulo C4-suppressed higher-order terms.  Apply Theorem~\ref{thm:kappa-int-binary-upper-bound} to each cross-coupling term: $E_{\mathrm{cross}}[d_i, d_j] \le I_{\mathrm{int}}^+ \|\phi_{d_i}\|_{L^1} \|\phi_{d_j}\|_{L^1}$ (with the kernel sup restricted to the $(d_i, d_j)$ pair, dominated by the family-wide $I_{\mathrm{int}}^+$).  Sum over $i \neq j$ to get the first inequality; bound each $L^1$ norm by $\max_i\|\phi_{d_i}\|_{L^1}$ for the second.
\end{proof}

\subsubsection*{Two-sided structural rigidity}

The lower bound from \S\ref{sec:composition-factorization} (unconditional, derived from MD via BW) and the upper bound from this section (conditional on C1--C5, derived from substrate locality and finite-range coupling) combine into a single two-sided structural rigidity statement.

\begin{theorem}[Two-sided structural bound on the interface term]\label{thm:kappa-int-two-sided-bound}
Let $S = \{d_1, \ldots, d_n\}$ be an admissible family in a finite physical regime satisfying the continuum-bridge regime assumptions C1--C5, with profiles $\phi_{d_1}, \ldots, \phi_{d_n}$ on substrate $\Sigma$.  Then the singleton-form interface term satisfies
\[
    \underbrace{n\,\varepsilon^*_\Gamma \;-\; \sum_{d\in S}\kappa_\Gamma(d)}_{\substack{\text{structural lower bound,}\\ \text{from MD via BW}}}
    \;\;\le\;\;
    \kappa_{\Gamma,\mathrm{int}}(S)
    \;\;\le\;\;
    \underbrace{I_{\mathrm{int}}^+ \cdot \!\!\sum_{i \neq j} \!\!\|\phi_{d_i}\|_{L^1}\,\|\phi_{d_j}\|_{L^1}}_{\substack{\text{kernel-norm upper bound,}\\ \text{from C3 + C1}}}.
\]
The lower bound is unconditional and holds in any finite physical regime; the upper bound is conditional on the continuum-bridge regime C1--C5.  In any regime where both apply, $\kappa_{\Gamma,\mathrm{int}}$ is bounded between explicit structurally derived quantities and is not a free functional.
\end{theorem}
\begin{proof}
The lower bound is Theorem~\ref{thm:kappa-int-singleton-shape}.  The upper bound is Corollary~\ref{cor:kappa-int-singleton-upper-bound}.
\end{proof}
\coderef{check\_T\_kappa\_int\_two\_sided\_rigidity}{kappa\_int\_bounds.py}

\paragraph{Sharpness.}
The lower bound is attained in the full-amortization limit: when every per-distinction surplus $\kappa_\Gamma(d_i) - \varepsilon^*_\Gamma$ is fully paid by shared substrate and the joint cost reduces to $n\varepsilon^*_\Gamma$.  The upper bound is attained when every cross-distinction kernel uniformly saturates its positive-part sup on the joint support and no kernel cancellation occurs across the family.  In the far-separation limit (Corollary~\ref{cor:kappa-int-far-separation}), the bound is tight when the kernel saturates the exponential decay envelope on the inter-support distance.  Real APF interfaces typically lie strictly between these extremes, but the two-sided bound certifies that the residue cannot escape the structural envelope determined by the spine plus the substrate's coupling kernel.

\subsubsection*{Reading}

Theorem~\ref{thm:kappa-int-two-sided-bound} packages the structural picture in one statement: in the continuum-bridge regime, $\kappa_{\Gamma,\mathrm{int}}$ is bounded between $n\varepsilon^*_\Gamma - \sum_d \kappa_\Gamma(d)$ from below (unconditional, MD via BW) and $I_{\mathrm{int}}^+ \cdot \sum_{i\neq j} \|\phi_{d_i}\|_{L^1}\|\phi_{d_j}\|_{L^1}$ from above (conditional on C1--C5).  In the typical APF interface --- finite cost floor, local substrate response, smooth small-load profiles --- the residue is a structurally constrained finite quantity, not a free functional.  The audit-flagged complaint that $\kappa_{\Gamma,\mathrm{int}}$ ``absorbs whatever deviation from additivity the joint protocol exhibits'' (v8.24, \S\ref{sec:distinctions-cost} Remark~\ref{rem:kappa-int-placeholder}) is closed: in any regime where both bounds apply, the residue's shape is determined by the spine plus the substrate's coupling kernel, with no remaining structural freedom.  Far-separated anchors get exponential suppression of the coupling overhead (Corollary~\ref{cor:kappa-int-far-separation}), recovering the factorization limit of Theorem~\ref{thm:ledger-additivity-protocol-factorization} as the formal $L \to \infty$ limit.

The continuum-bridge upper bound is conditional on C1--C5.  Outside this regime --- saturation regimes where the second-order expansion C4 fails and higher-order terms in $E_{\mathrm{rec}}$ dominate, far-from-equilibrium regimes where the linear-response approximation C5 fails (Paper~25 \S8 cosmogenesis), or substrate-collapse regimes where the locality assumption C3 fails --- the upper bound does not transfer automatically.  Paper~25 \S\S\,2--3 supplies the recruitment-radius extension's full apparatus for substrate-rich regimes that exit C1--C5, and the Paper~6 supplement \S1.5 closes the H1/H2 hypotheses for the geometric substrate.  Downstream papers reasoning about coupling overhead in those regimes should supply the relevant upper bound from substrate-level commitments rather than relying on the continuum-bridge form alone.

\bigskip

\begin{center}
\fbox{\begin{minipage}{0.92\linewidth}
\textbf{Scope of Section~\ref{sec:recruitment-radius}.}
Sections~1--13 prove the finite-interface continuation and Sep/IJC
results.  Section~\ref{sec:recruitment-radius} does not add a premise
to those proofs and is not needed by them.  It gives an effective
continuum representation of the same ledger idea under additional
coarse-graining, smoothness, locality, and linear-response assumptions.
Claims about specific field equations, radiation laws, gauge sectors,
gravity, or cosmological residuals belong to Papers~24 and~25
(\emph{The Recruitment-Radius Extension to APF, Paper~I ---
Foundations} and \emph{Paper~II --- Applications, Closures, and
Domain Boundaries}) and to their supplements, and require their own
model-specific hypotheses.  Paper~24 develops the recruitment cost
functional, the master equation for radiation production, the
sixteen-case unification, and the DCE $Q$-dependence prediction.
Paper~25 closes Paper~6 H1/H2 and Paper~8~\S J via the master
equation, gives the geometric-substrate proposal at linearized and
nonlinear levels, and names the framework's two domain boundaries
(Newton's constant; cosmogenesis far-from-equilibrium dynamics).
The bridge in Section~\ref{sec:recruitment-radius} is the link
between Paper~1's finite-interface formalism and the substantive
recruitment-radius work in Papers~24 and~25; it is not a substitute
for either.
\end{minipage}}
\end{center}

\section{Relation to GPT, contextuality, and sequential frameworks}\label{sec:related-frameworks}

\paragraph{Terminology crosswalk.}
For readers familiar with neighboring frameworks, the following
compact mapping records each APF term, its closest standard term, and
the substantive difference.

\begin{center}
\begin{tabular}{p{0.27\linewidth}p{0.30\linewidth}p{0.33\linewidth}}
\toprule
\textbf{APF term} & \textbf{Closest standard term} & \textbf{Difference} \\
\midrule
Continuation identity & Operational equivalence & Defined by finite admissible continuations. \\
Physical descent & Quotient-respecting observable & Filters/predicates must factor through continuation quotient. \\
Defender & Classical/noncontextual completion & Must preserve tested sequential words and recovery maps. \\
Sep & Classical completability & Relative to finite tested word set and operational resolution. \\
IJC & Failure of classical/global completion & Failure of every faithful commuting continuation defender. \\
Ledger ($\kappa_\Gamma$) & Resource/accounting structure & Tracks enforcement cost and trigger conditions. \\
No-added-resource filter & Free/passive post-processing & Formalized as resource-nonincreasing protocol simulation. \\
\bottomrule
\end{tabular}
\end{center}


We use APF terminology throughout, but the construction sits near several established operational frameworks.

\begin{center}
\begin{longtable}{@{}p{0.25\linewidth}p{0.30\linewidth}p{0.35\linewidth}@{}}
\toprule
APF notion & Neighboring standard notion & Difference \\
\midrule
Admissible possibility space & Operational theory / GPT state space & APF starts from continuations, then derives quotient states and convexity. \\
Physical distinction & Effect, proposition, event & APF requires finite positive enforcement cost. \\
Continuation profile & Operational equivalence class & APF includes future tests, records, composition, and maintenance operations. \\
Sep defender & Noncontextual hidden-variable model / joint distribution & APF defender must reproduce sequential continuation words, not only static marginals.  In Bell/CHSH nondisturbing cases, this reduces to Fine's joint-distribution criterion. \\
IJC & Contextuality, incompatibility, order effects & APF frames nonclassicality as failure of faithful commuting continuation completion. \\
Continuation-word algebra & Sequential instrument algebra / process algebra & APF quotients words by context-complete continuation identity. \\
Ledger capacity & Resource monotone / distinguishability cost & APF treats enforceability as primitive finite support cost. \\
Joint measurability/compatibility & Common jointly measurable observable or commuting dilation & Sep defender is a finite common-record completion for the tested interface, with sequential recovery data. \\
Effect algebras and effectus theory & Partial event structures and operational maps & APF adds continuation identity and positive enforcement cost before convex/probabilistic completion. \\
\bottomrule
\end{longtable}
\end{center}


\subsection*{Generalized probabilistic theories: full finite bridge}
\label{subsec:gpt-full-bridge}

Generalized probabilistic theories (GPTs) usually begin with an ordered real vector space \(V\), a positive cone \(V_+\), an order unit \(u\), a normalized state space
\[
  \Omega=\{x\in V_+ : u(x)=1\},
\]
and an effect set \(E\subseteq [0,u]\subseteq V^*\), with probabilities given by \(e(x)\).  Classical finite probability theory is the special case in which \(\Omega\) is a simplex \(\Delta(\Lambda)\), equivalently \(V\cong\R^\Lambda\) with pointwise order.  This is the standard GPT starting point \cite{Hardy2001,Barrett2007,ChiribellaPurification2010,Chiribella2011}.  APF reaches this structure only after a finite interface, resolution, quotient, and admissible classical-control records have been specified.

\begin{definition}[APF-induced accessible GPT fragment]\label{def:apf-induced-gpt-fragment}
Restrict APF to a finite static prepare--measure fragment with finite resolved preparations \(S\), finite resolved effects \(E\), no tested update words, and exact or \(\delta\)-resolved probabilities \(p(e|s)\).  Assume the context includes admissible classical randomization records for preparations.  Define the tagged evaluation vector
\[
  \widehat\sigma_s := \bigl(1,(p(e|s))_{e\in E}\bigr)\in\R^{1+|E|}.
\]
Let
\[
  V_{\Gamma,Q}^{\rm GPT}:=\operatorname{span}\{\widehat\sigma_s:s\in S\}\subseteq\R^{1+|E|},
  \qquad
  u(t,x):=t,
\]
and let \(\Omega_{\Gamma,Q}^{\rm acc}\) be the convex hull of the vectors \(\widehat\sigma_s\) generated by admissible classical control.  The accessible state cone is
\[
  C_{\Gamma,Q}^{\rm acc}:=\operatorname{cone}\bigl(\Omega_{\Gamma,Q}^{\rm acc}\bigr).
\]
Each tested effect \(e\in E\) determines the coordinate functional
\[
  \widehat e(t,x):=x_e,
\]
so that \(0\le \widehat e\le u\) on the accessible cone.  Preparations and effects are then quotient-identified exactly when their finite evaluation profiles agree at the chosen operational resolution.  The tuple
\[
  \mathsf G_{\Gamma,Q}^{\rm acc}:=\bigl(V_{\Gamma,Q}^{\rm GPT},C_{\Gamma,Q}^{\rm acc},u,\Omega_{\Gamma,Q}^{\rm acc},\{\widehat e:e\in E\}\bigr)
\]
is the APF-induced accessible GPT fragment.
\end{definition}

This construction is an \emph{accessible fragment}, not a no-restriction GPT.  APF does not assert that every positive functional in \([0,u]\) is physically implementable, nor that every point of an ambient cone is preparable.  Only those preparations, effects, mixtures, and quotients generated by the finite interrogation and its admissible classical controls are included.

\begin{proposition}[Static APF--GPT equivalence]\label{prop:static-apf-gpt-equivalence}
For every finite static APF prepare--measure interface with admissible classical randomization, Definition~\ref{def:apf-induced-gpt-fragment} gives a finite accessible GPT fragment preserving the tested probabilities and quotient operational equivalences.  Conversely, any finite accessible GPT fragment with finite exposed preparations/effects, a resolution, and enforcement-cost assignments for resolved separators can be presented as a static finite APF interface whose continuation identity is equality of exposed state--effect evaluations at that resolution.
\end{proposition}

\begin{proof}
The forward construction sends each resolved preparation to its finite probability vector and each resolved effect to the corresponding coordinate functional.  Classical randomization records give convex mixtures, and quotient descent identifies exactly those preparations or effects that no exposed evaluation separates at the stated resolution.  Thus all tested probabilities are preserved.  Conversely, a finite accessible GPT fragment already specifies finite operational probabilities for a chosen preparation/effect family.  Declare two raw labels continuation-identical when every exposed effect gives the same value within the stated resolution, and declare a resolved physical distinction when some exposed effect separates the corresponding evaluation profiles.  Supplying a positive cost assignment for such separators completes the APF static interface.  The construction preserves the finite fragment, not any unexposed global GPT completion.
\end{proof}

\begin{proposition}[Sep equals simplex embeddability of the induced fragment]\label{prop:sep-simplex-embeddability}
In a finite static nondisturbing APF prepare--measure fragment, the following are equivalent:
\begin{enumerate}[label=(\roman*)]
\item the APF interface is Sep;
\item the APF-induced accessible GPT fragment \(\mathsf G_{\Gamma,Q}^{\rm acc}\) admits a finite classical simplex embedding;
\item there exist a finite set \(\Lambda\), probability distributions \(\mu_s\in\Delta(\Lambda)\), and response functions \(r_e:\Lambda\to[0,1]\) such that
\[
  p(e|s)=\sum_{\lambda\in\Lambda}\mu_s(\lambda)r_e(\lambda)
\]
for every tested preparation/effect pair, with the represented effects separating the resolved quotient exactly as required by faithfulness.
\end{enumerate}
Equivalently, Sep is simplex embeddability of the accessible GPT fragment, with quotient faithfulness added as the APF separation requirement.
\end{proposition}

\begin{proof}
A Sep defender is a finite Boolean record representation.  Its atoms \(\Lambda\) form a classical simplex: preparations are probability weights on atoms and tested effects are pointwise response functions.  This gives (iii), hence a positive order-unit-preserving embedding into the finite classical GPT \(\R^\Lambda\).  Conversely, a simplex embedding supplies exactly such atoms, weights, and response functions; the Boolean algebra of subsets of \(\Lambda\), together with the coordinate functions \(r_e\), is a commutative record representation.  Faithfulness is precisely the requirement that the embedding not identify operationally separated quotient classes.  Thus the APF defender criterion and simplex embeddability coincide in the finite static nondisturbing regime.
\end{proof}

\begin{remark}[Cone equivalence reading]
Accessible-fragment cone equivalence compares finite operational fragments up to positive order-isomorphisms preserving the exposed probabilities.  Under Proposition~\ref{prop:static-apf-gpt-equivalence}, APF quotient-equivalent static fragments induce the same accessible cone data up to such an isomorphism, while Proposition~\ref{prop:sep-simplex-embeddability} identifies the Sep branch with the case in which that accessible cone embeds into a classical simplex.  APF adds two pieces not normally present in bare cone equivalence: the continuation quotient that defines the fragment, and the enforcement ledger \(\kappa_\Gamma\) that charges the maintenance of its resolved separators \cite{SelbySchmidWolfeSainzKunjwalSpekkens2023}.
\end{remark}

\begin{remark}[Where APF is not just GPT in different language]\label{rem:apf-not-just-gpt}
The bridge is exact only for the finite static fragment just described.  APF differs from a standard GPT presentation in six ways.
\begin{enumerate}[label=(\alph*)]
\item APF begins with admissible continuations and quotient identity, not with a pre-given convex state space.
\item Convexity enters only when the context includes enforceable classical control records implementing randomization.
\item APF assumes no no-restriction hypothesis: untested positive functionals need not be physical effects.
\item APF assumes no tensor product, local tomography, purification, or composite rule at this layer; those are additional hypotheses in GPT reconstructions and are deferred here.
\item Sequential APF interfaces replace a single static effect set with a typed partial continuation-word algebra.  The static GPT fragment is recovered by restricting to length-one prepare--measure words.
\item The enforcement ledger is not part of ordinary GPT kinematics.  It records the positive cost of keeping quotient distinctions physically maintained.
\end{enumerate}
Thus GPT supplies the right comparison language for static accessible fragments, while APF supplies the pre-GPT quotient, cost, and sequential-continuation machinery.
\end{remark}

\subsection*{Quantum logic and effect algebras}
The distinction-refinement structure resembles partial event/effect structures.  APF differs by making Booleanity branch-specific: in the Sep branch the queried distinctions embed into a Boolean algebra; in the IJC branch they do not.

\subsection*{Fine-style joint distributions}
For Bell/CHSH tables in the nondisturbing marginal regime, a Sep defender reduces to a joint distribution over all settings and outcomes.  Thus the APF defender criterion agrees with Fine's theorem in its static Bell form, while the continuation-word formulation extends the criterion to sequential and disturbance-sensitive protocols.

\subsection*{Joint measurability and commuting dilations}
For static measurement families, a Sep defender resembles a jointly measurable refinement or a commuting dilation: all queried events are represented inside one common classical record algebra.  APF is stricter about the finite interface being defended: the defender must also reproduce the tested sequential words, recovery maps, and continuation equivalences specified by $\mathcal T$.  Thus the comparison is exact for static compatibility tests and becomes a sequential commuting-extension criterion when update behavior is included.

\subsection*{Spekkens-style noncontextuality}

A Sep defender is close in spirit to a generalized noncontextual ontological model, but with a stronger sequential requirement.  It must reproduce continuation-word behavior, not merely preparation and measurement marginals.

\begin{proposition}[Finite static Spekkens equivalence]
\label{prop:finite-static-spekkens-equivalence}
Restrict APF to a finite static prepare--measure interface with no tested sequential update words, quotient preparations/effects by operational equivalence at the chosen resolution, and require exact reproduction of all tested probabilities.  Then an APF Sep defender is equivalent to a finite generalized-noncontextual ontological model for that tested fragment: the defender atoms are ontic states, preparations are probability distributions on those atoms, and effects are response functions depending only on the quotient effect class.
\end{proposition}
\begin{proof}
Given a defender \((\Lambda,\mathcal B,\rho)\), set the ontic state space to \(\Lambda\), represent each preparation by the probability weights \(p_s\), and represent each effect by the function \(\rho(e):\Lambda\to[0,1]\).  Quotient faithfulness ensures operationally equivalent preparations/effects receive the same representation and operationally distinct resolved effects remain separated.  Conversely, a finite noncontextual ontological model with response functions on a common ontic space defines a commutative record representation in \(\R^\Lambda\).  The equivalence is limited to the stated finite static fragment; APF adds further recovery and word-composition requirements when transformations or sequential words are included.
\end{proof}

\paragraph{Generalized noncontextuality and continuation quotienting.}
Spekkens-style generalized noncontextuality requires operationally
equivalent preparations, transformations, and effects to receive
identical ontological representations.  APF implements the same
discipline through continuation quotienting: raw labels with the
same finite admissible continuation profile are identified before
physical distinctions or filters are defined.  In the static
fragment, an APF defender is therefore a noncontextual ontological
model for the resolved finite experiment.  The APF extension is
that operational equivalence is organized by finite admissible
continuation identity, and the defender must also preserve tested
sequential words, recovery maps, and resource-trigger restrictions.
APF does not replace generalized noncontextuality; it gives a
continuation-algebraic and ledger-aware finite-interface version
of it.

\subsection*{Sheaf-theoretic contextuality}
The commuting-extension defender corresponds to a global compatible record assignment.  IJC corresponds to obstruction to such a global section, enriched by sequential continuation behavior.

\begin{proposition}[Finite static sheaf/global-section equivalence]
\label{prop:finite-static-sheaf-equivalence}
For a finite nondisturbing measurement scenario whose APF tested words are only context-wise joint readouts, existence of an APF defender is equivalent to existence of a global probability distribution over total outcome assignments whose marginals reproduce the context-wise empirical distributions.  In deterministic support language, this is the sheaf-theoretic global-section condition.
\end{proposition}
\begin{proof}
A defender supplies a finite Boolean atom \(\lambda\) carrying a value for every queried event in one common record algebra.  A state assigns a probability distribution \(p(\lambda)\) over those atoms.  Marginalizing \(p\) to each measured context gives the observed context-wise distribution.  Conversely, a global distribution over total assignments defines \(\Lambda\) as the set of total assignments and represents each queried event by the corresponding coordinate function.  This is exactly a Boolean defender for the static nondisturbing APF fragment.  Sequential APF interfaces add tested products and recovery maps, so the global-section comparison remains a static special case rather than the full theorem.
\end{proof}

\paragraph{Static sheaf fragment.}
Given a measurement scenario in the sense of sheaf-theoretic
contextuality, let $X$ be the finite set of measurements, let
$\mathcal M$ be the family of jointly queried contexts, and let $O_x$
be the finite outcome set for measurement $x$.  The APF static
fragment is obtained by taking $Q=X$, taking $\mathcal T$ to contain
only the context words $C\in\mathcal M$, and taking recovery maps to
be the usual marginal restrictions.  A finite Boolean defender then
consists of a common atom space
\[
    \Lambda \subseteq \prod_{x\in X} O_x
\]
with probability weights $p_s\in\Delta(\Lambda)$ whose restrictions to
each context reproduce the empirical model.  Thus, in this fragment,
APF Sep is exactly the existence of a global section or global
probability model, and APF IJC is sheaf contextuality.

The APF theorem extends this static fragment by allowing $\mathcal T$
to contain typed sequential words, tested order effects, recovery
operations, coarse-grainings, abort/failure words, and continuation
equivalence relations.  The sheaf-theoretic obstruction is therefore
recovered as the zero-depth or nondisturbing marginal case of a
broader finite continuation representation problem.

\begin{remark}[Exact and approximate noncontextual bridges]\label{rem:exact-approx-noncontextual-bridges}
Propositions~\ref{prop:finite-static-spekkens-equivalence} and~\ref{prop:finite-static-sheaf-equivalence} are exact finite-static equivalences.  Approximate \(\varepsilon\)-versions require choosing an operational norm, tolerance, nondisturbance convention, and finite tested word set.  Those choices are empirical resolution data, not part of the exact Sep/IJC representation theorem.
\end{remark}

\subsection*{Sequential instruments and order algebras}
The continuation-word algebra parallels sequential products of instruments.  The APF-specific step is quotienting by continuation identity: two words are identical exactly when no admissible context can distinguish their continuation behavior.

\paragraph{Temporal and device-independent sequential tests.}
Temporal Bell, Leggett--Garg, sequential contextuality, and
device-independent temporal tests fit naturally into APF as finite
tested word families.  A temporal test does not merely ask whether a
set of outcomes has a common static assignment.  It asks whether a
single classical record can reproduce the statistics of ordered
interrogation words such as
\[
    E_{i_k}\cdots E_{i_2} E_{i_1},
\]
together with the declared compatibility, noninvasiveness, recovery,
or coarse-graining assumptions of the protocol.

In APF language, such a protocol specifies a finite tested set
$\mathcal T$ of sequential continuation words.  A classical temporal
model is precisely a commutative record defender for that tested word
set: one finite Boolean atom space must reproduce all tested temporal
state--effect pairings while preserving the tested products and
recovery maps.  A temporal or device-independent inequality violation
is therefore a separating witness showing that no such commutative
defender exists for the specified $(\Gamma,Q,\mathcal T,\delta)$.

This formulation also makes explicit where assumptions enter.
Conditions often described as noninvasive measurability,
arrow-of-time constraints, instrument compatibility, memory bounds,
or device-independent causal constraints appear in APF as restrictions
on the allowed tested words, recovery maps, hidden records, and
no-added-resource simulations.  Thus APF does not identify every
temporal disturbance with IJC.  A temporal disturbance is an IJC
witness only when it rules out every faithful commuting continuation
defender satisfying the declared protocol constraints.  In the
nondisturbing static limit, this temporal formulation collapses back
to the usual global-section criterion; with genuine sequential words
included, the obstruction is a failure of global Boolean
representation for a process table rather than only for a
marginal-outcome table.

\subsection*{Accessible GPT fragments and cone equivalence}
The full bridge in Section~\ref{subsec:gpt-full-bridge} makes the accessible-fragment comparison explicit.  The static APF quotient induces an accessible ordered-vector-space fragment, Sep is simplex embeddability of that fragment, and cone-equivalent accessible fragments represent the same finite state--effect data up to positive order isomorphism.  The APF extension is that \(\mathcal T\) may include sequential words, recovery maps, and continuation equivalences, not just prepare--measure pairings; the ledger also distinguishes physically free identifications from resource-adding simulations.

\paragraph{GPT fragment.}
A finite GPT fragment consists of a finite family of preparations,
effects, transformations, and operational probabilities closed under
the tests under discussion.  APF maps such a fragment to a finite
tested continuation interface by treating preparations as tested input
states, effects as tested readout functionals, and transformations as
admissible continuation fragments.  The finite table of GPT
probabilities becomes the collection of state--effect values
$\omega_s(W)$ on tested continuation words $W\in\mathcal T$.

Under this translation, an APF commutative record defender is a
classical simplex model for the tested GPT fragment: states are
probability weights on a finite atom space $\Lambda$, effects are
functions on $\Lambda$, and tested transformations act through the
represented continuation calculus.  APF Sep therefore corresponds to
classical representability of the tested GPT fragment.  APF IJC
corresponds to failure of such a classical simplex representation once
the tested products, recovery maps, and operational equivalences are
imposed.

The APF contribution is not to replace GPT convex analysis.  It adds
the continuation quotient, the explicit finite tested word algebra,
and the ledger discipline specifying which distinctions, filters, and
composites are physically admissible at the interface.

\subsection*{Resource-theoretic hierarchy and univalent simulations}
The resource-theoretic hierarchy of GPT contextuality treats classical systems and univalent simulations as free structure.  APF's no-added-resource coarse-graining is the local analogue of a free operation: it may forget, identify, restrict, or lower resolution, but it may not smuggle in an uncharged memory, stabilizer, or ancilla.  Under that restriction, $\kappa_\Gamma$ behaves like a ledger monotone.  A full comparison with classical excess would require choosing a common norm and allowing approximate defenders, which is why this supplement states the exact finite theorem first and leaves quantitative monotones to the approximate layer.

\begin{proposition}[No-added-resource maps are ledger monotone]
Let finite interfaces and no-added-resource filters form the admissible morphisms of a local resource theory.  Then the pullback cost of any target distinction cannot be smaller than the target cost:
\[
  \kappa_{\Gamma'}(d')\le \kappa_\Gamma(F^*d'),
\]
whenever the pullback distinction is defined.  Hence \(\kappa\) is a monotone for the distinctions that survive the free map.  Additivity is not automatic; it still requires factorized support or ledger independence.
\end{proposition}
\begin{proof}
This is Theorem~\ref{thm:cost-monotonicity-no-added} restated in resource-theoretic language.  The proof is exactly that a free map can reuse, forget, merge, or restrict existing records, but cannot introduce an uncharged record that makes a distinction easier to enforce.
\end{proof}

\paragraph{Relation to resource-theoretic contextuality.}
Resource theories of contextuality often specify free embeddings or
simulations that cannot generate contextuality from classical
systems.  APF's no-added-resource condition plays an analogous local
role for finite interfaces: a filter may forget, coarse-grain,
relabel, or restrict already maintained records, but it may not
smuggle in an uncharged ancilla, stabilizer, memory, reference
frame, or adaptive controller.  The difference is one of scope.
Resource-theoretic contextuality usually orders whole operational
theories or GPT fragments.  APF works at the level of a fixed
finite tested interface $(\Gamma,Q,\mathcal T,\delta)$ and asks
whether that interface admits a faithful Boolean defender under the
declared resource triggers.  No-added-resource simulation is the
interface-level analogue of a free operation, while the ledger
records the costs and support assumptions under which monotonicity,
additivity, and independent counting are licensed.

\paragraph{Information-theoretic lower bounds for defender memory.}
When the defender must maintain a context-independent record $S$
capable of reproducing statistics across a tested context variable
$C$, standard information-theoretic quantities can provide lower
bounds on ledger cost.  If successful defense requires storing
context-relevant information with mutual information
\[
    I(S;C) \;>\; 0,
\]
then any physical implementation of the defender must provide
record support for at least that much external memory or
correlation.  Under a chosen calibration map $\chi_\Gamma$, such
information requirements may be translated into ledger lower
bounds, for example by charging memory, reset, stabilization, or
communication resources.  APF does not identify $\kappa_\Gamma$
universally with mutual information.  The relation is
calibration-dependent.  The point is that entropic monotones can
instantiate lower bounds on the abstract ledger in contexts where
the required defender memory, communication, or context-correlation
burden is explicitly modeled.

\subsection*{epBA minimal Boolean embeddings}
Exclusive partial Boolean algebra (epBA) constructions provide a logical route from partial event structures to a minimal classical counterpart.  For finite static event structures, APF's Boolean defender is the same kind of object: it is a global Boolean record completion of locally compatible tested events.  The difference is that APF builds the event structure from continuation identity and then asks for a faithful completion of tested sequential words.  Thus the epBA comparison is exact for finite static contexts and becomes a sequential Boolean-record completion problem when $\mathcal T$ contains products, recovery maps, or order-sensitive words.

\subsection*{External bookkeeping cost}
Information-theoretic obstruction models quantify the auxiliary context information required by a classical simulation that keeps a fixed shared state semantics.  This is directly aligned with the APF ledger reading: if a classical defender can reproduce a contextual table only by carrying an additional context label or memory, that label is not free.  In APF terms it must either be included in $\widetilde\Omega$ and charged as part of the defended interface, or the defender has failed to reproduce the tested continuation behavior without hidden resources.


\subsection*{NPA hierarchy and finite-level commuting certificates}
The NPA hierarchy supplies semidefinite outer approximations to quantum and commuting-operator correlations.  APF's finite-tested quotient has a different purpose: it classifies only the finite word behavior in $\mathcal T$ and asks whether it has a faithful Boolean-record completion.  Thus APF should not be read as a finite-level decision procedure for arbitrary commuting-operator completions.  Recent nonattainment and undecidability results for commuting-operator values reinforce the caution: finite certificates can be incomplete for global operator questions, whereas APF states an interface-relative finite theorem on purpose.

\begin{remark}[Practical finite-table certification]
For a finite Bell-type table, nonexistence of a Boolean defender is usually a finite linear-feasibility problem, not an NPA problem.  Enumerate deterministic Boolean atoms \(\lambda\) assigning outcomes to all queried settings and ask whether probabilities \(p_\lambda\ge0\), \(\sum_\lambda p_\lambda=1\), reproduce the observed tested marginals.  Failure of this linear program rules out an APF defender for that finite static table.  NPA-style SDPs become relevant when one asks the different question of whether the same table is compatible with quantum or general commuting-operator models.  APF's Sep/IJC split asks first whether a classical Boolean defender exists; Hilbert/Born admissibility is downstream.
\end{remark}

\subsection*{Ordered $*$-algebra representation results}
Ordered $*$-algebra representation theorems show that faithful operator or function representations require order-theoretic regularity conditions.  APF does not assume a $*$-algebra or a full positivity cone at this stage, but the moral is aligned: a commutative function representation becomes meaningful only after specifying the order/positivity data it must preserve.  That is why Definition~\ref{def:faithful-commutative-representation} explicitly includes tested unit and positivity preservation, while the stronger record-positivity assumptions needed for Hilbert/Born structure are deferred.

\subsection*{Operadic and multicategorical process frameworks}
Operadic and multicategorical accounts of quantum processes emphasize compositional structure, higher-order processes, and quotienting by operational equivalence or no-go ideals.  APF uses a simpler finite continuation-word algebra, but the comparison is natural: no-go ideals in process formalisms play a role analogous to APF contextual null spaces or excluded continuation patterns.  The APF-specific contribution is that the quotient is generated by finite admissible continuation identity and charged by an explicit enforcement ledger.

\subsection*{Causal order and process-theoretic reading}
A finite APF context can be organized as a partial category whose objects are admissible contexts and whose arrows are finite continuation fragments.  Sequential composition is categorical composition when defined.  Parallel composition may be added as a symmetric monoidal product only when the ledger supports factorized or explicitly interfaced composition; this is why additivity requires factorized support rather than being automatic.  A Sep defender is then a structure-preserving map from the tested partial process category into finite Boolean records: on static fragments it resembles a functor into finite sets/Boolean algebras, while on sequential fragments it must preserve the tested products and recovery maps.  Indefinite-causal-order or adaptive protocols can be included only by adding the relevant controlled words and control records to $\mathcal T$ and charging them in the ledger.

\begin{proposition}[Defenders as Boolean-record functors]
\label{prop:defenders-as-boolean-functors}
For a finite tested subcategory \(\mathsf{Cont}_{\Gamma,Q,\mathcal T}\), a Sep defender is equivalently a partial structure-preserving functorial representation
\[
  F_\rho:\mathsf{Cont}_{\Gamma,Q,\mathcal T}
  \longrightarrow \mathsf{FinBoolRec},
\]
where \(\mathsf{FinBoolRec}\) has finite Boolean record algebras as objects.  Its arrows are defined Boolean post-processing or recovery maps.  The representation preserves identities, tested composites, recovery/coarse-graining arrows, positivity/order of events, and the stated state--effect pairings.  Faithfulness says that \(F_\rho\) has kernel exactly the contextual continuation equivalence relation on the tested quotient.
\end{proposition}
\begin{proof}
A defender assigns one finite Boolean record algebra to the defended interface and assigns every tested continuation word a Boolean-record representative.  The defender axioms are precisely preservation of identity words, tested products, recovery/coarse-graining maps, and tested state--effect values.  Conversely, such a partial functorial representation into finite Boolean records supplies the atom space, record algebra, represented words, and recovery maps required by Definition~\ref{def:commuting-extension-defender}.  Equality of the functorial images exactly on context-equivalent words is the faithfulness condition.
\end{proof}

\subsection*{Technical bridge to existing frameworks}
In an ontological-model comparison, a Sep defender plays the role of a noncontextual hidden-variable completion, but with transformations included: the hidden record must reproduce finite words $E_{i_k}\cdots E_{i_1}$, not merely single-shot outcome marginals.  In a sheaf comparison, the Boolean defender is the global section object; IJC is the obstruction to extending the locally tested record assignments to a single global Boolean record.  In a sequential-instrument comparison, the generators $E_i$ behave like finite instruments or effects with update rules, while $\mathcal N_{\mathrm{cont}}$ identifies exactly those process words that no admissible continuation can separate.

This is why the Sep/IJC theorem is stronger than a static contextuality statement but weaker than a Hilbert-space reconstruction.  It certifies the failure of faithful commuting continuation bookkeeping.  It does not yet select complex Hilbert space, density matrices, or the Born trace rule.


\subsection*{Commutative operator algebras and Gelfand duality}
In the Sep branch, the finite reduced commutative algebra $\R^\Lambda$ is the real finite analogue of a commutative function algebra.  Thus the representation theorem is compatible with the familiar operator-algebraic slogan that classical record algebras are function algebras on a space of atoms.  APF does not claim novelty for that equivalence by itself.  The new ingredient is the route to the finite atom space through continuation identity, finite enforcement cost, and sequential word quotienting.

\subsection*{Process theories, categorical quantum mechanics, and effectus-style approaches}
The partial process-category structure of finite continuations is close in spirit to process-theoretic approaches: systems are compared by composable tests and transformations rather than by hidden underlying states.  APF differs by enforcing a quotient principle first: two process fragments are identical exactly when no admissible continuation can separate them in context.  The Sep/IJC theorem then asks whether the resulting finite process behavior has a faithful Boolean record completion.

\subsection*{Novelty boundary}
The theorem should not be read as rediscovering, in different notation, that finite Boolean models are commutative algebras.  That part is standard.  The APF contribution is the combination of (i) continuation identity as the definition of physical possibility, (ii) positive finite enforcement cost for distinctions, (iii) explicit capacity/ledger bookkeeping, and (iv) sequential continuation-word quotienting.  The representation theorem then says exactly when that APF-generated finite word behavior collapses to one shared Boolean record.
\section{Downstream Hilbert/Born preview}\label{sec:downstream-hilbert-born-preview}

The present supplement stops at the exact Sep/IJC branch theorem and the noncommutative continuation algebra forced by IJC.  The finite Hilbert/Born representation requires additional structure.

\begin{warningbox}
Hilbert space, complex amplitudes, tensor products, density matrices, and Born probabilities are not assumed and not fully derived in this supplement.
\end{warningbox}

The downstream quantum-structure theorem has the following conditional form:
\[
\begin{gathered}
\text{finite robust IJC continuation algebra}\;+ \\
\text{finite record positivity}\;+ \\
\text{operational duality}\;+ \\
\text{interface-complete composition} \\
\Longrightarrow \\
\bigoplus_k M_{n_k}(\C)\text{ representation} \\
\Longrightarrow \\
\omega(E)=\sum_k p_k\mathrm{Tr}(\rho_k E_k).
\end{gathered}
\]

The representation supplement must therefore supply:
\begin{enumerate}[leftmargin=2em]
\item positive normalized record calculus from finite records;
\item preparation/effect duality and the finite self-dual record regime;
\item finite-dimensional ordered $*$-algebra representation;
\item interface-complete composition and local/correlation tomography;
\item field selection from real/complex/quaternionic alternatives;
\item the finite matrix proof of the Born trace rule.
\end{enumerate}

\medskip
\noindent\textit{Correspondence with Paper~10's seven-gate list.}\quad
Paper~10 (\emph{The Calculus of Finite Continuability}) v1.11, \S5, names a seven-gate Hilbert--Born conditional representation theorem.  The seven gates describe the same conditional structure as the six items above, in a finer decomposition: Paper~10 gate~1 (finite tested word algebra modulo continuation equivalence) is the present supplement's continuation-word algebra of \S\ref{sec:continuation-word-algebra}; gates~2 (positivity and effect cones), 3 (reciprocal calibration), and 4 (adjoint or involution compatible with the tested order) decompose item~1 (positive normalized record calculus) and item~2 (preparation/effect duality) above; gate~5 (radical collapse / stable-sector completeness) and gate~6 (composite closure sufficient to select complex sectors) correspond to items~4 (interface-complete composition) and~5 (field selection); gate~7 (Gleason/Busch-style probability representation) corresponds to item~6 (finite matrix proof of the Born trace rule).  The two presentations are decompositions of the same conditional, not competing theorem statements; downstream representation supplements may use either.

\section{Assumption/theorem audit table}\label{sec:audit-table}

The table below records the role of each ingredient in the assumption economy.  The key point is that the supplement has a small primitive spine; most later machinery is either forced construction or a local regime condition for a theorem.

\begin{center}
\begin{longtable}{@{}p{0.27\linewidth}p{0.22\linewidth}p{0.40\linewidth}@{}}
\toprule
Ingredient & Role & Comment \\
\midrule
Finite admissible continuations $\Cont_\Gamma$ & Primitive spine & Defines what a possibility can still do. \\
Physical identity $x\sim y$ & Primitive spine / definition & Equality of finite continuation profiles. \\
Physical distinctions & Primitive spine & Finite enforceable separators of quotient classes. \\
Positive enforcement cost & Primitive spine & No zero-cost physical distinctions; not derived in this supplement. \\
Trigger-condition failure modes & Audit discipline & Each local theorem states the condition under which its conclusion is licensed. \\
Raw substrate $X_\Gamma$ & Prephysical bookkeeping & Candidate descriptions before quotienting. \\
Physical quotient $\Omega_\Gamma$ & Forced construction & Removes raw redundancy. \\
Physical descent & Forced construction & Distinctions and filters must factor through the quotient. \\
Ledger capacity $C_\Gamma$ & Regime parameter & Finite context budget for theorems involving capacity bounds. \\
Ledger admissibility & Forced bookkeeping & Records when finite distinction families fit inside a context budget. \\
Finite APF interrogation protocol & Normal form theorem & Supplies finite $Q$, finite $\mathcal T$, finite depth, finite resolution, and resolved floor. \\
Compact robust interface R1--R4 & Compact-extension hypothesis & Extends the positive-floor theorem beyond explicitly finite protocols. \\
Resolved finite-protocol floor & Normal form theorem & Finite resolved tested distinctions have a positive minimum cost. \\
Compact robust floor & Regime theorem & Compactness/LSC/robustness imply $\mu_\Gamma(Q)>0$ for compact families. \\
No-added-resource coarse-graining & Trigger condition/definition & Defines the pure coarse-graining class; excludes hidden amplifiers, ancillas, catalysts, or stabilizers. \\
Cost monotonicity & Regime theorem & Coarse-graining cannot increase required separation without added resources. \\
Classical control register & Regime construction & Available when finite classical alternatives are recordable. \\
Convexity & Forced theorem from control & Mixtures arise by coarse-graining over admissible control records. \\
Composition/factorization & Structural definitions & Say when two continuation spaces can be jointly supported. \\
Factorized support & Trigger condition/definition & Additivity is asserted only when enforcement support factorizes. \\
Interface cost & Forced bookkeeping & Tracks nonfactorized support without implying IJC. \\
Continuation words & Forced construction & Finite sequential process fragments generated by queries. \\
Contextual congruence & Forced construction & Words are identified when no admissible context separates them. \\
Continuation-word algebra & Forced construction & Quotient of formal word algebra by contextual continuation identity. \\
Formal linear span & Representation device & Not a physical mixture or superposition assumption. \\
Operationally resolved word classes & Regime definition & Counts only word classes distinguishable at fixed resolution. \\
Ledger-independent resolution & Trigger condition/definition & Counting theorem applies only to independently enforced resolved burdens. \\
Finite word resolution & Regime theorem & Bounds resolved independent classes by $C_\Gamma/\mu_\Gamma(Q)$. \\
Sep branch & Classification definition & A faithful commuting-extension defender exists. \\
IJC branch & Classification definition & No faithful commuting-extension defender exists. \\
Sep iff commutative representation & Main classification theorem & The finite word algebra has a faithful commutative record representation exactly in Sep. \\
IJC forces noncommutative faithful bookkeeping & Corollary & Any faithful algebraic realization of queried word composition has a noncommutative image algebra. \\
IJC certification for real systems & Empirical/structural witness layer & Bell, KS, complementarity, controlled order effects. \\
Approximate defenders & Deferred quantitative layer & Requires word depth, norm, noise model, and stability bounds. \\
Finite-margin robustness & Regime theorem sketch & Positive witness margin protects exact Sep/IJC classification under small finite-table perturbations. \\
Register-channel no-added-resource model & Operational model & Grounds cost monotonicity in finite records, channels, and absence of hidden stabilizers or ancillas. \\
Practical ledger-independence criteria & Operational criteria & Prevents double-counting shared records and clarifies borderline implementations. \\
Ledger-independence examples & Expository verification & Show shared coarse records, independent records, and coupled discriminators. \\
Worked Sep/IJC/coupled-classical examples & Expository verification & Demonstrates nonzero cost without IJC, controlled order obstruction, and coupled-but-Sep behavior. \\
Hilbert/Born & Deferred representation theory & Requires positive records, duality, composition, and tomography. \\
Paper~10 PLEC reduction & Cross-paper interpretive bridge & A1, MD, A2, BW are derivable consequences of the primitive spine under Paper~10 v1.11 \S3.5; not separate axioms here. \\
Marginal floor & Regime refinement & Coincides with this supplement's in-isolation cost floor when substrate has no shared structural commitments; the operative invariant in broken-symmetry / confined-phase regimes (Paper~10 v1.11 \S3). \\
Paper~10 seven-gate Hilbert/Born & Cross-decomposition & The seven gates of Paper~10 v1.11 \S5 describe the same conditional as \S\ref{sec:downstream-hilbert-born-preview}'s six-item list, in a finer decomposition. \\
\bottomrule
\end{longtable}
\end{center}

\section{Conclusion}\label{sec:conclusion}

\paragraph{Summary of the finite-interface result.}
The finite Sep/IJC result can now be read as a representation theorem
for tested continuation behavior.  Starting from finite admissible
continuation identity and positive-cost enforceable distinctions, a
finite interrogation protocol induces a resolved tested quotient
\[
    A^{\Gamma,\delta}_{Q,\mathcal T}.
\]
The Sep branch is exactly the case in which this quotient admits a
$\delta$-faithful representation inside one finite commutative Boolean
record algebra.  The IJC branch is exactly the failure of the
corresponding finite Boolean-record feasibility problem.  In static
nondisturbing fragments this reduces to the Fine--Boole--sheaf
global-section criterion; in sequential fragments it extends the
obstruction to tested continuation words, recovery maps, order effects,
and protocol-level resource constraints.  The ledger formalism
supplies the finite resource discipline under which distinctions,
filters, composites, and independent counts are licensed.

We have established the minimal Paper~1 spine without treating every technical hypothesis as a new physical input. The primitive claim is simple: physical states are finite admissible continuation profiles, and physical distinctions are finite positive-cost separators of those profiles. From that claim the quotient, descent requirements, ledger bookkeeping, controlled convexity, and continuation-word algebra are forced constructions.

A finite act of APF interrogation already supplies the finite query family, finite tested word set, finite depth, finite resolution, and positive resolved distinction floor used by the Sep/IJC theorem. What remains are local trigger conditions for specialized claims. Compact robust-interface hypotheses extend the floor theorem beyond explicit finite protocols. No-added-resource filters define the pure coarse-graining class for monotonicity. Factorized support licenses additive ledgers. Ledger-independent resolution licenses independent counting. None of these conditions enlarges the primitive spine; each states exactly where the stronger monotonicity, additivity, or counting conclusion applies. When a trigger condition fails, the corresponding theorem is not invoked: shared records are charged jointly rather than independently, resource-adding filters are not monotone, and infinite or noncompact families are not treated by the finite-interface theorem without further analytic completion.

The exact Sep/IJC representation theorem, for a fixed finite tested word set $\mathcal T$, then states:
\[
 (\Gamma,Q,\mathcal T)\in\Sep
 \quad\Longleftrightarrow\quad
 \AQtest\text{ admits a faithful commutative continuation representation}.
\]
IJC is the failure of all faithful commuting continuation completions. At such interfaces, any faithful algebraic bookkeeping of the queried finite continuation-word behavior must have a noncommutative image algebra. The quantum-structure paper takes this result as its starting point for the separate Hilbert/Born representation theorem.

\begin{thebibliography}{99}

\bibitem{Hardy2001}
L. Hardy, ``Quantum Theory From Five Reasonable Axioms,'' arXiv:quant-ph/0101012.

\bibitem{ChiribellaPurification2010}
G. Chiribella, G. M. D'Ariano, and P. Perinotti, ``Probabilistic theories with purification,'' \emph{Physical Review A} \textbf{81}, 062348 (2010).

\bibitem{Chiribella2011}
G. Chiribella, G. M. D'Ariano, and P. Perinotti, ``Informational derivation of quantum theory,'' \emph{Physical Review A} \textbf{84}, 012311 (2011).

\bibitem{Spekkens2005}
R. W. Spekkens, ``Contextuality for preparations, transformations, and unsharp measurements,'' \emph{Physical Review A} \textbf{71}, 052108 (2005).

\bibitem{AbramskyBrandenburger2011}
S. Abramsky and A. Brandenburger, ``The sheaf-theoretic structure of non-locality and contextuality,'' \emph{New Journal of Physics} \textbf{13}, 113036 (2011).

\bibitem{KochenSpecker1967}
S. Kochen and E. P. Specker, ``The problem of hidden variables in quantum mechanics,'' \emph{Journal of Mathematics and Mechanics} \textbf{17}, 59--87 (1967).

\bibitem{Fine1982}
A. Fine, ``Hidden variables, joint probability, and the Bell inequalities,'' \emph{Physical Review Letters} \textbf{48}, 291--295 (1982).

\bibitem{Bell1964}

J. S. Bell, ``On the Einstein Podolsky Rosen paradox,'' \emph{Physics} \textbf{1}, 195--200 (1964).

\bibitem{FoulisRandall1981}
D. J. Foulis and C. H. Randall, ``Empirical logic and tensor products,'' in \emph{Interpretations and Foundations of Quantum Theory}, Bibliographisches Institut, Mannheim (1981).

\bibitem{DaviesLewis1970}
E. B. Davies and J. T. Lewis, ``An operational approach to quantum probability,'' \emph{Communications in Mathematical Physics} \textbf{17}, 239--260 (1970).

\bibitem{Gudder1979}
S. Gudder, \emph{Stochastic Methods in Quantum Mechanics}, North-Holland (1979).


\bibitem{GelfandNaimark1943}
I. M. Gelfand and M. A. Naimark, ``On the imbedding of normed rings into the ring of operators in Hilbert space,'' \emph{Matematicheskii Sbornik} \textbf{12}, 197--213 (1943).

\bibitem{CoeckeKissinger2017}
B. Coecke and A. Kissinger, \emph{Picturing Quantum Processes}, Cambridge University Press (2017).

\bibitem{Jacobs2015}
B. Jacobs, ``New directions in categorical logic, for classical, probabilistic and quantum logic,'' \emph{Logical Methods in Computer Science} \textbf{11}, 1--76 (2015).

\bibitem{SelbySchmidWolfeSainzKunjwalSpekkens2023}
J. H. Selby, D. Schmid, E. Wolfe, A. B. Sainz, R. Kunjwal, and R. W. Spekkens, ``Accessible fragments of generalized probabilistic theories, cone equivalence, and applications to witnessing nonclassicality,'' \emph{Physical Review A} \textbf{107}, 062203 (2023); arXiv:2112.04521.

\bibitem{CataniGalleyGonda2026}
L. Catani, T. D. Galley, and T. Gonda, ``Resource-theoretic hierarchy of contextuality for general probabilistic theories,'' arXiv:2406.00717.

\bibitem{LiuWangWangHeWang2025}
S. Liu, Y. Wang, B. Wang, C. He, and J. Wang, ``The logical structure of contextuality and nonclassicality,'' arXiv:2506.15071.

\bibitem{Kim2026}
S.-J. Kim, ``Contextuality as an External Bookkeeping Cost under Fixed Shared-State Semantics,'' arXiv:2601.20167.

\bibitem{Barrett2007}
J. Barrett, ``Information processing in generalized probabilistic theories,'' \emph{Physical Review A} \textbf{75}, 032304 (2007).

\bibitem{NavascuesPironioAcin2008}
M. Navascu\'es, S. Pironio, and A. Ac\'in, ``A convergent hierarchy of semidefinite programs characterizing the set of quantum correlations,'' \emph{New Journal of Physics} \textbf{10}, 073013 (2008).

\bibitem{FanizzaKroellMehtaPaddockRochetteSlofstraZhao2025}
M. Fanizza, L. Kroell, A. Mehta, C. Paddock, D. Rochette, W. Slofstra, and Y. Zhao, ``The NPA hierarchy does not always attain the commuting operator value,'' arXiv:2510.04943 (2025).

\bibitem{Schoetz2023}
M. Sch\"otz, ``Gelfand--Naimark theorems for ordered $*$-algebras,'' \emph{Canadian Journal of Mathematics} \textbf{75}, 1272--1292 (2023); arXiv:1906.08752.

\bibitem{Chang2025}
S.-Y. Chang, ``Multicategorical Adjoints, Monadicity, and Quantum Resources,'' arXiv:2512.15951 (2025).


\bibitem{Landauer1961}
R. Landauer, ``Irreversibility and heat generation in the computing process,'' \emph{IBM Journal of Research and Development} \textbf{5}, 183--191 (1961).

\bibitem{Bennett1982}
C. H. Bennett, ``The thermodynamics of computation---a review,'' \emph{International Journal of Theoretical Physics} \textbf{21}, 905--940 (1982).
\end{thebibliography}

\end{document}
